\documentclass[main.tex]{subfiles}
\ifSubfilesClassLoaded{\BandSetupStandalone{B00}}{}
\title{Bd. 00 - Überblick über die Bände}
\author{Martin Kunze}
\date{}
\hypersetup{pdftitle={Bd. 00 - Überblick über die Bände},pdfauthor={Martin Kunze}}
\begin{document}
\hypersetup{pageanchor=false}
\maketitle
\hypersetup{pageanchor=true}
\setcounter{tocdepth}{0}
\tableofcontents
\setcounter{file}{0}
\renewcommand{\FormulaBandID}{0}
% Zweispaltige Übersichten nach der Vorlage aus Band 27.
\newcommand{\UebFormel}[1]{%
  \par\smallskip\noindent\(\displaystyle #1\)\par}
\newcommand{\UebQuellen}[1]{%
  \par\nobreak\smallskip{\footnotesize #1}\par}
\newenvironment{BandUebersicht}{%
  \begingroup
  \small
  \setlength{\LTleft}{0pt}%
  \setlength{\LTright}{0pt}%
  \setlength{\LTpre}{.5\baselineskip}%
  \setlength{\LTpost}{0pt}%
  \renewcommand{\arraystretch}{1.08}%
  \begin{longtable}{@{}>{\raggedright\arraybackslash}p{.455\linewidth}%
    @{\hspace{.012\linewidth}\vrule width .35pt\hspace{.012\linewidth}}%
    >{\raggedright\arraybackslash}p{.51\linewidth}@{}}%
  \toprule
  \textbf{Definitionen und Strukturaxiome}&\textbf{Zentrale Resultate}\\
  \midrule\endfirsthead
  \toprule
  \textbf{Definitionen und Strukturaxiome}&\textbf{Zentrale Resultate (Fortsetzung)}\\
  \midrule\endhead
  \midrule
  \multicolumn{2}{@{}r@{}}{\footnotesize\itshape Fortsetzung auf der nächsten Seite}\\
  \endfoot
  \bottomrule\endlastfoot
}{%
  \end{longtable}%
  \endgroup
}


\chapter*{Zur Benutzung}
\addcontentsline{toc}{chapter}{Zur Benutzung}
Dieser Band stellt die Definitionen und Strukturaxiome der Bände 1 bis 47
den daraus entwickelten zentralen Resultaten gegenüber. Links stehen die
Grundbegriffe und ihre Voraussetzungen, rechts die zugehörigen Folgerungen.
Die Quellenverweise führen zur vollständigen Formulierung und zum Beweis
im jeweiligen Fachband. Für eine Folgerung werden gegebenenfalls auch die
im Fachband angegebenen früheren Grundlagen benutzt.

Die Darstellung folgt der Übersicht \emph{Band 27 auf einen Blick}.
Ihr mathematischer Inhalt wurde hierher übernommen. Definitionen legen
Begriffe fest; Strukturaxiome beschreiben die jeweils untersuchten Objekte.
Schlussregeln und unbewiesene Voraussetzungen werden ausdrücklich als solche
bezeichnet. Eine Vermutung wird nicht als bewiesener Satz aufgeführt.

Links von \(\vdash\) stehen die Voraussetzungen, rechts die Folgerung;
\(\eqvdash\) bezeichnet die entsprechende Äquivalenz. Die in einer Zeile
oder in ihrer Lesart genannten Träger, Abbildungen und Typannahmen gelten
für beide Spalten. Die bestehenden Bandnummern bleiben unverändert.

\chapter[Band 1: Grundlagen der Logik]{Band 1 auf einen Blick}
\noindent\textbf{Grundlagen der Logik}\par\medskip
\noindent\emph{Lesart.} \(P,Q,R\) sind Formeln, \(t,u\) Terme und
\(\Gamma\) eine Menge offener Annahmen. Dieser Grundlagenband legt die
Sprache und den Kalkül fest. Rechts stehen daher ausdrücklich auch
\emph{Schlussregeln}; sie sind keine bereits bewiesenen semantischen Sätze.
\begin{BandUebersicht}
\textbf{Terme und atomare Formeln}
\UebFormel{\begin{aligned}
&x,\ c,\ f(t_1,\ldots,t_n)\quad\text{Terme},\\
&P(t_1,\ldots,t_n),\ t=u\quad\text{Formeln}.
\end{aligned}}
Funktions- und Prädikatensymbole haben jeweils eine feste Stelligkeit.
\UebQuellen{Band 1, Abschnitte \emph{Terme} und \emph{Aussagen der Logik}.}
&
\textit{Gleichheitsregeln}
\UebFormel{\begin{aligned}
&\vdash t=t,\\
&t=u\dsep P(t)\vdash P(u).
\end{aligned}}
Bei der Ersetzung werden gebundene Variablen nicht eingefangen.
\UebQuellen{Band 1, Regeln \(=I\) und \(=E\).}\\
\addlinespace[.8ex]
\textbf{Zusammengesetzte Formeln}
\UebFormel{\begin{aligned}
&\neg P,\quad P\land Q,\quad P\lor Q,\\
&P\rightarrow Q,\quad P\leftrightarrow Q.
\end{aligned}}
Diese Bildungsregeln bestimmen die zulässigen Junktorenformeln.
&
\textit{Einführung und Elimination der Konjunktion}
\UebFormel{\begin{aligned}
&P\dsep Q\vdash P\land Q,\\
&P\land Q\vdash P,\qquad P\land Q\vdash Q.
\end{aligned}}
\UebQuellen{Band 1, Regeln \(\land I\), \(\land E_1\), \(\land E_2\).}\\
\addlinespace[.8ex]
\textbf{Beweis und Ableitbarkeit}
Eine endliche Ableitung endet mit der behaupteten Formel; jede Zeile
ist Annahme, Axiom oder durch eine Regel gerechtfertigt.
\UebFormel{\Gamma\vdash_{\mathcal C}P.}
Ein Satz hat keine offenen Annahmen. Inhaltliche Axiome werden später
als Voraussetzungen einer Theorie eingeführt.
&
\textit{Implikation und Entlassung einer Annahme}
\UebFormel{\begin{aligned}
&[P]\vdots Q\quad\vdash\quad P\rightarrow Q,\\
&P\rightarrow Q\dsep P\vdash Q.
\end{aligned}}
Im ersten Schema wird die ausgezeichnete Annahme \(P\) entlassen;
andere offene Annahmen bleiben bestehen.
\UebQuellen{Band 1, Regeln \(\rightarrow I\) und \(\rightarrow E\).}\\
\addlinespace[.8ex]
\textbf{Fallunterscheidung und Widerspruch}
Die Metanotation \([P]\vdots R\) bezeichnet einen Hilfsbeweis mit
ausgezeichneter Annahme \(P\). \(\bot\) bezeichnet Falschheit.
&
\textit{Disjunktions- und Negationsregeln}
\UebFormel{\begin{aligned}
&P\vdash P\lor Q,\qquad Q\vdash P\lor Q,\\
&P\lor Q\dsep[P]\vdots R\dsep[Q]\vdots R\\
&\hspace{18mm}\vdash R,\\
&P\dsep\neg P\vdash\bot,\\
&[P]\vdots\bot\quad\vdash\quad\neg P.
\end{aligned}}
\UebQuellen{Band 1, Regeln \(\lor I\), \(\lor E\), \(\bot I\), \(\neg I\).}\\
\addlinespace[.8ex]
\textbf{Quantifizierte Formeln}
\UebFormel{\forall x\,P(x),\qquad\exists x\,P(x).}
Ein Term muss für die ersetzte Variable einsetzbar sein. Bei den
Eigenvariablenregeln gelten die im Band angegebenen Frischebedingungen.
&
\textit{Quantorenregeln}
\UebFormel{\begin{aligned}
&\forall x\,P(x)\vdash P(t),\\
&P(t)\vdash\exists x\,P(x),\\
&P(x)\mid_x\quad\vdash\quad\forall x\,P(x),\\
&\exists x\,P(x)\dsep[P(x)]\vdots Q\mid_x\\
&\hspace{18mm}\vdash Q.
\end{aligned}}
Bei \(\forall I\) darf \(x\) in keiner offenen Annahme frei sein;
bei \(\exists E\) auch nicht in der Schlussformel \(Q\).
\UebQuellen{Band 1, Regeln \(\forall I\), \(\forall E\), \(\exists I\), \(\exists E\).}\\
\addlinespace[.8ex]
\textbf{Definitionen und eindeutige Beschreibung}
Explizite Definitionen führen Abkürzungen ein; implizite Definitionen
legen eine Struktur durch Axiome fest. Ein Iota-Term setzt eindeutige
Existenz voraus.
\UebFormel{a\coloneqq\iota x\,P(x).}
&
\textit{Verwendung einer eindeutigen Beschreibung}
Unter der Voraussetzung \(\exists!x\,P(x)\) kann der eindeutig
beschriebene Gegenstand eingesetzt werden:
\UebFormel{\begin{aligned}
&P(\iota x\,P(x)),\\
&P(y)\vdash y=\iota x\,P(x).
\end{aligned}}
\UebQuellen{Band 1, Abschnitt \emph{Das Iota-Symbol und Iota-Definitionen}.}\\
\addlinespace[.8ex]
\textbf{Metanotation}
\UebFormel{\begin{aligned}
&a\neq b\quad\text{steht für}\quad\neg(a=b),\\
&P\eqvdash Q\quad\text{steht für}\\
&\qquad P\vdash Q\ \text{und}\ Q\vdash P.
\end{aligned}}
Metasymbole strukturieren die Darstellung von Ableitungen.
&
\textit{Ersetzung äquivalenter Teilausdrücke}
Bewiesene Äquivalenzen erlauben die im Band formulierte
metasprachliche Ersetzung in einem Formelschema. Dabei werden
die Voraussetzungen der verwendeten Äquivalenz mitgeführt.
\UebQuellen{Band 1, Abschnitt \emph{Meta-Regel: Ersetzung äquivalenter Teilausdrücke}.}\\
\end{BandUebersicht}

\chapter[Band 2: Theoreme der Logik]{Band 2 auf einen Blick}
\noindent\textbf{Theoreme der Logik}\par\medskip
\noindent\emph{Lesart.} \(P,Q,R\) sind beliebige Formeln;
\(P(x)\) und \(Q(x)\) dürfen freie Variablen besitzen. Die
Grundjunktoren und Quantoren stammen aus Band 1; dieser Band beweist
ihre Rechengesetze und definiert weitere logische Operationen.
\begin{BandUebersicht}
\textbf{Konjunktion und Disjunktion}
Die iterierte Schreibweise wird rechts geklammert:
\UebFormel{\begin{aligned}
P\land Q\land R&\coloneqq P\land(Q\land R),\\
P\lor Q\lor R&\coloneqq P\lor(Q\lor R).
\end{aligned}}
\UebQuellen{\FormulaRefAuto{P \land Q \land R \coloneqq P \land (Q \land R)}[def];
\FormulaRefAuto{P \lor Q \lor R \coloneqq P \lor (Q \lor R)}[def]}
&
\textit{Idempotenz, Kommutativität und Distribution}
\UebFormel{\begin{aligned}
&P\land P\eqvdash P,\quad P\lor P\eqvdash P,\\
&P\land Q\eqvdash Q\land P,\\
&P\land(Q\lor R)\\
&\hspace{8mm}\eqvdash(P\land Q)\lor(P\land R).
\end{aligned}}
\UebQuellen{\FormulaRefAuto{P \land P \eqvdash P}[thm];
\FormulaRefAuto{P \lor P \eqvdash P}[thm];
\FormulaRefAuto{P\land Q\eqvdash Q\land P}[thm];
\FormulaRefAuto{P \land (Q\lor R) \eqvdash (P \land Q) \lor (P\land R)}[thm]}\\
\addlinespace[.8ex]
\textbf{Negation und Wahrheitswerte}
\UebFormel{\top\coloneqq\neg\bot.}
\(\bot\) und \(\top\) erlauben die Formulierung der Wahrheitstafeln
als beweisbare Äquivalenzen.
\UebQuellen{\FormulaRefAuto{\top \coloneqq \neg\bot}[def]}
&
\textit{Klassische Negationsgesetze}
\UebFormel{\begin{aligned}
&P\eqvdash\neg\neg P,\qquad P\lor\neg P,\\
&\neg(P\lor Q)\eqvdash\neg P\land\neg Q,\\
&\neg(P\land Q)\eqvdash\neg P\lor\neg Q,\\
&\bot\vdash P.
\end{aligned}}
\UebQuellen{\FormulaRefAuto{P \eqvdash \neg\neg P}[thm];
\FormulaRefAuto{P \lor \neg P}[thm];
\FormulaRefAuto{\neg(P \lor Q) \eqvdash \neg P \land \neg Q}[thm];
\FormulaRefAuto{\neg(P \land Q) \eqvdash \neg P \lor \neg Q}[thm];
\FormulaRefAuto{\bot\vdash P}[thm]}\\
\addlinespace[.8ex]
\textbf{Implikation und Bikonditional}
Die Grundzeichen \(\rightarrow\) und \(\leftrightarrow\) werden
mittels ihrer Einführungs- und Eliminationsregeln verwendet.
&
\textit{Charakterisierungen und Kontraposition}
\UebFormel{\begin{aligned}
&P\rightarrow Q\eqvdash\neg P\lor Q,\\
&P\rightarrow Q\eqvdash\neg Q\rightarrow\neg P,\\
&P\leftrightarrow Q\\
&\hspace{8mm}\eqvdash(P\rightarrow Q)\land(Q\rightarrow P),\\
&(P\land Q)\rightarrow R\\
&\hspace{8mm}\eqvdash P\rightarrow(Q\rightarrow R).
\end{aligned}}
\UebQuellen{\FormulaRefAuto{P \rightarrow Q \eqvdash \neg P \lor Q}[thm];
\FormulaRefAuto{P \rightarrow Q \eqvdash \neg Q \rightarrow \neg P}[thm];
\FormulaRefAuto{P \leftrightarrow Q \eqvdash (P \rightarrow Q)\land (Q\rightarrow P)}[thm];
\FormulaRefAuto{(P \land Q) \rightarrow R \eqvdash P \rightarrow (Q \rightarrow R)}[thm]}\\
\addlinespace[.8ex]
\textbf{Exklusives Oder}
\UebFormel{P\lxor Q\coloneqq\neg(P\leftrightarrow Q).}
\UebQuellen{\FormulaRefAuto{P \lxor Q \coloneqq \neg(P \leftrightarrow Q)}[def]}
&
\textit{Genau einer der beiden Fälle}
\UebFormel{\begin{aligned}
&P\lxor Q\eqvdash(\neg P\land Q)\lor(P\land\neg Q),\\
&P\lxor Q\eqvdash Q\lxor P,\qquad\neg(P\lxor P).
\end{aligned}}
\UebQuellen{\FormulaRefAuto{P \lxor Q \eqvdash (\neg P\land Q)\lor(P\land \neg Q)}[thm];
\FormulaRefAuto{P\lxor Q \eqvdash Q\lxor P}[thm];
\FormulaRefAuto{\neg(P\lxor P)}[thm]}\\
\addlinespace[.8ex]
\textbf{NAND und NOR}
\UebFormel{\begin{aligned}
P\lnand Q&\coloneqq\neg(P\land Q),\\
P\lnor Q&\coloneqq\neg(P\lor Q).
\end{aligned}}
\UebQuellen{\FormulaRefAuto{P \lnand Q \coloneqq \neg(P \land Q)}[def];
\FormulaRefAuto{P \lnor Q \coloneqq \neg(P \lor Q)}[def]}
&
\textit{Darstellung anderer Junktoren}
\UebFormel{\begin{aligned}
&P\lnand P\eqvdash\neg P,\\
&P\land Q\eqvdash(P\lnand Q)\lnand(P\lnand Q),\\
&P\lnor P\eqvdash\neg P,\\
&P\lor Q\eqvdash(P\lnor Q)\lnor(P\lnor Q).
\end{aligned}}
\UebQuellen{\FormulaRefAuto{P\lnand P \eqvdash \neg P}[thm];
\FormulaRefAuto{P\land Q \eqvdash (P\lnand Q)\lnand(P\lnand Q)}[thm];
\FormulaRefAuto{P\lnor P \eqvdash \neg P}[thm];
\FormulaRefAuto{P\lor Q \eqvdash (P\lnor Q)\lnor(P\lnor Q)}[thm]}\\
\addlinespace[.8ex]
\textbf{Halbaddierer}
\UebFormel{\begin{aligned}
\HAddSum{P}{Q}&\coloneqq P\lxor Q,\\
\HAddCarry{P}{Q}&\coloneqq P\land Q.
\end{aligned}}
\UebQuellen{\FormulaRefAuto{\HAddSum{P}{Q}\coloneqq P\lxor Q}[def];
\FormulaRefAuto{\HAddCarry{P}{Q}\coloneqq P\land Q}[def]}
&
\textit{Summe und Übertrag schließen einander aus}
\UebFormel{\begin{aligned}
&\HAddSum{P}{Q}\vdash\neg\HAddCarry{P}{Q},\\
&P\dsep Q\vdash
\neg\HAddSum{P}{Q}\land\HAddCarry{P}{Q}.
\end{aligned}}
\UebQuellen{\FormulaRefAuto{\HAddSum{P}{Q}\vdash \neg\HAddCarry{P}{Q}}[thm];
\FormulaRefAuto{P\dsep Q\vdash \neg\HAddSum{P}{Q}\land \HAddCarry{P}{Q}}[thm]}\\
\addlinespace[.8ex]
\textbf{Volladdierer}
\UebFormel{\begin{aligned}
&\FAddSum{P}{Q}{R}\coloneqq\\
&\qquad\HAddSum{\HAddSum{P}{Q}}{R},\\
&\FAddCarry{P}{Q}{R}\coloneqq\HAddCarry{P}{Q}\\
&\qquad\lor\HAddCarry{\HAddSum{P}{Q}}{R}.
\end{aligned}}
\UebQuellen{\FormulaRefAuto{\FAddSum{P}{Q}{R}\coloneqq \HAddSum{\HAddSum{P}{Q}}{R}}[def];
\FormulaRefAuto{\FAddCarry{P}{Q}{R}\coloneqq \HAddCarry{P}{Q}\lor \HAddCarry{\HAddSum{P}{Q}}{R}}[def]}
&
\textit{Explizite Formeln}
\UebFormel{\begin{aligned}
&\FAddSum{P}{Q}{R}\eqvdash(P\lxor Q)\lxor R,\\
&\FAddCarry{P}{Q}{R}\\
&\hspace{6mm}\eqvdash(P\land Q)\lor((P\lxor Q)\land R).
\end{aligned}}
\UebQuellen{\FormulaRefAuto{\FAddSum{P}{Q}{R}\eqvdash (P\lxor Q)\lxor R}[thm];
\FormulaRefAuto{\FAddCarry{P}{Q}{R}\eqvdash (P\land Q)\lor((P\lxor Q)\land R)}[thm]}\\
\addlinespace[.8ex]
\textbf{Identität und Ungleichheit}
Die Gleichheit besitzt die Regeln aus Band 1; \(a\neq b\)
ist die Negation von \(a=b\).
&
\textit{Kongruenz von Funktionstermen}
\UebFormel{\begin{aligned}
&a=b\vdash f(c,a)=f(c,b),\\
&a=b\vdash f(a,c)=f(b,c).
\end{aligned}}
\UebQuellen{\FormulaRefAuto{BinaryFunctionRightArgumentCongruence}[thm];
\FormulaRefAuto{BinaryFunctionLeftArgumentCongruence}[thm]}\\
\addlinespace[.8ex]
\textbf{Quantoren und Eindeutigkeit}
\UebFormel{\begin{aligned}
&\existsleone x\,P(x)\coloneqq\\
&\quad\forall y\forall z\,((P(y)\land P(z))\rightarrow y=z).
\end{aligned}}
\UebQuellen{\FormulaRefAuto{\existsleone x\,P(x) \coloneqq \forall y\forall z((P(y)\land P(z))\rightarrow y=z)}[def]}
&
\textit{Quantorennegation und eindeutige Existenz}
\UebFormel{\begin{aligned}
&\neg\forall x\,P(x)\eqvdash\exists x\,\neg P(x),\\
&\forall x\,\neg P(x)\eqvdash\neg\exists x\,P(x),\\
&\exists!x\,P(x)\eqvdash
\exists x\,P(x)\land\existsleone x\,P(x).
\end{aligned}}
\UebQuellen{\FormulaRefAuto{\neg\forall x(P(x)) \eqvdash \exists x (\neg P(x))}[thm];
\FormulaRefAuto{\forall x(\neg P(x)) \eqvdash \neg\exists x (P(x))}[thm];
\FormulaRefAuto{\exists! x\,P(x) \eqvdash \exists x\,P(x)\land \existsleone x\,P(x)}[thm]}\\
\addlinespace[.8ex]
\textbf{Konditionaler Aussagenoperator und Fallterm}
\UebFormel{\begin{aligned}
\IfWff{P}{Q}{R}&\coloneqq(P\land Q)\lor(\neg P\land R),\\
&\IfThenElse{P}{a}{b}.
\end{aligned}}
Der Fallterm bezeichnet \(a\) im Fall \(P\), andernfalls \(b\).
\UebQuellen{\FormulaRefAuto{if-}[def];\FormulaRefAuto{Fallterm}[def]}
&
\textit{Implikationsform und Auswertung}
\UebFormel{\begin{aligned}
&\IfWff{P}{Q}{R}\\
&\hspace{6mm}\eqvdash(P\rightarrow Q)\land(\neg P\rightarrow R),\\
&P\vdash\IfThenElse{P}{a}{b}=a,\\
&\neg P\vdash\IfThenElse{P}{a}{b}=b.
\end{aligned}}
\UebQuellen{\FormulaRefAuto{if-Implikationsform}[thm];
\FormulaRefAuto{P \vdash \IfThenElse{P}{a}{b}=a}[thm];
\FormulaRefAuto{\neg P \vdash \IfThenElse{P}{a}{b}=b}[thm]}\\
\end{BandUebersicht}

\chapter[Band 3: Mengenlehre]{Band 3 auf einen Blick}
\noindent\textbf{Mengenlehre}\par\medskip
\noindent\emph{Lesart.} Alle Gegenstände sind Mengen. \(P\) und
\(\varphi\) bezeichnen Formeln der Mengensprache; bei den Axiomenschemata
gelten die Frischebedingungen des Originals. Die Relations- und
Operatorbegriffe gehören ausdrücklich zu diesem Band.
\begin{BandUebersicht}
\textbf{ZFC und Extensionalität}
Die Mengenstruktur wird durch die ZFC-Axiome beschrieben.
\UebFormel{\begin{aligned}
&\forall x\,(x\in A\leftrightarrow x\in B)\\
&\hspace{15mm}\vdash A=B.
\end{aligned}}
\UebQuellen{\FormulaRefAuto{Menge}[def];
\FormulaRefAuto{\forall x\, (x \in A \leftrightarrow x \in B) \vdash A = B}[ax]}
&
\textit{Gleichheit wird durch Elemente bestimmt}
\UebFormel{\begin{aligned}
&A=B\eqvdash\forall x\,(x\in A\leftrightarrow x\in B),\\
&x\in A\dsep y\notin A\vdash x\neq y.
\end{aligned}}
\UebQuellen{\FormulaRefAuto{\forall x\, (x \in A \leftrightarrow x \in B) \eqvdash A = B}[thm];
\FormulaRefAuto{x\in A\dsep y\notin A \vdash x\neq y}[thm]}\\
\addlinespace[.8ex]
\textbf{Teilmengen und leere Menge}
\UebFormel{\begin{aligned}
A\subseteq B&\coloneq\forall x\,(x\in A\rightarrow x\in B),\\
A\subset B&\coloneq A\subseteq B\land A\neq B,\\
\varnothing&\coloneq\iota O\,\forall x\,(x\notin O).
\end{aligned}}
\UebQuellen{\FormulaRefAuto{A \subseteq B \coloneq \forall x\,(x\in A \rightarrow x\in B)}[def];
\FormulaRefAuto{A \subset B \coloneq A \subseteq B \land A\neq B}[def];
\FormulaRefAuto{\varnothing \coloneq \iota O\bigl(\forall x\,(x \not\in O)\bigr)}[def]}
&
\textit{Leere Menge und echte Erweiterungen}
\UebFormel{\begin{aligned}
&\varnothing\subseteq A,\qquad
A\subseteq\varnothing\vdash A=\varnothing,\\
&A\neq\varnothing\eqvdash\exists x\,(x\in A),\\
&A\subset B\vdash\exists x\in B\setminus A.
\end{aligned}}
\UebQuellen{\FormulaRefAuto{\varnothing\subseteq A}[thm];
\FormulaRefAuto{A\subseteq\varnothing\vdash A=\varnothing}[thm];
\FormulaRefAuto{A\neq\varnothing \eqvdash \exists x\,(x\in A)}[thm];
\FormulaRefAuto{A\subset B\vdash \exists x\in B\setminus A}[thm]}\\
\addlinespace[.8ex]
\textbf{Aussonderung}
Das Aussonderungsaxiom garantiert die Menge der Elemente von \(A\),
die \(P\) erfüllen:
\UebFormel{\{x\in A\mid P(x)\}.}
\UebQuellen{\FormulaRefAuto{\exists B\;\forall x\;\bigl(x \in B \;\leftrightarrow\; x \in A \,\land\, P(x)\bigr)}[ax];
\FormulaRefAuto{\{x\in A \,\mid\, P(x)\} \coloneq \iota B\bigl(\forall x\,(x\in B \leftrightarrow (x\in A \land P(x)))\bigr)}[def]}
&
\textit{Elementkriterium und Grenze der Komprehension}
\UebFormel{\begin{aligned}
&x\in\{u\in A\mid P(u)\}\\
&\hspace{8mm}\eqvdash x\in A\land P(x),\\
&\exists x\,(x\notin A).
\end{aligned}}
Insbesondere gibt es keine Menge aller Mengen.
\UebQuellen{\FormulaRefAuto{x \in \{u \in A \mid P(u)\} \eqvdash x \in A \land P(x)}[thm];
\FormulaRefAuto{FreshElementOutsideSet}[thm]}\\
\addlinespace[.8ex]
\textbf{Schnitt und Differenz}
\UebFormel{\begin{aligned}
A\cap B&\coloneq\{x\in A\mid x\in B\},\\
A\setminus B&\coloneq\{x\in A\mid x\notin B\}.
\end{aligned}}
\UebQuellen{\FormulaRefAuto{A \cap B \coloneq \{ x \in A \mid x \in B \}}[def];
\FormulaRefAuto{A \setminus B \coloneq \{ x \in A \mid x \notin B \}}[def]}
&
\textit{Relative Komplemente}
Für \(X,Y\subseteq M\):
\UebFormel{\begin{aligned}
&M\setminus(M\setminus X)=X,\\
&X\subseteq Y\eqvdash M\setminus Y\subseteq M\setminus X.
\end{aligned}}
\UebQuellen{\FormulaRefAuto{RelCompInvolutive}[thm];
\FormulaRefAuto{RelCompAntitone}[thm]}\\
\addlinespace[.8ex]
\textbf{Paare und geordnete Paare}
\UebFormel{\begin{aligned}
\{a\}&\coloneq\{a,a\},\\
(a,b)&\coloneq\{\{a\},\{a,b\}\}.
\end{aligned}}
Das Paarmengenaxiom liefert \(\{a,b\}\).
\UebQuellen{\FormulaRefAuto{\{a\} \coloneq \{a,a\}}[def];
\FormulaRefAuto{(a, b) \coloneq \{ \{ a \}, \{ a, b \} \}}[def]}
&
\textit{Elemente einer Paarmenge und feste Paar-Tags}
\UebFormel{\begin{aligned}
&x\in\{A,B\}\eqvdash x=A\lor x=B,\\
&(p,x)=(p,y)\vdash x=y,\\
&p\neq q\vdash(p,x)\neq(q,y).
\end{aligned}}
\UebQuellen{\FormulaRefAuto{x \in \{A,B\}\;\eqvdash\;(x = A \lor x = B)}[thm];
\FormulaRefAuto{TaggedPairRightInjective}[thm];
\FormulaRefAuto{TaggedPointsDifferent}[thm]}\\
\addlinespace[.8ex]
\textbf{Vereinigung und Familiendurchschnitt}
\UebFormel{A\cup B\coloneq\bigcup\{A,B\}.}
\(\bigcup M\) sammelt die Elemente der Familienglieder;
\(\bigcap M\) wird für \(M\neq\varnothing\) definiert.
\UebQuellen{\FormulaRefAuto{A \cup B \coloneq \bigcup \{A, B\}}[def];
\FormulaRefAuto{M\neq\varnothing \vdash \bigcap M \coloneq \bigcap\{X\mid X\in M\}}[def]}
&
\textit{Distributivität und gemeinsamer Inhalt}
\UebFormel{\begin{aligned}
&A\cap(B\cup C)=(A\cap B)\cup(A\cap C),\\
&M\neq\varnothing\ \vdash\\
&\quad x\in\bigcap M\leftrightarrow\forall X\in M\,(x\in X).
\end{aligned}}
\UebQuellen{\FormulaRefAuto{A \cap (B \cup C) = (A \cap B) \cup (A \cap C)}[thm];
\FormulaRefAuto{M\neq\varnothing \vdash x\in\bigcap M \leftrightarrow \forall X\in M\,(x\in X)}[thm]}\\
\addlinespace[.8ex]
\textbf{Mengenfamilien}
\UebFormel{\begin{aligned}
&\CoverFam(M,A)\coloneqq A\subseteq\bigcup M,\\
&\PartFam(M,A)\coloneqq\\
&\qquad\DisjFam(M)\land A=\bigcup M.
\end{aligned}}
\(\DisjFam(M)\): nichtleere, paarweise disjunkte Glieder.
\(\UnionClosedFam(M)\): mit \(X,Y\in M\) liegt auch \(X\cup Y\) in \(M\).
Das Disjunktheitsaxiom lautet
\UebFormel{X,Y\in M\dsep X\neq Y\vdash X\cap Y=\varnothing.}
\UebQuellen{\FormulaRefAuto{\CoverFam(M,A) \coloneqq A \subseteq \bigcup M}[def];
\FormulaRefAuto{\PartFam(M,A) \coloneqq \DisjFam(M)\land A=\bigcup M}[def];
\FormulaRefAuto{Disjunkte Familie}[def];\FormulaRefAuto{Vereinigungsabgeschlossene Familie}[def]}
&
\textit{Eindeutigkeit des Familiengliedes}
Für \(\DisjFam(M)\):
\UebFormel{\begin{aligned}
&X,Y\in M\dsep a\in X\dsep a\in Y\vdash X=Y,\\
&X,Y\in M\dsep X\neq Y\dsep a\in X\\
&\hspace{10mm}\vdash a\notin Y.
\end{aligned}}
\UebQuellen{\FormulaRefAuto{X \in M \dsep Y \in M \dsep a \in X \dsep a \in Y \vdash X = Y}[thm];
\FormulaRefAuto{X \in M \dsep Y \in M \dsep X \neq Y \dsep a \in X \vdash a \notin Y}[thm]}
\\\addlinespace[.8ex]
\textbf{Regularität}
\UebFormel{\begin{aligned}
&A\neq\varnothing\\
&\hspace{8mm}\vdash\exists x\in A\,(x\cap A=\varnothing).
\end{aligned}}
\UebQuellen{\FormulaRefAuto{A \neq \varnothing \vdash \exists x \in A \,(x \cap A = \varnothing)}[ax]}
&
\textit{Keine Selbstmitgliedschaft}
\UebFormel{\begin{aligned}
&a\notin a,\qquad a\in b\vdash b\notin a,\\
&a\cup\{a\}=b\cup\{b\}\eqvdash a=b.
\end{aligned}}
\UebQuellen{\FormulaRefAuto{a\not\in a}[thm];
\FormulaRefAuto{a\in b\vdash b\not\in a}[thm];
\FormulaRefAuto{a\cup\{a\}=b\cup\{b\}\eqvdash a=b}[thm]}\\
\addlinespace[.8ex]
\textbf{Potenzmenge und Mengensystem}
\(\powerset(A)\) ist die durch das Potenzmengenaxiom bestimmte Menge
aller Teilmengen von \(A\).
\UebFormel{\begin{aligned}
&\mathsf{SetSys}(M;X)\coloneqq M\subseteq\powerset(X),\\
&\mathfrak I[P,Q]\coloneqq\\
&\qquad\{H\in\powerset(Q)\mid P\subseteq H\}.
\end{aligned}}
\UebQuellen{\FormulaRefAuto{SetSystemOverDef}[def];
\FormulaRefAuto{BooleanSetIntervalDef}[def]}
&
\textit{Elementkriterium und Monotonie}
\UebFormel{\begin{aligned}
&x\in\powerset(A)\eqvdash x\subseteq A,\\
&A\subseteq B\vdash\powerset(A)\subseteq\powerset(B),\\
&\powerset(\varnothing)=\{\varnothing\}.
\end{aligned}}
\UebQuellen{\FormulaRefAuto{x \in \powerset(A)\;\eqvdash\; x \subseteq A}[thm];
\FormulaRefAuto{A \subseteq B \vdash \powerset(A) \subseteq \powerset(B)}[thm];
\FormulaRefAuto{\powerset(\varnothing)=\{\varnothing\}}[thm]}\\
\addlinespace[.8ex]
\textbf{Kartesisches Produkt, Relation und Graph}
\(A\times B\) enthält die geordneten Paare aus \(A\) und \(B\).
\UebFormel{\begin{aligned}
&F\subseteq A\times B\quad\text{Relation},\\
&\GraphStruct{V,E}\coloneqq E\subseteq V\times V.
\end{aligned}}
\UebQuellen{\FormulaRefAuto{F \subseteq A \times B}[def];
\FormulaRefAuto{GraphDef}[def]}
&
\textit{Paar- und Produktkriterien}
\UebFormel{\begin{aligned}
&(a,b)\in A\times B\eqvdash a\in A\land b\in B,\\
&A\times\varnothing=\varnothing,\\
&\varnothing\times A=\varnothing.
\end{aligned}}
Eine Relation muss hier weder total noch funktional sein; diese
Axiome folgen erst in den eigenständigen Bänden 4 und 5.
\UebQuellen{\FormulaRefAuto{(a,b)\in A\times B\eqvdash a\in A\land b\in B}[thm];
\FormulaRefAuto{A\times\varnothing=\varnothing}[thm];
\FormulaRefAuto{\varnothing\times A=\varnothing}[thm]}\\
\addlinespace[.8ex]
\textbf{Binärer Operator}
\UebFormel{\mathsf{BinOp}(A,\star).}
Das Strukturaxiom fordert Abgeschlossenheit:
\UebFormel{x,y\in A\vdash x\star y\in A.}
\UebQuellen{\FormulaRefAuto{BinaryOperationDef}[def];
\FormulaRefAuto{BinaryOperationClosureAxiom}[ax]}
&
\textit{Reichweite dieses Grundbegriffs}
Der Band stellt damit die Typisierung algebraischer Terme bereit.
Assoziativität, Kommutativität, Idempotenz und Kürzbarkeit werden
hier noch nicht vorausgesetzt. Die entsprechenden eigenständigen
Strukturbände ergänzen jeweils ihre eigenen Axiome.\\
\addlinespace[.8ex]
\textbf{Ersetzung, Auswahl und Unendlichkeit}
Ersetzung sammelt eindeutig durch eine Formel bestimmte Werte.
Das Auswahlaxiom trifft jedes Glied einer disjunkten nichtleeren
Familie genau einmal. Unendlichkeit fordert eine induktive Menge:
\UebFormel{\exists A\,(\varnothing\in A\land
\forall x\in A\,(x\cup\{x\}\in A)).}
\UebQuellen{\FormulaRefAuto{ReplacementSchema}[ax];
\FormulaRefAuto{ChoiceFamilyAxiom}[ax];
\FormulaRefAuto{\exists A\;\bigl(\varnothing \in A \land \forall x \in A\,(x \cup\{x\} \in A)\bigr)}[ax]}
&
\textit{Durch einen Mengenterm beschriebene Bildmenge}
Für einen Mengenterm \(t(j)\) mit festen Parametern:
\UebFormel{\begin{aligned}
&\exists!X\,\forall x\,\bigl(\\
&\quad x\in X\leftrightarrow
\exists j\,(j\in J\land x=t(j))\bigr).
\end{aligned}}
Dies folgt aus Ersetzung und Extensionalität und setzt noch
keinen Funktionsgraphen voraus.
\UebQuellen{\FormulaRefAuto{TermImageSetUniqueExists}[thm]}\\
\end{BandUebersicht}

\chapter[Band 4: Totale Relationen]{Band 4 auf einen Blick}
\noindent\textbf{Totale Relationen}\par\medskip
\noindent\emph{Lesart.} \(A,B,M,R\) sind Mengen. Totalität verlangt
mindestens einen zugeordneten Wert, ohne Eindeutigkeit zu fordern.
\begin{BandUebersicht}
\textbf{Totale Relation}
\UebFormel{\begin{aligned}
&\TotRel{R,A,B}:\\
&R\subseteq A\times B,\\
&x\in A\vdash\exists y\,(x,y)\in R.
\end{aligned}}
Die letzten beiden Zeilen sind die definierenden Strukturaxiome.
\UebQuellen{\FormulaRefAuto{Totale Relation}[def];
\FormulaRefAuto{x\in A \vdash \exists y\,(x,y)\in F}[ax]}
&
\textit{Modell aus einer disjunkten Familie}
Die Mitgliedschaftsrelation rechts unten erfüllt diese beiden
Anforderungen:
\UebFormel{\DisjFam(M)\vdash\TotRel{\MemRel(M),M,\bigcup M}.}
\UebQuellen{\FormulaRefAuto{\DisjFam(M)\vdash \TotRel{\MemRel(M),M,\bigcup M}}[thm]}\\
\addlinespace[.8ex]
\textbf{Mitgliedschaftsrelation}
\UebFormel{\begin{aligned}
\MemRel(M)\coloneqq{}&\{(X,a)\in M\times\bigcup M\\
&\mid a\in X\}.
\end{aligned}}
\UebQuellen{\FormulaRefAuto{\MemRel(M) \coloneqq \{\, (X,a)\in M\times \bigcup M \mid a\in X \,\}}[def]}
&
\textit{Elemente und Totalität}
\UebFormel{\begin{aligned}
&\MemRel(M)\subseteq M\times\bigcup M,\\
&X\in M\dsep a\in X\vdash(X,a)\in\MemRel(M),\\
&\DisjFam(M)\dsep X\in M\\
&\hspace{8mm}\vdash\exists a\,((X,a)\in\MemRel(M)).
\end{aligned}}
\UebQuellen{\FormulaRefAuto{\MemRel(M) \subseteq M\times \bigcup M}[thm];
\FormulaRefAuto{X\in M \dsep a\in X \vdash (X,a)\in \MemRel(M)}[thm];
\FormulaRefAuto{\DisjFam(M)\dsep X\in M \vdash \exists a\,\bigl((X,a)\in \MemRel(M)\bigr)}[thm]}\\
\addlinespace[.8ex]
\textbf{Strikte Obermengenrelation}
Für \(B\subseteq\powerset(M)\):
\UebFormel{\StrictSupRel{B}\coloneqq
\{(x,y)\in B\times B\mid x\subset y\}.}
\UebQuellen{\FormulaRefAuto{B\subseteq \powerset(M)\vdash \StrictSupRel{B}\coloneqq \{(x,y)\in B\times B\mid x\subset y\}}[def]}
&
\textit{Typisierung und Totalitätskriterium}
Unter derselben Voraussetzung:
\UebFormel{\begin{aligned}
&\StrictSupRel{B}\subseteq B\times B,\\
&\forall x\in B\,\exists y\in B\,(x\subset y)\\
&\hspace{10mm}\vdash\TotRel{\StrictSupRel{B},B,B}.
\end{aligned}}
\UebQuellen{\FormulaRefAuto{B\subseteq \powerset(M)\vdash \StrictSupRel{B}\subseteq B\times B}[thm];
\FormulaRefAuto{B\subseteq \powerset(M)\dsep \forall x\in B\,\exists y\in B\,\bigl(x\subset y\bigr) \vdash \TotRel{\StrictSupRel{B},B,B}}[thm]}\\
\end{BandUebersicht}

\chapter[Band 5: Funktionen]{Band 5 auf einen Blick}
\noindent\textbf{Funktionen}\par\medskip
\noindent\emph{Lesart.} Sofern eine Zeile nichts anderes angibt,
ist \(F\colon A\to B\) eine Funktion, \(C,D\subseteq A\) und
\(S,T\subseteq B\). Funktionsgraphen werden als Mengen geordneter Paare
verstanden.
\begin{BandUebersicht}
\textbf{Funktion}
\UebFormel{\begin{aligned}
&F\colon A\to B:\quad\TotRel{F,A,B},\\
&(x,y)\in F\dsep(x,z)\in F\vdash y=z.
\end{aligned}}
Zur Totalität tritt die funktionale Eindeutigkeit.
\UebQuellen{\FormulaRefAuto{Funktion}[def];
\FormulaRefAuto{(x,y)\in F\dsep (x,z)\in F \vdash y=z}[ax]}
&
\textit{Eindeutiger Wert und eindeutiger Definitionsbereich}
\UebFormel{\begin{aligned}
&x\in A\vdash\exists!y\;(x,y)\in F,\\
&F\colon A\to B\dsep F\colon C\to D\vdash A=C.
\end{aligned}}
\UebQuellen{\FormulaRefAuto{F\colon A\to B\dsep x\in A\vdash \exists! y\;(x,y)\in F}[thm];
\FormulaRefAuto{FunctionDomainUnique}[thm]}\\
\addlinespace[.8ex]
\textbf{Funktionswert und Graph}
Für \(x\in A\) ist \(F(x)\) das eindeutig bestimmte \(y\) mit
\((x,y)\in F\).
\UebFormel{\Graph(F)\coloneq
\{(x,y)\in A\times B\mid y=F(x)\}.}
\UebQuellen{\FormulaRefAuto{\Graph(F)\coloneq\{(x,y) \in A \times B \mid y=F(x)\}}[def]}
&
\textit{Typisierung und Funktionsextensionalität}
\UebFormel{\begin{aligned}
&x\in A\vdash F(x)\in B,\\
&F,G\colon A\to B\dsep\\
&\quad\forall x\in A\,(F(x)=G(x))\vdash F=G.
\end{aligned}}
\UebQuellen{\FormulaRefAuto{F\colon A\to B\dsep x\in A\vdash F(x)\in B}[thm];
\FormulaRefAuto{FunctionExtensionality}[thm]}\\
\addlinespace[.8ex]
\textbf{Bildmenge}
\(F[C]\) wird durch eindeutige Beschreibung eingeführt:
\UebFormel{y\in F[C]\leftrightarrow\exists x\in C\,(y=F(x)).}
\UebQuellen{\FormulaRefAuto{C\subseteq A\dsep F\colon A\to B\vdash F[C]\coloneqq \iota D\Bigl(\forall y\,\bigl(y\in D \leftrightarrow \exists x\in C\,y=F(x)\bigr)\Bigr)}[def]}
&
\textit{Existenz, Nichtleerheit und Vereinigungen}
\UebFormel{\begin{aligned}
&\exists!Y\,\forall y\,\bigl(y\in Y\\
&\qquad\leftrightarrow\exists x\in C\,(y=F(x))\bigr),\\
&C\in\PowNE{A}\vdash F[C]\in\PowNE{B},\\
&F[C\cup D]=F[C]\cup F[D].
\end{aligned}}
\(\PowNE{A}\) bezeichnet die nichtleeren Teilmengen von \(A\).
\UebQuellen{\FormulaRefAuto{SubsetFunctionImageSetUniqueExists}[thm];
\FormulaRefAuto{NonemptyDirectImage}[thm];
\FormulaRefAuto{F\colon M\to N \dsep A\subseteq M \dsep B\subseteq M \vdash F[A\cup B]=F[A]\cup F[B]}[thm]}\\
\addlinespace[.8ex]
\textbf{Urbild und Faser}
\UebFormel{\begin{aligned}
F^{-1}[T]&\coloneqq\{x\in A\mid F(x)\in T\},\\
\Fib_F(y)&\coloneqq F^{-1}[\{y\}]\quad(y\in B).
\end{aligned}}
\UebQuellen{\FormulaRefAuto{F\colon A\to B\dsep C\subseteq B \vdash F^{-1}[C] \coloneqq \{\,x\in A \mid F(x)\in C\,\}}[def];
\FormulaRefAuto{FiberMapDef}[def]}
&
\textit{Urbildkriterium und Rechengesetze}
\UebFormel{\begin{aligned}
&\forall x\in A\,(F(x)\in T\leftrightarrow x\in C)\\
&\hspace{10mm}\vdash F^{-1}[T]=C,\\
&F^{-1}[S\cup T]=F^{-1}[S]\cup F^{-1}[T],\\
&F^{-1}[S\cap T]=F^{-1}[S]\cap F^{-1}[T].
\end{aligned}}
Die Fasern zerlegen den Definitionsbereich nach gleichen Werten;
leere Fasern sind bei allgemeinen Funktionen erlaubt.
\UebQuellen{\FormulaRefAuto{PreimageEqFromPointwiseIff}[thm];
\FormulaRefAuto{F\colon A\to B \dsep C\subseteq B \dsep D\subseteq B \vdash F^{-1}[C\cup D]=F^{-1}[C]\cup F^{-1}[D]}[thm];
\FormulaRefAuto{F\colon A\to B \dsep C\subseteq B \dsep D\subseteq B \vdash F^{-1}[C\cap D]=F^{-1}[C]\cap F^{-1}[D]}[thm]}\\
\addlinespace[.8ex]
\textbf{Funktion aus einem Graphen}
Eine Relation \(R\subseteq A\times B\) kann eine Funktion beschreiben,
wenn ihre Werte für jedes Argument eindeutig existieren.
&
\textit{Konstruktionsprinzip}
\UebFormel{\begin{aligned}
&R\subseteq A\times B\dsep\\
&\forall x\in A\,\exists!y\in B\,((x,y)\in R)\\
&\hspace{10mm}\vdash R\colon A\to B.
\end{aligned}}
\UebQuellen{\FormulaRefAuto{UniqueValuedGraphFunction}[thm]}\\
\addlinespace[.8ex]
\textbf{Konstante und Fallabbildung}
Für \(b\in B\), \(f\colon C\to B\) und \(g\colon D\to B\):
\UebFormel{\begin{aligned}
&\ConstFun{A}{B}{b}(x)\coloneqq b\quad(x\in A),\\
&\CaseFun{C}{D}{f}{g}(x)\coloneqq\\
&\qquad\IfThenElse{x\in C}{f(x)}{g(x)}.
\end{aligned}}
Die zweite Formel gilt auf \(C\cup D\).
\UebQuellen{\FormulaRefAuto{Konstante Funktion}[def];
\FormulaRefAuto{CaseFunctionDef}[def]}
&
\textit{Elementarer Graph und Auswertung}
\UebFormel{\begin{aligned}
&b\in B\vdash\{(a,b)\}\colon\{a\}\to B,\\
&b\in B\vdash\{(a,b)\}(a)=b,\\
&f\colon A\to C\dsep g\colon B\to C\dsep x\in A\\
&\hspace{6mm}\vdash\CaseFun{A}{B}{f}{g}(x)=f(x).
\end{aligned}}
\UebQuellen{\FormulaRefAuto{SingletonGraphFunction}[thm];
\FormulaRefAuto{SingletonGraphValue}[thm];
\FormulaRefAuto{f\colon A\to C\dsep g\colon B\to C\dsep x\in A \vdash \CaseFun{A}{B}{f}{g}(x)=f(x)}[thm]}\\
\addlinespace[.8ex]
\textbf{Einschränkung}
\UebFormel{\begin{aligned}
&F\restriction_C\colon C\to B,\\
&(F\restriction_C)(x)\coloneqq F(x)\quad(x\in C).
\end{aligned}}
\UebQuellen{\FormulaRefAuto{RestrictionDef}[def]}
&
\textit{Gleichheitskriterium}
Für \(G\colon C\to B\):
\UebFormel{\begin{aligned}
F\restriction_C=G\ \leftrightarrow\ 
\forall x\in C\,(F(x)=G(x)).
\end{aligned}}
\UebQuellen{\FormulaRefAuto{FunctionRestrictionEqualityCriterion}[thm]}\\
\addlinespace[.8ex]
\textbf{Komposition}
Für \(F\colon A\to B\), \(G\colon B\to C\):
\UebFormel{\begin{aligned}
&G\circ F\colon A\to C,\\
&(G\circ F)(x)\coloneqq G(F(x))\quad(x\in A).
\end{aligned}}
\UebQuellen{\FormulaRefAuto{CompositionDef}[def]}
&
\textit{Bilder unter einer Linksinversen}
Für \(F\colon A\to B\), \(G\colon B\to A\), \(K\subseteq A\):
\UebFormel{\begin{aligned}
&\forall x\in A\,G(F(x))=x\\
&\hspace{12mm}\vdash G[F[K]]=K.
\end{aligned}}
\UebQuellen{\FormulaRefAuto{LeftInverseImages}[thm]}\\
\addlinespace[.8ex]
\textbf{Induzierte Potenzmengenabbildung}
\UebFormel{\begin{aligned}
&\PowMap{F}\colon\powerset(A)\to\powerset(B),\\
&\PowMap{F}(X)\coloneqq F[X],\\
&\PowMapNE{F}\coloneqq\PowMap{F}\restriction_{\PowNE{A}}.
\end{aligned}}
\UebQuellen{\FormulaRefAuto{PowerMapDef}[def];
\FormulaRefAuto{NonemptyPowerMapDef}[def]}
&
\textit{Nichtleere Teilmengen bleiben nichtleer}
\UebFormel{\begin{aligned}
&\PowMapNE{F}\colon\PowNE{A}\to\PowNE{B},\\
&(\PowMap{F})^{-1}[\PowNE{B}]=\PowNE{A}.
\end{aligned}}
\UebQuellen{\FormulaRefAuto{NonemptyPowerMapTyped}[thm];
\FormulaRefAuto{NonemptyPowerMapPreimage}[thm]}\\
\addlinespace[.8ex]
\textbf{Spurabbildung}
Für \(F\colon A\to\powerset(B)\):
\UebFormel{\begin{aligned}
&\TraceMap{F}\colon B\to\powerset(A),\\
&\TraceMap{F}(b)\coloneqq\{x\in A\mid b\in F(x)\}.
\end{aligned}}
\UebQuellen{\FormulaRefAuto{TraceMapDef}[def]}
&
\textit{Gleiche Spuren}
Für \(b,c\in B\), \(x\in A\):
\UebFormel{\begin{aligned}
&\TraceMap{F}(b)=\TraceMap{F}(c)\\
&\hspace{8mm}\vdash b\in F(x)\leftrightarrow c\in F(x).
\end{aligned}}
\UebQuellen{\FormulaRefAuto{TraceMapEqualityMembership}[thm]}\\
\addlinespace[.8ex]
\textbf{Erweiterung und überdeckende Familie}
Für \(a\notin A\), \(b\in B\) erweitert \(\ExtMap{F}{a}{b}\)
die Funktion durch den neuen Wert \(a\mapsto b\).
\UebFormel{\ExtMap{F}{a}{b}(x)\coloneqq
\IfThenElse{x=a}{b}{F(x)}.}
\UebQuellen{\FormulaRefAuto{ExtensionMapDef}[def]}
&
\textit{Vereinigung einer Überdeckung}
Für \(D\colon I\to\powerset(W)\):
\UebFormel{\begin{aligned}
&\forall w\in W\,\exists i\in I\,(w\in D(i))\\
&\hspace{12mm}\vdash\bigcup D[I]=W.
\end{aligned}}
\UebQuellen{\FormulaRefAuto{IndexedSubsetFamilyCoverUnion}[thm]}\\
\end{BandUebersicht}

\chapter[Band 6: Injektive Funktionen]{Band 6 auf einen Blick}
\noindent\textbf{Injektive Funktionen}\par\medskip
\noindent\emph{Lesart.} Dieser eigenständige Band fügt dem Funktionsbegriff
das Axiom der Injektivität hinzu. Für die Bild- und Urbildformeln gelten
die jeweils genannten Träger- und Teilmengenbedingungen.
\begin{BandUebersicht}
\textbf{Injektive Funktion}
\(F\colon A\inj B\) bedeutet: \(F\colon A\to B\) und
\UebFormel{x,y\in A\dsep F(x)=F(y)\vdash x=y.}
\UebQuellen{\FormulaRefAuto{Injektive Funktion}[def];
\FormulaRefAuto{Injektivität}[ax]}
&
\textit{Bildmitgliedschaft erkennt das ursprüngliche Element}
Für \(Y\subseteq A\), \(x\in A\):
\UebFormel{F(x)\in F[Y]\vdash x\in Y.}
\UebQuellen{\FormulaRefAuto{F\colon A\inj B \dsep Y\in\powerset(A)\dsep x\in A\dsep F(x)\in F[Y] \vdash x\in Y}[thm]}\\
\addlinespace[.8ex]
\textbf{Bilder unter einer Injektion}
Die Bildkonstruktion aus Band 5 wird auf eine injektive Funktion
angewandt. Sei \(F\colon M\inj N\), \(A,B\subseteq M\).
&
\textit{Schnitte werden erhalten; gleiche Bilder sind eindeutig}
\UebFormel{\begin{aligned}
&F[A\cap B]=F[A]\cap F[B],\\
&F[A]=F[B]\vdash A=B.
\end{aligned}}
\UebQuellen{\FormulaRefAuto{F\colon M\inj N \dsep A\subseteq M \dsep B\subseteq M \vdash F[A\cap B]=F[A]\cap F[B]}[thm];
\FormulaRefAuto{F\colon A\inj B \dsep X\in\powerset(A)\dsep Y\in\powerset(A)\dsep F[X]=F[Y] \vdash X=Y}[thm]}\\
\addlinespace[.8ex]
\textbf{Einschränkungen und Zielmengenwechsel}
\UebFormel{(F\restriction_C)(x)\coloneqq F(x)\quad(x\in C).}
Sei \(F\colon A\inj B\), \(C\subseteq A\), \(T\subseteq B\).
\UebQuellen{\FormulaRefAuto{RestrictionDef}[def]}
&
\textit{Injektivität bleibt erhalten}
\UebFormel{\begin{aligned}
&F\restriction_C\colon C\inj B,\\
&F\restriction_{F^{-1}[T]}\colon F^{-1}[T]\inj T,\\
&B\subseteq D\vdash F\colon A\inj D.
\end{aligned}}
\UebQuellen{\FormulaRefAuto{F\colon A\inj B\dsep C\subseteq A \vdash F\restriction_C\colon C\inj B}[thm];
\FormulaRefAuto{F\colon A\inj B\dsep T\subseteq B \vdash F\restriction_{F^{-1}[T]}\colon F^{-1}[T]\inj T}[thm];
\FormulaRefAuto{InjectiveCodomainExtension}[thm]}\\
\addlinespace[.8ex]
\textbf{Inklusionsabbildung und Komposition}
Für \(A\subseteq B\):
\UebFormel{\iota_{A,B}(x)\coloneqq x\quad(x\in A).}
Für kompatible Funktionen gilt
\UebFormel{(G\circ F)(x)\coloneqq G(F(x)).}
\UebQuellen{\FormulaRefAuto{InclusionMapDef}[def];
\FormulaRefAuto{CompositionDef}[def]}
&
\textit{Komposition von Injektionen}
\UebFormel{\begin{aligned}
&F\colon A\inj B\dsep G\colon B\inj C\\
&\hspace{12mm}\vdash G\circ F\colon A\inj C,\\
&B\subseteq C\dsep F\colon A\to B\dsep x\in A\\
&\hspace{8mm}\vdash(\iota_{B,C}\circ F)(x)=F(x).
\end{aligned}}
\UebQuellen{\FormulaRefAuto{F\colon A\inj B \dsep G\colon B\inj C\vdash G\circ F\colon A\inj C}[thm];
\FormulaRefAuto{InclusionCompositionValue}[thm]}\\
\addlinespace[.8ex]
\textbf{Fallabbildung auf disjunkten Bereichen}
Für \(A\cap B=\varnothing\), \(f\colon A\inj C\),
\(g\colon B\inj C\) verbindet \(\CaseFun{A}{B}{f}{g}\)
die beiden Zweige.
\UebQuellen{\FormulaRefAuto{CaseFunctionDef}[def]}
&
\textit{Disjunkte Bilder verhindern Kollisionen}
\UebFormel{\begin{aligned}
&f[A]\cap g[B]=\varnothing\\
&\hspace{8mm}\vdash\CaseFun{A}{B}{f}{g}\colon A\cup B\inj C.
\end{aligned}}
\UebQuellen{\FormulaRefAuto{CaseFunInjectiveDisjointImages}[thm]}\\
\addlinespace[.8ex]
\textbf{Adjunktion in einer Mengenfamilie}
\UebFormel{\begin{aligned}
\FamContaining{M}{a}&\coloneqq\{X\in M\mid a\in X\},\\
\FamAvoiding{M}{a}&\coloneqq\{X\in M\mid a\notin X\}.
\end{aligned}}
Für \(\UnionClosedFam(M)\), \(\{a\}\in M\) wird auf der
vermeidenden Teilfamilie definiert:
\UebFormel{\FranklAdj{M}{a}(X)\coloneqq X\cup\{a\}.}
\UebQuellen{\FormulaRefAuto{Enthaltende Teilfamilie}[def];
\FormulaRefAuto{Vermeidende Teilfamilie}[def];
\FormulaRefAuto{Adjunktionsabbildung einer Mengenfamilie}[def]}
&
\textit{Adjunktion ist injektiv}
Unter den links genannten Voraussetzungen:
\UebFormel{\FranklAdj{M}{a}\colon
\FamAvoiding{M}{a}\inj\FamContaining{M}{a}.}
\UebQuellen{\FormulaRefAuto{\UnionClosedFam(M)\dsep \{a\}\in M \vdash \FranklAdj{M}{a}\colon \FamAvoiding{M}{a}\inj \FamContaining{M}{a}}[thm]}\\
\addlinespace[.8ex]
\textbf{Injektiver Größenvergleich}
Für \(B\subseteq A\):
\UebFormel{\AtMostHalf{B}{A}\coloneqq
\exists F\colon B\inj A\setminus B.}
Dies definiert eine injektive Fassung von \enquote{höchstens der Hälfte}.
\UebQuellen{\FormulaRefAuto{AtMostHalf}[def]}
&
\textit{Elementarer injektiver Vergleich}
\UebFormel{b\in B\vdash\exists F\colon\{a\}\inj B.}
Eine Einermenge injiziert also in jede nichtleere Menge.
\UebQuellen{\FormulaRefAuto{SingletonInjectionIntoNonempty}[thm]}\\
\end{BandUebersicht}

\chapter[Band 7: Surjektive Funktionen]{Band 7 auf einen Blick}
\noindent\textbf{Surjektive Funktionen}\par\medskip
\noindent\emph{Lesart.} Surjektivität ist ein eigenes Strukturaxiom
für \(F\colon A\to B\); sie hängt ausdrücklich von der Zielmenge \(B\) ab.
\begin{BandUebersicht}
\textbf{Surjektive Funktion}
\(F\colon A\sur B\) bedeutet \(F\colon A\to B\) zusammen mit
\UebFormel{y\in B\vdash\exists x\in A\;F(x)=y.}
\UebQuellen{\FormulaRefAuto{Surjektive Funktion}[def];
\FormulaRefAuto{Surjektivität}[ax]}
&
\textit{Surjektivität ist gleichbedeutend mit vollem Bild}
Für \(F\colon A\to B\):
\UebFormel{\begin{aligned}
&F\colon A\sur B\leftrightarrow F[A]=B,\\
&F\colon A\sur F[A].
\end{aligned}}
\UebQuellen{\FormulaRefAuto{SurjectiveIffImageEqualsCodomain}[thm];
\FormulaRefAuto{CorestrictionToImageSurjective}[thm]}\\
\addlinespace[.8ex]
\textbf{Potenzmenge und Diagonalmenge}
Für \(F\colon A\to\powerset(A)\) betrachtet der Beweis
\UebFormel{D\coloneqq\{x\in A\mid x\notin F(x)\}.}
Dies ist eine Aussonderung in \(A\), keine zusätzliche Axiomannahme.
&
\textit{Cantors Satz}
\UebFormel{\begin{aligned}
&F\colon A\to\powerset(A)\\
&\hspace{10mm}\vdash\neg(F\colon A\sur\powerset(A)).
\end{aligned}}
Keine Funktion auf einer Menge trifft alle ihre Teilmengen.
\UebQuellen{\FormulaRefAuto{CantorNoPowerSetSurjection}[thm]}\\
\addlinespace[.8ex]
\textbf{Surjektion aus einer Injektion}
Für \(F\colon A\inj B\), \(a_0\in A\) wird
\(\Gsurjfrominj{F}{a_0}\) durch Rückzuordnung auf \(F[A]\)
und den festen Wert \(a_0\) außerhalb von \(F[A]\) definiert:
\UebFormel{\begin{aligned}
&(b,a)\in\Gsurjfrominj{F}{a_0}\ \leftrightarrow\\
&b\in B\land a\in A\land\bigl(\\
&\quad(b\in F[A]\land F(a)=b)\\
&\quad\lor(b\notin F[A]\land a=a_0)\bigr).
\end{aligned}}
\UebQuellen{\FormulaRefAuto{F\colon A\inj B \vdash \Gsurjfrominj{F}{a_0} \coloneqq \{\, (b,a)\in B\times A \mid (b\in F[A]\land F(a)=b)\lor (b\notin F[A]\land a=a_0)\,\}}[def]}
&
\textit{Die Rückzuordnung ist surjektiv}
\UebFormel{\begin{aligned}
&F\colon A\inj B\dsep a_0\in A\\
&\hspace{10mm}\vdash\Gsurjfrominj{F}{a_0}\colon B\sur A.
\end{aligned}}
Die ausdrücklich verlangte Wahl von \(a_0\) behandelt die Werte
außerhalb des Bildes.
\UebQuellen{\FormulaRefAuto{F\colon A\inj B \dsep a_0\in A \vdash \Gsurjfrominj{F}{a_0}\colon B\sur A}[thm]}\\
\addlinespace[.8ex]
\textbf{Projektionen}
Für \((x,y)\in A\times B\):
\UebFormel{\pi_1(x,y)\coloneqq x,\qquad
\pi_2(x,y)\coloneqq y.}
\UebQuellen{\FormulaRefAuto{FirstProjectionDef}[def];
\FormulaRefAuto{SecondProjectionDef}[def]}
&
\textit{Surjektivität bei nichtleerem anderem Faktor}
\UebFormel{\begin{aligned}
&b_0\in B\vdash\pi_1\colon A\times B\sur A,\\
&a_0\in A\vdash\pi_2\colon A\times B\sur B.
\end{aligned}}
\UebQuellen{\FormulaRefAuto{FirstProjectionSurjective}[thm];
\FormulaRefAuto{SecondProjectionSurjective}[thm]}\\
\addlinespace[.8ex]
\textbf{Fasern und Einschränkung auf eine Relation}
\UebFormel{\Fib_F(y)\coloneqq F^{-1}[\{y\}].}
Für eine totale Relation \(\TotRel{R,A,B}\) wird außerdem
\(\pi_1\) auf \(R\) eingeschränkt.
\UebQuellen{\FormulaRefAuto{FiberMapDef}[def];
\FormulaRefAuto{RestrictionDef}[def]}
&
\textit{Nichtleere Fasern und surjektive Relationsprojektion}
\UebFormel{\begin{aligned}
&F\colon A\sur B\dsep y\in B\vdash\Fib_F(y)\neq\varnothing,\\
&\TotRel{R,A,B}\vdash\pi_1\restriction_R\colon R\sur A.
\end{aligned}}
\UebQuellen{\FormulaRefAuto{F\colon A\sur B \dsep y\in B \vdash \Fib_F(y)\neq\varnothing}[thm];
\FormulaRefAuto{\pi_1\restriction_R\colon R\sur A}[thm]}\\
\addlinespace[.8ex]
\textbf{Komposition}
Für \(F\colon A\sur B\), \(G\colon B\sur C\):
\UebFormel{(G\circ F)(x)\coloneqq G(F(x))\quad(x\in A).}
\UebQuellen{\FormulaRefAuto{CompositionDef}[def]}
&
\textit{Surjektivität bleibt erhalten}
\UebFormel{G\circ F\colon A\sur C.}
Auch der leere Randfall ist ausgearbeitet:
\UebFormel{\varnothing\colon\varnothing\sur\varnothing.}
\UebQuellen{\FormulaRefAuto{G\circ F\colon A\sur C}[thm];
\FormulaRefAuto{\varnothing\colon \varnothing\sur \varnothing}[thm]}\\
\end{BandUebersicht}

\chapter[Band 8: Bijektive Funktionen]{Band 8 auf einen Blick}
\noindent\textbf{Bijektive Funktionen}\par\medskip
\noindent\emph{Lesart.} Bijektionen verbinden die beiden getrennt
entwickelten Axiome der Injektivität und Surjektivität. \(F^{-1}\)
bezeichnet hier die Umkehrfunktion einer Bijektion; \(F^{-1}[T]\)
bleibt die Urbildmenge.
\begin{BandUebersicht}
\textbf{Bijektive Funktion}
\UebFormel{\begin{aligned}
&F\colon A\bij B:\quad F\colon A\to B,\\
&F\colon A\inj B,\qquad F\colon A\sur B.
\end{aligned}}
\UebQuellen{\FormulaRefAuto{Bijektive Funktion}[def]}
&
\textit{Komposition von Bijektionen}
\UebFormel{\begin{aligned}
&F\colon A\bij B\dsep G\colon B\bij C\\
&\hspace{12mm}\vdash G\circ F\colon A\bij C.
\end{aligned}}
\UebQuellen{\FormulaRefAuto{BijectiveFunctionComposition}[thm]}\\
\addlinespace[.8ex]
\textbf{Identität und Umkehrfunktion}
\(\Id_A\) ordnet jedem \(x\in A\) sich selbst zu. Für
\(F\colon A\bij B\) wird definiert:
\UebFormel{F^{-1}(y)\coloneqq
\iota x\,(x\in A\land F(x)=y)\quad(y\in B).}
\UebQuellen{\FormulaRefAuto{InverseFunctionDef}[def]}
&
\textit{Beidseitige Umkehrung}
\UebFormel{\begin{aligned}
&F\colon A\bij B\vdash F^{-1}\colon B\bij A,\\
&F\colon A\bij B\dsep x\in A\vdash{}\\[-2pt]
&\qquad F^{-1}(F(x))=x,\\
&F\colon A\bij B\dsep y\in B\vdash{}\\[-2pt]
&\qquad F(F^{-1}(y))=y.
\end{aligned}}
\UebQuellen{\FormulaRefAuto{InverseFunctionBijection}[thm];
\FormulaRefAuto{InverseFunctionLeftInverse}[thm];
\FormulaRefAuto{InverseFunctionRightInverse}[thm]}\\
\addlinespace[.8ex]
\textbf{Bijektive Einschränkungen}
Sei \(F\colon A\bij B\), \(S\subseteq A\), \(T\subseteq B\).
Die Einschränkung verwendet weiterhin die Werte von \(F\).
\UebQuellen{\FormulaRefAuto{RestrictionDef}[def]}
&
\textit{Punktweises Kriterium}
\UebFormel{\begin{aligned}
&F\restriction_S\colon S\bij T\ \leftrightarrow\\
&\quad\forall x\in A\,(x\in S\leftrightarrow F(x)\in T),\\
&F\restriction_{F^{-1}[T]}\colon F^{-1}[T]\bij T.
\end{aligned}}
\UebQuellen{\FormulaRefAuto{BijectiveRestrictionIffPointwiseCompatibility}[thm];
\FormulaRefAuto{F\colon A\bij B \dsep T\subseteq B \vdash F\restriction_{F^{-1}[T]}\colon F^{-1}[T]\bij T}[thm]}\\
\addlinespace[.8ex]
\textbf{Induzierte Potenzmengenabbildung}
\UebFormel{\PowMap{F}(X)\coloneqq F[X].}
Hier ist \(F\colon A\to B\) zunächst nur eine Funktion.
\UebQuellen{\FormulaRefAuto{PowerMapDef}[def]}
&
\textit{Bijektivität wird übertragen und erkannt}
\UebFormel{\begin{aligned}
&F\colon A\bij B\vdash
\PowMap{F}\colon\powerset(A)\bij\powerset(B),\\
&\PowMap{F}\colon\powerset(A)\bij\powerset(B)\\
&\hspace{12mm}\vdash F\colon A\bij B.
\end{aligned}}
\UebQuellen{\FormulaRefAuto{F\colon A\bij B \vdash \PowMap{F}\colon \powerset(A)\bij \powerset(B)}[thm];
\FormulaRefAuto{PowerMapBijectiveReflectsBijectivity}[thm]}\\
\addlinespace[.8ex]
\textbf{Verkleben disjunkter Zweige}
Für \(A_0\cap A_1=\varnothing\), \(B_0\cap B_1=\varnothing\)
und \(f_i\colon A_i\bij B_i\) wird die Fallabbildung verwendet.
\UebQuellen{\FormulaRefAuto{CaseFunctionDef}[def]}
&
\textit{Bijektive Fallabbildung}
\UebFormel{\begin{aligned}
&\CaseFun{A_0}{A_1}{f_0}{f_1}\\
&\hspace{10mm}\colon A_0\cup A_1\bij B_0\cup B_1.
\end{aligned}}
\UebQuellen{\FormulaRefAuto{CaseFunBijectiveDisjointTargets}[thm]}\\
\addlinespace[.8ex]
\textbf{Kernfamilien und Kerntransport}
\UebFormel{\begin{aligned}
&\CoreFamily{P}{A}\coloneqq\\
&\qquad\{H\in\powerset(P\cup A)\mid P\subseteq H\},\\
&\CoreTransport{P}{Q}{f}(H)\coloneqq Q\cup f[H\setminus P].
\end{aligned}}
Für den Transport gilt \(P\cap A=Q\cap B=\varnothing\),
\(f\colon A\to B\), \(H\in\CoreFamily{P}{A}\).
\UebQuellen{\FormulaRefAuto{CoreFamilyDef}[def];
\FormulaRefAuto{CoreTransportDef}[def]}
&
\textit{Bijektiver Transport}
Unter den links genannten Disjunktheitsbedingungen:
\UebFormel{\begin{aligned}
&f\colon A\bij B\vdash\\
&\CoreTransport{P}{Q}{f}\colon
\CoreFamily{P}{A}\bij\CoreFamily{Q}{B}.
\end{aligned}}
\UebQuellen{\FormulaRefAuto{CoreFamilyTransportBijective}[thm]}\\
\addlinespace[.8ex]
\textbf{Schichtfortsetzung}
Für \(g\colon\powerset(D)\to\powerset(E)\),
\(X\in\powerset(A\cup D)\):
\UebFormel{\begin{aligned}
\LayerPowerMap{A}{D}{E}{g}(X)
\coloneqq(X\cap A)\cup g(X\cap D).
\end{aligned}}
\UebQuellen{\FormulaRefAuto{LayerPowerMapDef}[def]}
&
\textit{Bijektive Fortsetzung bei getrennten Schichten}
Für \(A\cap D=A\cap E=\varnothing\) und
\(g\colon\powerset(D)\bij\powerset(E)\):
\UebFormel{\begin{aligned}
\LayerPowerMap{A}{D}{E}{g}\colon
\powerset(A\cup D)\bij\powerset(A\cup E).
\end{aligned}}
Falls außerdem \(g(\varnothing)=\varnothing\), ist auch die
Einschränkung auf die nichtleeren Teilmengen bijektiv.
\UebQuellen{\FormulaRefAuto{LayerPowerMapBijective}[thm];
\FormulaRefAuto{LayerPowerMapNonemptyBijective}[thm]}\\
\addlinespace[.8ex]
\textbf{Cantor--Schröder--Bernstein}
Für \(F\colon A\inj B,\ G\colon B\inj A\) ist
\(\CBPart F G\) die Menge aller \(x\in A\), die zu jedem
\(X\subseteq A\) mit
\UebFormel{\begin{aligned}
&A\setminus G[B]\subseteq X,\\
&G[F[X]]\subseteq X
\end{aligned}}
gehören.
\UebQuellen{\FormulaRefAuto{CantorBernsteinPartDef}[def]}
&
\textit{Fixpunkt und Zusammenfügung}
Mit \(C=\CBPart F G\) gilt
\UebFormel{C=(A\setminus G[B])\cup G[F[C]].}
Auf \(C\) wird \(F\) verwendet, auf \(A\setminus C\) der
eindeutige inverse Wert von \(G\). Dies liefert
\UebFormel{\exists H\colon A\bij B.}
\UebQuellen{\FormulaRefAuto{CantorBernsteinPartFixedPoint}[thm];
\FormulaRefAuto{CantorBernsteinFixedInjections}[thm]}
\\ \addlinespace[.8ex]

\textbf{Erweiterung um ein neues Paar}
Für \(F\colon A\bij B\), \(a\notin A\), \(b\notin B\)
wird die Erweiterungsabbildung auf den vergrößerten Zielbereich angewandt.
\UebQuellen{\FormulaRefAuto{ExtensionMapDef}[def]}
&
\textit{Frische Erweiterung einer Bijektion}
\UebFormel{\begin{aligned}
&\ExtMap{\iota_{B,B\cup\{b\}}\circ F}{a}{b}\\
&\hspace{12mm}\colon A\cup\{a\}\bij B\cup\{b\}.
\end{aligned}}
\UebQuellen{\FormulaRefAuto{FreshExtMapBijective}[thm]}\\
\addlinespace[.8ex]
\textbf{Vorgeschriebener Wert}
Eine vorhandene Bijektion kann durch eine Vertauschung im Ziel
an ein ausgewähltes Paar angepasst werden.
&
\textit{Punktierte Bijektion}
\UebFormel{\begin{aligned}
&\exists u\,(u\colon A\bij B)\dsep a\in A\dsep b\in B\\
&\hspace{6mm}\vdash\exists f\,(f\colon A\bij B\land f(a)=b).
\end{aligned}}
\UebQuellen{\FormulaRefAuto{PointedBijectionExists}[thm]}\\
\end{BandUebersicht}
\noindent\emph{Lesepfad zu Cantor--Schröder--Bernstein.}
Die \CSBReadingLink{} erläutert den kleinsten abgeschlossenen Bereich,
seine Fixpunktgleichung und die Zusammenfügung der beiden Bijektionen.
Die zugeordneten \CSBProofsLink{} enthalten die ausführlichen
Ableitungen zu denselben Resultaten.

\chapter[Band 9: Auswahlprinzip]{Band 9 auf einen Blick}
\noindent\textbf{Auswahlprinzip}\par\medskip
\noindent\emph{Lesart.} Die Surjektions-, Relations- und Familienform
des Auswahlprinzips werden miteinander verknüpft. Die Aussagen über
beliebige gleichzeitige Auswahlen verwenden ausdrücklich das Auswahlaxiom;
die Injektivität einer bereits vorhandenen Rechtsinversen benötigt es nicht.
\begin{BandUebersicht}
\textbf{Rechtsinverse}
Für \(F\colon A\to B\), \(G\colon B\to A\) bedeutet
\UebFormel{F\circ G=\Id_B,}
dass \(G\) eine Rechtsinverse von \(F\) ist.
&
\textit{Eine Rechtsinverse ist injektiv}
\UebFormel{F\circ G=\Id_B\vdash G\colon B\inj A.}
\UebQuellen{\FormulaRefAuto{F\circ G=\Id_B\vdash G\colon B\inj A}[thm]}\\
\addlinespace[.8ex]
\textbf{Auswahlaxiom in Surjektionsform}
Für jede Surjektion \(G\colon B\sur A\) wird gefordert:
\UebFormel{\exists F\colon A\to B\;G\circ F=\Id_A.}
\(\ACsur\) bezeichnet die entsprechende global quantifizierte Aussage.
\UebQuellen{\FormulaRefAuto{Auswahlaxiom}[ax];
\FormulaRefAuto{ACsur}[def]}
&
\textit{Auswahl aus einer totalen Relation}
Unter AC und \(\TotRel{R,A,B}\):
\UebFormel{\begin{aligned}
&\exists F\colon A\to B\,\forall x\in A\\
&\hspace{12mm}((x,F(x))\in R).
\end{aligned}}
\UebQuellen{\FormulaRefAuto{\exists F\colon A \to B\,\forall x\in A\,\bigl((x,F(x))\in R\bigr)}[thm]}\\
\addlinespace[.8ex]
\textbf{Disjunkte Familien und Auswahlmengen}
Für \(\DisjFam(M)\) soll \(L\) jedes Glied von \(M\) genau einmal
treffen. Die Mitgliedschaftsrelation aus Band 4 vermittelt zwischen
Familien- und Relationsform.
\UebQuellen{\FormulaRefAuto{Disjunkte Familie}[def]}
&
\textit{Familienform des Auswahlprinzips}
Unter AC und \(\DisjFam(M)\):
\UebFormel{\exists L\,\forall X\in M\,\exists!a\in X\,(a\in L).}
Der Band beweist außerdem die Rückrichtung zur Surjektionsform.
\UebQuellen{\FormulaRefAuto{\exists L\,\forall X\in M\,\exists! a\in X\,\bigl(a\in L\bigr)}[thm];
\FormulaRefAuto{\forall M\,\bigl(\DisjFam(M)\rightarrow \exists L\,\forall X\in M\,\exists! a\,(a\in X \land a\in L)\bigr) \vdash \exists F\colon A\to B\,\bigl(G\circ F=\Id_A\bigr)}[thm]}\\
\addlinespace[.8ex]
\textbf{Sektion aus einer Auswahlmenge}
Sei \(G\colon B\sur A\); \(L\) treffe jede Faser
\(X\in\Fib_G[A]\) genau einmal. Dann:
\UebFormel{\Sec_{G,L}(x)\coloneqq
\iota b\,(b\in\Fib_G(x)\land b\in L).}
\UebQuellen{\FormulaRefAuto{ChoiceSectionDef}[def]}
&
\textit{Die gewählte Sektion liefert eine Rechtsinverse}
\UebFormel{\begin{aligned}
&\forall X\in\Fib_G[A]\,\exists!a\,(a\in X\land a\in L)\\
&\hspace{6mm}\vdash\exists F\colon A\to B\,(G\circ F=\Id_A).
\end{aligned}}
\UebQuellen{\FormulaRefAuto{ChoiceSetRightInverse}[thm]}\\
\addlinespace[.8ex]
\textbf{Surjektive Größenvergleiche}
Die Surjektionsform von AC garantiert die Existenz einer Sektion;
deren Injektivität wurde am Anfang des Bandes unabhängig gezeigt.
&
\textit{Injektive Rechtsinverse}
Unter AC:
\UebFormel{\begin{aligned}
&F\colon A\sur B\vdash\exists H\colon B\inj A\\
&\hspace{18mm}\forall y\in B\,(F(H(y))=y).
\end{aligned}}
\UebQuellen{\FormulaRefAuto{F\colon A\sur B\vdash \exists H\colon B\inj A\,\forall y\in B\,\bigl(F(H(y))=y\bigr)}[thm]}\\
\addlinespace[.8ex]
\textbf{Retraktionen und ihre Fasern}
Seien \(G\subseteq S\), \(H\subseteq T\),
\(r\colon S\to G\), \(s\colon T\to H\) mit
\(r(g)=g\), \(s(h)=h\), und \(\varphi\colon G\bij H\).
\UebFormel{Q_g\coloneqq r^{-1}[\{g\}],\qquad
R_h\coloneqq s^{-1}[\{h\}].}
&
\textit{Bijektives Verkleben der Retraktionsfasern}
Unter AC und den links genannten Voraussetzungen:
\UebFormel{\begin{aligned}
&\forall g\in G\,\exists u\,(u\colon Q_g\bij R_{\varphi(g)})\\
&\vdash\exists f\,\bigl(f\colon S\bij T\\
&\quad\land\forall x\in S\,s(f(x))=\varphi(r(x))\\
&\quad\land\forall g\in G\,f(g)=\varphi(g)\bigr).
\end{aligned}}
\UebQuellen{\FormulaRefAuto{RetractionFiberBijection}[thm]}\\
\end{BandUebersicht}

\chapter[Band 10: Natürliche Zahlen]{Band 10 auf einen Blick}
\noindent\textbf{Natürliche Zahlen}\par\medskip
\noindent\emph{Lesart.} Der Band unterscheidet die mengentheoretische
Konstruktion und den anschließenden Aufbau mit der Peano-Schnittstelle.
In den arithmetischen Zeilen sind \(a,b,c,m,n\in\mathbb N\);
\(u\colon\mathbb N\to\mathbb N\) ist eine Folge natürlicher Zahlen.
\begin{BandUebersicht}
\textbf{Induktive Menge und von-Neumann-Konstruktion}
\UebFormel{\begin{aligned}
&\Induktiv(A):\quad\varnothing\in A,\\
&\forall x\in A\,(x\cup\{x\}\in A),\\
&\mathbb N\coloneq\bigcap\{A\mid\Induktiv(A)\},\\
&0\coloneqq\varnothing.
\end{aligned}}
\UebQuellen{\FormulaRefAuto{InductiveSetDef}[def];
\FormulaRefAuto{\mathbb{N} \coloneq \bigcap \{ A \mid \Induktiv(A) \}}[def];
\FormulaRefAuto{0\coloneqq\varnothing}[def]}
&
\textit{Kleinste induktive Menge}
\UebFormel{\begin{aligned}
&\Induktiv(\mathbb N),\\
&\Induktiv(A)\vdash\mathbb N\subseteq A.
\end{aligned}}
Das Unendlichkeitsaxiom sichert, dass induktive Mengen existieren;
die Schnittkonstruktion liefert ihre gemeinsame kleinste Struktur.
\UebQuellen{\FormulaRefAuto{\Induktiv(\mathbb{N})}[thm];
\FormulaRefAuto{\Induktiv(A) \vdash \mathbb{N} \subseteq A}[thm]}\\
\addlinespace[.8ex]
\textbf{Peano-Schnittstelle}
\(0\in\mathbb N\), \(\Succ\colon\mathbb N\to\mathbb N\),
kein Nachfolger ist null, und \(\Succ\) ist injektiv. Dazu tritt Induktion:
\UebFormel{\begin{aligned}
&P(0)\dsep\forall n\in\mathbb N\,(P(n)\rightarrow P(\Succ(n)))\\
&\hspace{10mm}\vdash\forall n\in\mathbb N\,P(n).
\end{aligned}}
\UebQuellen{\FormulaRefAuto{Peano-Axiome}[def];
\FormulaRefAuto{PeanoInductionAxiom}[ax]}
&
\textit{Eindeutiger Vorgänger jeder Nichtnullzahl}
\UebFormel{\begin{aligned}
&n\neq0\vdash\exists!k\in\mathbb N\,(n=\Succ(k)),\\
&\Succ\colon\mathbb N\bij\mathbb N\setminus\{0\}.
\end{aligned}}
\UebQuellen{\FormulaRefAuto{PeanoNonzeroUniquePredecessor}[thm];
\FormulaRefAuto{PeanoSuccNonzeroBijection}[thm]}\\
\addlinespace[.8ex]
\textbf{Rekursionskern und Rekursionsabbildung}
Für \(m_0\in M\), \(f\colon M\to M\) wird der kleinste
Graphkandidat mit Startpaar \((0,m_0)\) und Abschluss unter
\((n,a)\mapsto(\Succ(n),f(a))\) konstruiert.
Die daraus gewonnene Funktion heißt \(\RecFun{m_0}{f}\).
\UebQuellen{\FormulaRefAuto{RecFunDef}[def]}
&
\textit{Dedekindscher Rekursionssatz}
\UebFormel{\begin{aligned}
&m_0\in M\dsep f\colon M\to M\vdash\\
&\exists!G\colon\mathbb N\to M\,\bigl(G(0)=m_0\\
&\quad\land\forall n\in\mathbb N\,
G(\Succ(n))=f(G(n))\bigr).
\end{aligned}}
\UebQuellen{\FormulaRefAuto{DedekindRecursionTheorem}[thm]}\\
\addlinespace[.8ex]
\textbf{Addition durch Rekursion}
\UebFormel{a+b\coloneqq\RecFun{a}{\Succ}(b).}
\UebQuellen{\FormulaRefAuto{Addition auf \(\mathbb{N}\)}[def]}
&
\textit{Assoziativität und Kommutativität}
\UebFormel{\begin{aligned}
&(a+b)+c=a+(b+c),\\
&a+b=b+a.
\end{aligned}}
Die rekursive Definition führt damit zu den üblichen Rechengesetzen.
\UebQuellen{\FormulaRefAuto{PeanoAddAssociative}[thm];
\FormulaRefAuto{PeanoAddCommutative}[thm]}\\
\addlinespace[.8ex]
\textbf{Nachfolger- und Vorgängerschreibweise}
\(1\) bezeichnet \(\Succ(0)\); \(\Pred=\Succ^{-1}\) wird auf
\(\mathbb N\setminus\{0\}\) definiert. Für \(n\neq0\) ist
\(n-1\) die entsprechende Vorgängerschreibweise.
\UebQuellen{\FormulaRefAuto{PeanoPred}[def];
\FormulaRefAuto{PeanoMinusOneNotation}[def]}
&
\textit{Induktion mit dem Schritt \(n\mapsto n+1\)}
\UebFormel{\begin{aligned}
&P(0)\dsep\forall n\in\mathbb N\,(P(n)\rightarrow P(n+1))\\
&\hspace{12mm}\vdash\forall n\in\mathbb N\,P(n),\\
&n\neq0\vdash(n-1)+1=n.
\end{aligned}}
\UebQuellen{\FormulaRefAuto{PeanoInductionAddOne}[thm];
\FormulaRefAuto{PeanoMinusOneAddOne}[thm]}\\
\addlinespace[.8ex]
\textbf{Natürliche Ordnung}
Der Band führt für die von-Neumann-Zahlen ein:
\UebFormel{\begin{aligned}
m<_{\mathbb N}n&\coloneqq m\in n,\\
m\leq_{\mathbb N}n&\coloneqq m\subseteq n.
\end{aligned}}
Im arithmetischen Aufbau werden die additiven Charakterisierungen
und konkreten Ordnungsgesetze ausgearbeitet.
\UebQuellen{\FormulaRefAuto{NaturalStrictOrderOperator}[def];
\FormulaRefAuto{NaturalOrderOperator}[def]}
&
\textit{Trichotomie und Wohlordnungsprinzip}
\UebFormel{\begin{aligned}
&m<n\lor m=n\lor n<m,\\
&B\subseteq\mathbb N\dsep B\neq\varnothing\\
&\hspace{8mm}\vdash\exists a\in B\,\forall b\in B\,(a\leq b).
\end{aligned}}
\UebQuellen{\FormulaRefAuto{PeanoOrderTrichotomy}[thm];
\FormulaRefAuto{NaturalWellOrderingPrinciple}[thm]}\\
\addlinespace[.8ex]
\textbf{Anfangsabschnitte}
\UebFormel{\NatSeg n\coloneqq\{i\in\mathbb N\mid i<n\}.}
\UebQuellen{\FormulaRefAuto{PeanoNatSeg}[def]}
&
\textit{Endliche Abschnitte fassen nicht ganz \(\mathbb N\)}
\UebFormel{\begin{aligned}
&i\in\mathbb N\vdash
i\in\NatSeg n\leftrightarrow i<n,\\
&\neg\exists F\colon\mathbb N\inj\NatSeg n.
\end{aligned}}
\UebQuellen{\FormulaRefAuto{PeanoNatSegMembership}[thm];
\FormulaRefAuto{PeanoNatSegNoNatInjection}[thm]}\\
\addlinespace[.8ex]
\textbf{Multiplikation durch wiederholte Addition}
\UebFormel{a\cdot b\coloneqq\RecFun{0}{\AddStep{a}}(b).}
Dabei ist \(\AddStep{a}\) der Additionsschritt mit festem Summanden \(a\).
\UebQuellen{\FormulaRefAuto{Multiplikation auf \(\mathbb{N}\)}[def]}
&
\textit{Multiplikative und distributive Gesetze}
\UebFormel{\begin{aligned}
&a\cdot b=b\cdot a,\\
&(a\cdot b)\cdot c=a\cdot(b\cdot c),\\
&a\cdot(b+c)=a\cdot b+a\cdot c.
\end{aligned}}
\UebQuellen{\FormulaRefAuto{PeanoMulCommutative}[thm];
\FormulaRefAuto{PeanoMulAssociative}[thm];
\FormulaRefAuto{PeanoMulLeftDistributive}[thm]}\\
\addlinespace[.8ex]
\textbf{Endliche Summen und Produkte}
\(R_u\) und \(P_u\) akkumulieren die gelesenen Werte durch
Addition beziehungsweise Multiplikation. Definiert wird:
\UebFormel{\begin{aligned}
\sum_{i<n}u(i)&\coloneqq\pi_2(R_u(n)),\\
\prod_{i<n}u(i)&\coloneqq\pi_2(P_u(n)).
\end{aligned}}
\UebQuellen{\FormulaRefAuto{Summenzeichen}[def];
\FormulaRefAuto{Produktzeichen}[def]}
&
\textit{Rekursive Auswertung der Summe}
\UebFormel{\begin{aligned}
&\sum_{i<0}u(i)=0,\\
&\sum_{i<n+1}u(i)=\left(\sum_{i<n}u(i)\right)+u(n).
\end{aligned}}
\UebQuellen{\FormulaRefAuto{u\colon\mathbb{N}\to\mathbb{N}\vdash \sum_{i<0}u(i)=0}[thm];
\FormulaRefAuto{u\colon\mathbb{N}\to\mathbb{N}\dsep n\in\mathbb{N}\vdash \sum_{i<n+1}u(i)=\left(\sum_{i<n}u(i)\right) + u(n)}[thm]}\\
\addlinespace[.8ex]
\textbf{Potenzen und Fakultäten}
\UebFormel{\begin{aligned}
b^n&\coloneqq\prod_{i<n}\ConstFun{\mathbb N}{\mathbb N}{b}(i),\\
n!&\coloneqq\prod_{i<n}(i+1).
\end{aligned}}
\UebQuellen{\FormulaRefAuto{Potenz natürlicher Zahlen}[def];
\FormulaRefAuto{FactorialFunction}[def]}
&
\textit{Null- und Schrittregeln}
\UebFormel{\begin{aligned}
&b^0=1,\qquad b^{n+1}=b^n\cdot b,\\
&0!=1,\qquad(n+1)!=n!\cdot(n+1).
\end{aligned}}
\UebQuellen{\FormulaRefAuto{PeanoPowerZero}[thm];
\FormulaRefAuto{PeanoPowerAddOne}[thm];
\FormulaRefAuto{FactorialRecursion}[thm]}\\
\addlinespace[.8ex]
\textbf{Stellenwertsysteme}
\UebFormel{\begin{aligned}
\BaseOK{b}&\coloneqq b\in\mathbb N\land2\leq b,\\
\DigitSet{b}&\coloneqq\{d\in\mathbb N\mid d<b\},\\
\PlaceTermSeq{b}{a}(i)&\coloneqq a(i)\cdot b^i.
\end{aligned}}
Zahlzeichen werden durch endliche Summen dieser Stellenwerte
ausgewertet; Dual- und Dezimalsystem sind die Fälle \(b=2,10\).
\UebQuellen{\FormulaRefAuto{Basis eines Stellenwertsystems}[def];
\FormulaRefAuto{Ziffernmenge zur Basis \(b\)}[def];
\FormulaRefAuto{Funktionsvorschrift der Stellenwertfolge}[def]}
&
\textit{Typisierung der Stellenwertfolge}
\UebFormel{\begin{aligned}
&\BaseOK{b}\dsep a\colon\mathbb N\to\mathbb N\\
&\hspace{12mm}\vdash\PlaceTermSeq{b}{a}\colon\mathbb N\to\mathbb N.
\end{aligned}}
Die normalisierten Ziffernfolgen werden definiert; ein allgemeiner
Existenz- und Eindeutigkeitssatz für jede Basis ist hier noch nicht
als bewiesenes Hauptresultat ausgewiesen.
\UebQuellen{\FormulaRefAuto{\BaseOK{b}\dsep a\colon\mathbb{N}\to\mathbb{N}\vdash \PlaceTermSeq{b}{a}\colon\mathbb{N}\to\mathbb{N}}[thm]}\\
\end{BandUebersicht}

\chapter[Band 11: Äquivalenzrelationen]{Band 11 auf einen Blick}
\noindent\textbf{Äquivalenzrelationen}\par\medskip
\noindent\emph{Lesart.} Für Klassen- und Quotientenformeln gilt
\(\EqRel{A,\sim}\). Es werden Relationsoperatoren auf einem ausdrücklich
angegebenen Träger benutzt. Gleichmächtigkeit wird auf Mengenfamilien
als Beispiel einer Äquivalenzrelation betrachtet.
\begin{BandUebersicht}
\textbf{Äquivalenzrelation}
Für \(x,y,z\in A\) lauten die Strukturaxiome:
\UebFormel{\begin{aligned}
&x\sim x,\\
&x\sim y\vdash y\sim x,\\
&x\sim y\dsep y\sim z\vdash x\sim z.
\end{aligned}}
\UebQuellen{\FormulaRefAuto{EqRelDef}[def];
\FormulaRefAuto{EqRelReflexivity}[ax];
\FormulaRefAuto{EqRelSymmetry}[ax];
\FormulaRefAuto{EqRelTransitivity}[ax]}
&
\textit{Ein konkretes Modell: Ununterscheidbarkeit}
Die in der nächsten Zeile definierte Relation erfüllt alle drei Axiome:
\UebFormel{\EqRel{\bigcup M,\SameOccRel{M}}.}
\UebQuellen{\FormulaRefAuto{SameOccRelEqRel}[thm]}\\
\addlinespace[.8ex]
\textbf{Ununterscheidbarkeit in einer Familie}
Für \(a,b\in\bigcup M\):
\UebFormel{\begin{aligned}
&a\mathrel{\SameOccRel{M}}b\coloneqq\\
&\qquad\forall X\in M\,(a\in X\leftrightarrow b\in X).
\end{aligned}}
\UebQuellen{\FormulaRefAuto{SameOccRelDef}[def]}
&
\textit{Die Quotientenabbildung trennt Familienglieder}
Mit der später eingeführten Quotientenbildabbildung:
\UebFormel{\begin{aligned}
&X,Y\in M\dsep
\OccQuotMap{M}(X)=\OccQuotMap{M}(Y)\\
&\hspace{12mm}\vdash X=Y.
\end{aligned}}
\UebQuellen{\FormulaRefAuto{OccQuotSetSeparatesMembers}[thm]}\\
\addlinespace[.8ex]
\textbf{Gleichmächtigkeit}
\UebFormel{A\EqCard B\coloneqq\exists F\colon A\bij B.}
\UebQuellen{\FormulaRefAuto{Gleichmächtigkeit}[def]}
&
\textit{Äquivalenzgesetze und Potenzmengentransport}
\UebFormel{\begin{aligned}
&A\EqCard A,\qquad A\EqCard B\vdash B\EqCard A,\\
&A\EqCard B\dsep B\EqCard C\vdash A\EqCard C,\\
&A\EqCard B\vdash\powerset(A)\EqCard\powerset(B).
\end{aligned}}
\UebQuellen{\FormulaRefAuto{EqCardReflexive}[thm];
\FormulaRefAuto{EqCardSymmetric}[thm];
\FormulaRefAuto{EqCardTransitive}[thm];
\FormulaRefAuto{EqCardPowerSetForward}[thm]}\\
\addlinespace[.8ex]
\textbf{Kardinaler Größenvergleich}
\UebFormel{\begin{aligned}
A\CardLeq B&\coloneqq\exists F\colon A\inj B,\\
A\CardLt B&\coloneqq
(A\CardLeq B\land\neg(A\EqCard B)).
\end{aligned}}
\UebQuellen{\FormulaRefAuto{CardLeqDef}[def];
\FormulaRefAuto{CardLtDef}[def]}
&
\textit{Cantor--Schröder--Bernstein}
\UebFormel{\begin{aligned}
&A\CardLeq B,\ B\CardLeq A\vdash A\EqCard B,\\
&A\EqCard B\leftrightarrow
(A\CardLeq B\land B\CardLeq A).
\end{aligned}}
Der Vergleich ist reflexiv und transitiv; sein strikter Teil
ist irreflexiv.
\UebQuellen{\FormulaRefAuto{CantorSchroederBernstein}[thm];
\FormulaRefAuto{EqCardMutualCardLeqEquiv}[thm];
\FormulaRefAuto{CardLeqReflexive}[thm];
\FormulaRefAuto{CardLeqTransitive}[thm];
\FormulaRefAuto{CardLtIrreflexive}[thm]}
\\ \addlinespace[.8ex]

\textbf{Äquivalenzklasse}
Für \(x\in A\):
\UebFormel{\EqClass{x}{\sim}{A}\coloneqq
\{y\in A\mid x\sim y\}.}
\UebQuellen{\FormulaRefAuto{EqClassDef}[def]}
&
\textit{Repräsentanten und gleiche Klassen}
Für \(x,y\in A\):
\UebFormel{\begin{aligned}
&x\in\EqClass{x}{\sim}{A},\\
&\EqClass{x}{\sim}{A}=\EqClass{y}{\sim}{A}
\leftrightarrow x\sim y.
\end{aligned}}
\UebQuellen{\FormulaRefAuto{EqClassRepresentative}[thm];
\FormulaRefAuto{EqClassEqualIffRel}[thm]}\\
\addlinespace[.8ex]
\textbf{Quotientenmenge und Projektion}
\UebFormel{\begin{aligned}
\QuotientSet{A}{\sim}\coloneqq{}&
\{C\in\powerset(A)\mid{}\\
&\exists x\in A\,C=\EqClass{x}{\sim}{A}\},\\
\QuotProj{A}{\sim}(x)\coloneqq{}&\EqClass{x}{\sim}{A}.
\end{aligned}}
\UebQuellen{\FormulaRefAuto{QuotientSetDef}[def];
\FormulaRefAuto{QuotProjDef}[def]}
&
\textit{Surjektivität und identifizierte Punkte}
\UebFormel{\begin{aligned}
&\QuotProj{A}{\sim}\colon A\sur\QuotientSet{A}{\sim},\\
&x,y\in A\vdash\\
&\quad\QuotProj{A}{\sim}(x)=\QuotProj{A}{\sim}(y)
\leftrightarrow x\sim y.
\end{aligned}}
\UebQuellen{\FormulaRefAuto{QuotProjSurjective}[thm];
\FormulaRefAuto{QuotProjEqualIffRel}[thm]}\\
\addlinespace[.8ex]
\textbf{Quotientenbild einer Mengenfamilie}
Für \(M\subseteq\powerset(A)\):
\UebFormel{\QuotFamMap{A}{\sim}{M}\coloneqq
\PowMap{\QuotProj{A}{\sim}}\restriction_M.}
\UebQuellen{\FormulaRefAuto{QuotFamMapDef}[def]}
&
\textit{Vereinigungen werden erhalten}
Für \(X,Y,X\cup Y\in M\):
\UebFormel{\begin{aligned}
&\QuotFamMap{A}{\sim}{M}(X\cup Y)\\
&\quad=\QuotFamMap{A}{\sim}{M}(X)
\cup\QuotFamMap{A}{\sim}{M}(Y).
\end{aligned}}
\UebQuellen{\FormulaRefAuto{QuotFamMapPreservesUnion}[thm]}\\
\addlinespace[.8ex]
\textbf{Ununterscheidbarkeitsquotient}
\(\OccQuotCarrier{M}\) ist der Träger der Ununterscheidbarkeitsklassen,
\(\OccQuotFam{M}\) die Familie der Bilder der Glieder von \(M\).
\(\OccQuotMap{M}\) ist die zugehörige Quotientenbildabbildung.
&
\textit{Bijektive Reduktion der Mengenfamilie}
\UebFormel{\begin{aligned}
&\OccQuotMap{M}\colon M\bij\OccQuotFam{M},\\
&\UnionClosedFam(M)\\
&\hspace{10mm}\vdash\UnionClosedFam(\OccQuotFam{M}),\\
&\TraceMap{\OccQuotMap{M}}\colon
\OccQuotCarrier{M}\inj\powerset(M).
\end{aligned}}
\UebQuellen{\FormulaRefAuto{OccQuotMapBijective}[thm];
\FormulaRefAuto{OccQuotUnionClosed}[thm];
\FormulaRefAuto{OccQuotTraceMapInjective}[thm]}\\
\end{BandUebersicht}
\noindent\emph{Lesepfad zur Gleichmächtigkeit.}
Die \CSBReadingLink{} erklärt, wie aus Injektionen in beide Richtungen
eine Bijektion entsteht und damit Gleichmächtigkeit folgt;
die \CSBProofsLink{} dokumentieren die zugrunde liegende Konstruktion.

\chapter[Band 12: Halbordnungen]{Band 12 auf einen Blick}
\noindent\textbf{Halbordnungen}\par\medskip
\noindent\emph{Lesart.} Es gilt \(\PartOrd{A,\leq}\),
\(T\subseteq A\). Ein \emph{minimales} Element wird von keinem
verschiedenen Element unterschritten; ein \emph{Minimum} liegt unter
allen Elementen. Diese Begriffe werden ausdrücklich unterschieden.
\begin{BandUebersicht}
\textbf{Partielle Ordnung und duale Relation}
Für \(x,y,z\in A\):
\UebFormel{\begin{aligned}
&x\leq x,\\
&x\leq y\dsep y\leq z\vdash x\leq z,\\
&x\leq y\dsep y\leq x\vdash x=y,\\
&x\geq y\leftrightarrow y\leq x.
\end{aligned}}
Die ersten drei Zeilen sind Strukturaxiome, die letzte definiert die Dualität.
\UebQuellen{\FormulaRefAuto{Partielle Ordnung}[def];
\FormulaRefAuto{OrderAntisymmetry}[ax];
\FormulaRefAuto{OrderDualRelation}[def]}
&
\textit{Dualität bewahrt die partielle Ordnung}
\UebFormel{\PartOrd{A,\leq}\vdash\PartOrd{A,\geq}.}
\UebQuellen{\FormulaRefAuto{OrderDualPartialOrder}[thm]}\\
\addlinespace[.8ex]
\textbf{Minimales Element und Minimum}
\UebFormel{\begin{aligned}
&\operatorname{MinEl}_{A,\leq}(m,T)\coloneqq\\
&\quad m\in T\land\forall x\in T\,(x\leq m\rightarrow x=m).
\end{aligned}}
Existiert ein kleinstes Element, wird definiert:
\UebFormel{\min_{A,\leq}(T)\coloneqq
\iota m\in T\,\forall x\in T\,(m\leq x).}
\UebQuellen{\FormulaRefAuto{OrderMinimalElement}[def];
\FormulaRefAuto{Minimum}[def]}
&
\textit{Eindeutigkeit und Minimalität des Minimums}
\UebFormel{\begin{aligned}
&\exists m\in T\,\forall x\in T\,(m\leq x)\\
&\quad\vdash\exists!m\in T\,\forall x\in T\,(m\leq x).
\end{aligned}}
Unter dieser Existenzvoraussetzung gilt außerdem
\UebFormel{\operatorname{MinEl}_{A,\leq}(\min_{A,\leq}(T),T).}
\UebQuellen{\FormulaRefAuto{MinimumUnique}[thm];
\FormulaRefAuto{MinimumIsMinimalElement}[thm]}\\
\addlinespace[.8ex]
\textbf{Maximales Element und Maximum}
\UebFormel{\begin{aligned}
&\operatorname{MaxEl}_{A,\leq}(m,T)\coloneqq\\
&\quad m\in T\land\forall x\in T\,(m\leq x\rightarrow x=m).
\end{aligned}}
Existiert ein größtes Element, wird definiert:
\UebFormel{\max_{A,\leq}(T)\coloneqq
\iota m\in T\,\forall x\in T\,(x\leq m).}
\UebQuellen{\FormulaRefAuto{OrderMaximalElement}[def];
\FormulaRefAuto{Maximum}[def]}
&
\textit{Eindeutigkeit und Maximalität des Maximums}
\UebFormel{\begin{aligned}
&\exists m\in T\,\forall x\in T\,(x\leq m)\\
&\quad\vdash\exists!m\in T\,\forall x\in T\,(x\leq m).
\end{aligned}}
Unter dieser Existenzvoraussetzung gilt
\UebFormel{\operatorname{MaxEl}_{A,\leq}(\max_{A,\leq}(T),T).}
\UebQuellen{\FormulaRefAuto{MaximumUnique}[thm];
\FormulaRefAuto{MaximumIsMaximalElement}[thm]}\\
\addlinespace[.8ex]
\textbf{Restriktion einer Ordnung}
Für \(x,y\in T\):
\UebFormel{x\leq_T y\leftrightarrow x\leq y.}
\UebQuellen{\FormulaRefAuto{OrderRestrictionOperator}[def]}
&
\textit{Die Teilmenge erbt eine partielle Ordnung}
\UebFormel{\begin{aligned}
&T\subseteq A\dsep\PartOrd{A,\leq}\\
&\hspace{12mm}\vdash\PartOrd{T,\leq_T}.
\end{aligned}}
\UebQuellen{\FormulaRefAuto{OrderRestrictionPartialOrder}[thm]}\\
\addlinespace[.8ex]
\textbf{Hauptfilter}
\UebFormel{\PrincipalFilter{A,\leq}{j}\coloneqq
\{x\in A\mid j\leq x\}.}
\UebQuellen{\FormulaRefAuto{PrincipalFilter}[def]}
&
\textit{Typisierung}
\UebFormel{\PrincipalFilter{A,\leq}{j}\subseteq A.}
Der Hauptfilter sammelt genau die Elemente oberhalb von \(j\).
\UebQuellen{\FormulaRefAuto{PrincipalFilterSubsetCarrier}[thm]}\\
\end{BandUebersicht}

\chapter[Band 13: Schranken, Infima und Suprema]{Band 13 auf einen Blick}
\noindent\textbf{Schranken, Infima und Suprema}\par\medskip
\noindent\emph{Lesart.} Es gilt \(\PartOrd{A,\leq}\),
\(T\subseteq A\). In dieser Übersicht kürzen wir die festen
Trägerindizes ab: \(L(T)=\LB_{A,\leq}(T)\),
\(U(T)=\UB_{A,\leq}(T)\), \(\inf T=\inf_{A,\leq}(T)\),
\(\sup T=\sup_{A,\leq}(T)\). Infimum und Supremum werden nur
unter ihrer jeweiligen Existenzvoraussetzung verwendet.
\begin{BandUebersicht}
\textbf{Untere und obere Schranken}
\UebFormel{\begin{aligned}
L(T)&\coloneqq\{a\in A\mid\forall x\in T\,(a\leq x)\},\\
U(T)&\coloneqq\{a\in A\mid\forall x\in T\,(x\leq a)\}.
\end{aligned}}
\UebQuellen{\FormulaRefAuto{LowerBoundSet}[def];
\FormulaRefAuto{UpperBoundSet}[def]}
&
\textit{Elementkriterien und Typisierung}
\UebFormel{\begin{aligned}
&u\in L(T)\leftrightarrow u\in A\land\forall x\in T\,(u\leq x),\\
&u\in U(T)\leftrightarrow u\in A\land\forall x\in T\,(x\leq u),\\
&L(T)\subseteq A,\qquad U(T)\subseteq A.
\end{aligned}}
\UebQuellen{\FormulaRefAuto{LowerBoundSetCharacterization}[thm];
\FormulaRefAuto{UpperBoundSetCharacterization}[thm];
\FormulaRefAuto{LowerBoundSetSubsetCarrier}[thm];
\FormulaRefAuto{UpperBoundSetSubsetCarrier}[thm]}\\
\addlinespace[.8ex]
\textbf{Schranken eines Elementpaares}
Die allgemeinen Schrankenmengen werden auf \(T=\{x,y\}\)
mit \(x,y\in A\) angewandt.
&
\textit{Zwei Ungleichungen genügen}
\UebFormel{\begin{aligned}
&u\in U(\{x,y\})\\
&\quad\leftrightarrow u\in A\land x\leq u\land y\leq u,\\
&d\in L(\{x,y\})\\
&\quad\leftrightarrow d\in A\land d\leq x\land d\leq y.
\end{aligned}}
\UebQuellen{\FormulaRefAuto{PairUpperBoundSetCharacterization}[thm];
\FormulaRefAuto{PairLowerBoundSetCharacterization}[thm]}\\
\addlinespace[.8ex]
\textbf{Supremum}
Existiert eine kleinste obere Schranke, setzt man
\UebFormel{\sup T\coloneqq\min_{A,\leq}(U(T)).}
Die ausdrückliche Voraussetzung ist
\UebFormel{\exists s\in U(T)\,\forall u\in U(T)\,(s\leq u).}
\UebQuellen{\FormulaRefAuto{Supremum}[def]}
&
\textit{Kleinste obere Schranke}
Unter der linken Existenzvoraussetzung:
\UebFormel{\begin{aligned}
&\sup T\in A,\qquad\sup T\in U(T),\\
&x\in T\vdash x\leq\sup T,\\
&u\in U(T)\vdash\sup T\leq u.
\end{aligned}}
\UebQuellen{\FormulaRefAuto{SupremumInCarrier}[thm];
\FormulaRefAuto{SupremumBelongsUpperBounds}[thm];
\FormulaRefAuto{SupremumDominatesMember}[thm];
\FormulaRefAuto{SupremumLeastUpperBound}[thm]}\\
\addlinespace[.8ex]
\textbf{Infimum}
Existiert eine größte untere Schranke, setzt man
\UebFormel{\inf T\coloneqq\max_{A,\leq}(L(T)).}
Die ausdrückliche Voraussetzung ist
\UebFormel{\exists i\in L(T)\,\forall d\in L(T)\,(d\leq i).}
\UebQuellen{\FormulaRefAuto{Infimum}[def]}
&
\textit{Größte untere Schranke}
Unter der linken Existenzvoraussetzung:
\UebFormel{\begin{aligned}
&\inf T\in A,\qquad\inf T\in L(T),\\
&x\in T\vdash\inf T\leq x,\\
&d\in L(T)\vdash d\leq\inf T.
\end{aligned}}
\UebQuellen{\FormulaRefAuto{InfimumInCarrier}[thm];
\FormulaRefAuto{InfimumBelongsLowerBounds}[thm];
\FormulaRefAuto{InfimumBelowMember}[thm];
\FormulaRefAuto{InfimumGreatestLowerBound}[thm]}\\
\addlinespace[.8ex]
\textbf{Nachweis eines Kandidaten}
Um einen vorgegebenen Wert als Supremum oder Infimum zu
identifizieren, prüft man Schranken- und Extremaleigenschaft.
&
\textit{Eindeutige Identifikation}
\UebFormel{\begin{aligned}
&u\in U(T)\dsep\forall a\in U(T)\,(u\leq a)\\
&\hspace{12mm}\vdash u=\sup T,\\
&u\in L(T)\dsep\forall a\in L(T)\,(a\leq u)\\
&\hspace{12mm}\vdash u=\inf T.
\end{aligned}}
\UebQuellen{\FormulaRefAuto{SupremumCandidateEqualsSupremum}[thm];
\FormulaRefAuto{InfimumCandidateEqualsInfimum}[thm]}\\
\addlinespace[.8ex]
\textbf{Duale Ordnung}
\UebFormel{x\geq y\leftrightarrow y\leq x\quad(x,y\in A).}
Beim Umkehren der Ordnung wechseln untere und obere Schranken ihre Rolle.
\UebQuellen{\FormulaRefAuto{OrderDualRelation}[def]}
&
\textit{Dualität der Extremalbegriffe}
\UebFormel{\begin{aligned}
&\LB_{A,\geq}(T)=\UB_{A,\leq}(T),\\
&\UB_{A,\geq}(T)=\LB_{A,\leq}(T),\\
&\sup_{A,\leq}(T)=\inf_{A,\geq}(T),\\
&\inf_{A,\leq}(T)=\sup_{A,\geq}(T).
\end{aligned}}
Die letzten beiden Gleichungen setzen die Existenz des links
stehenden Extremalwertes voraus.
\UebQuellen{\FormulaRefAuto{DualLowerBoundsAreUpperBounds}[thm];
\FormulaRefAuto{DualUpperBoundsAreLowerBounds}[thm];
\FormulaRefAuto{DualSupremumEqualsInfimum}[thm];
\FormulaRefAuto{DualInfimumEqualsSupremum}[thm]}\\
\end{BandUebersicht}

\chapter[Band 14: Paarinfima und Paarsuprema]{Band 14 auf einen Blick}
\noindent\textbf{Paarinfima und Paarsuprema}\par\medskip
\noindent\emph{Lesart.} Es gilt \(\PartOrd{A,\leq}\).
Zur knappen Darstellung schreiben wir in dieser Übersicht
\(s(x,y)=\sup_{A,\leq}(\{x,y\})\) und
\(i(x,y)=\inf_{A,\leq}(\{x,y\})\), jeweils nur wenn der
betreffende Wert existiert. Eine partielle Ordnung garantiert
seine Existenz im Allgemeinen noch nicht.
\begin{BandUebersicht}
\textbf{Existenz eines Paarsupremums}
\(\operatorname{SupEx}_{A,\leq}(x,y)\) bedeutet:
\(x,y\in A\), und die oberen Schranken von \(\{x,y\}\)
besitzen ein kleinstes Element.
\UebFormel{\begin{aligned}
&\exists s\in\UB_{A,\leq}(\{x,y\})\\
&\quad\forall u\in\UB_{A,\leq}(\{x,y\})\,(s\leq u).
\end{aligned}}
\UebQuellen{\FormulaRefAuto{PairSupremumExistence}[def]}
&
\textit{Universelle Eigenschaft}
Unter dieser Existenzvoraussetzung:
\UebFormel{\begin{aligned}
&x\leq s(x,y),\qquad y\leq s(x,y),\\
&u\in A\dsep x\leq u\dsep y\leq u\vdash s(x,y)\leq u.
\end{aligned}}
\UebQuellen{\FormulaRefAuto{PairSupremumFirstOperand}[thm];
\FormulaRefAuto{PairSupremumSecondOperand}[thm];
\FormulaRefAuto{PairSupremumUniversal}[thm]}\\
\addlinespace[.8ex]
\textbf{Existenz eines Paarinfimums}
\(\operatorname{InfEx}_{A,\leq}(x,y)\) bedeutet:
\(x,y\in A\), und die unteren Schranken von \(\{x,y\}\)
besitzen ein größtes Element.
\UebFormel{\begin{aligned}
&\exists i\in\LB_{A,\leq}(\{x,y\})\\
&\quad\forall d\in\LB_{A,\leq}(\{x,y\})\,(d\leq i).
\end{aligned}}
\UebQuellen{\FormulaRefAuto{PairInfimumExistence}[def]}
&
\textit{Duale universelle Eigenschaft}
Unter dieser Existenzvoraussetzung:
\UebFormel{\begin{aligned}
&i(x,y)\leq x,\qquad i(x,y)\leq y,\\
&d\in A\dsep d\leq x\dsep d\leq y\vdash d\leq i(x,y).
\end{aligned}}
\UebQuellen{\FormulaRefAuto{PairInfimumFirstOperand}[thm];
\FormulaRefAuto{PairInfimumSecondOperand}[thm];
\FormulaRefAuto{PairInfimumUniversal}[thm]}\\
\addlinespace[.8ex]
\textbf{Wiederholte und vertauschte Operanden}
Die Paarmengen \(\{x,x\}\) und \(\{x,y\}=\{y,x\}\)
verbinden Mengengesetze mit den Extremaldefinitionen.
&
\textit{Idempotenz und Kommutativität}
Für \(x\in A\) existieren die wiederholten Paarwerte:
\UebFormel{s(x,x)=x,\qquad i(x,x)=x.}
Für existierende Werte gilt
\UebFormel{s(x,y)=s(y,x),\qquad i(x,y)=i(y,x).}
\UebQuellen{\FormulaRefAuto{PairSupremumIdempotent}[thm];
\FormulaRefAuto{PairInfimumIdempotent}[thm];
\FormulaRefAuto{PairSupremumCommutative}[thm];
\FormulaRefAuto{PairInfimumCommutative}[thm]}\\
\addlinespace[.8ex]
\textbf{Geschachtelte Paarwerte}
Vorausgesetzt werden für die Supremumsform alle vier Existenzen
\(s(x,y)\), \(s(y,z)\), \(s(s(x,y),z)\), \(s(x,s(y,z))\).
Für die Infimumsform werden entsprechend alle vier Infima vorausgesetzt.
&
\textit{Assoziativität unter den Existenzbedingungen}
\UebFormel{\begin{aligned}
&s(s(x,y),z)=s(x,s(y,z)),\\
&i(i(x,y),z)=i(x,i(y,z)).
\end{aligned}}
\UebQuellen{\FormulaRefAuto{PairSupremumAssociative}[thm];
\FormulaRefAuto{PairInfimumAssociative}[thm]}\\
\addlinespace[.8ex]
\textbf{Inklusionsordnung auf einer Mengenfamilie}
Sei \(M\subseteq\powerset(U)\). Die Relation ist die
gewöhnliche Teilmengenrelation; für \(X,Y\in M\) sind
\(X\cap Y\) und \(X\cup Y\) zunächst nur Teilmengen von \(U\).
&
\textit{Schnitt und Vereinigung als Paarwerte}
\UebFormel{\begin{aligned}
&\PartOrd{M,\subseteq},\\
&X\cap Y\in M\vdash\inf_{M,\subseteq}(\{X,Y\})=X\cap Y,\\
&X\cup Y\in M\vdash\sup_{M,\subseteq}(\{X,Y\})=X\cup Y.
\end{aligned}}
\UebQuellen{\FormulaRefAuto{InclusionOrderOnFamily}[thm];
\FormulaRefAuto{InclusionFamilyPairInfimum}[thm];
\FormulaRefAuto{InclusionFamilyPairSupremum}[thm]}\\
\addlinespace[.8ex]
\textbf{Potenzmenge als abgeschlossener Träger}
Für \(M=\powerset(U)\) sind Schnitt und Vereinigung stets
wieder Elemente des Trägers.
&
\textit{Paarwerte in der Potenzmenge}
Für \(X,Y\subseteq U\):
\UebFormel{\begin{aligned}
&\inf_{\powerset(U),\subseteq}(\{X,Y\})=X\cap Y,\\
&\sup_{\powerset(U),\subseteq}(\{X,Y\})=X\cup Y.
\end{aligned}}
\UebQuellen{\FormulaRefAuto{PowerSetPairInfimumIntersection}[thm];
\FormulaRefAuto{PowerSetPairSupremumUnion}[thm]}\\
\end{BandUebersicht}

\chapter[Band 15: Totale Ordnungen]{Band 15 auf einen Blick}
\noindent\textbf{Totale Ordnungen}\par\medskip
\noindent\emph{Lesart.} Es gilt \(\TotOrd{A,\leq}\),
\(x,y,z\in A\). Bei den binären Operationen \(\min\) und \(\max\)
werden in dieser Übersicht die festen Indizes \({}_{A,\leq}\)
weggelassen.
\begin{BandUebersicht}
\textbf{Totale Ordnung}
Die partielle Ordnung wird um das Vergleichbarkeitsaxiom ergänzt:
\UebFormel{x,y\in A\vdash x\leq y\lor y\leq x.}
\UebQuellen{\FormulaRefAuto{Totale Ordnung}[def];
\FormulaRefAuto{OrderComparability}[ax]}
&
\textit{Minimale und maximale Elemente sind Extremwerte}
Für \(T\subseteq A\):
\UebFormel{\begin{aligned}
&\operatorname{MinEl}_{A,\leq}(m,T)\vdash m=\min_{A,\leq}(T),\\
&\operatorname{MaxEl}_{A,\leq}(m,T)\vdash m=\max_{A,\leq}(T).
\end{aligned}}
\UebQuellen{\FormulaRefAuto{TotalOrderMinimalElementIsMinimum}[thm];
\FormulaRefAuto{TotalOrderMaximalElementIsMaximum}[thm]}\\
\addlinespace[.8ex]
\textbf{Binäre Minimumsoperation}
\UebFormel{\min(x,y)\coloneqq\min_{A,\leq}(\{x,y\}).}
\UebQuellen{\FormulaRefAuto{TotalOrderBinaryMinimum}[def]}
&
\textit{Existenz, Funktionstyp und Infimum}
\UebFormel{\begin{aligned}
&\exists m\in\{x,y\}\,\forall u\in\{x,y\}\,(m\leq u),\\
&\min\colon A\times A\to A,\\
&\min(x,y)=\inf_{A,\leq}(\{x,y\}).
\end{aligned}}
\UebQuellen{\FormulaRefAuto{TotalOrderPairMinimumExistence}[thm];
\FormulaRefAuto{TotalOrderBinaryMinimumFunctionType}[thm];
\FormulaRefAuto{TotalOrderBinaryMinimumAsInfimum}[thm]}\\
\addlinespace[.8ex]
\textbf{Binäre Maximumsoperation}
\UebFormel{\max(x,y)\coloneqq\max_{A,\leq}(\{x,y\}).}
\UebQuellen{\FormulaRefAuto{TotalOrderBinaryMaximum}[def]}
&
\textit{Existenz, Funktionstyp und Supremum}
\UebFormel{\begin{aligned}
&\exists M\in\{x,y\}\,\forall u\in\{x,y\}\,(u\leq M),\\
&\max\colon A\times A\to A,\\
&\max(x,y)=\sup_{A,\leq}(\{x,y\}).
\end{aligned}}
\UebQuellen{\FormulaRefAuto{TotalOrderPairMaximumExistence}[thm];
\FormulaRefAuto{TotalOrderBinaryMaximumFunctionType}[thm];
\FormulaRefAuto{TotalOrderBinaryMaximumAsSupremum}[thm]}\\
\addlinespace[.8ex]
\textbf{Minimum und Maximum als algebraische Operationen}
Die auf allen Paaren definierten Operationen können wiederholt
und in verschiedenen Klammerungen angewandt werden.
&
\textit{Idempotenz, Kommutativität und Assoziativität}
Für \(\star=\min\) und ebenso \(\star=\max\):
\UebFormel{\begin{aligned}
&\star(x,x)=x,\qquad\star(x,y)=\star(y,x),\\
&\star(\star(x,y),z)=\star(x,\star(y,z)).
\end{aligned}}
\UebQuellen{\FormulaRefAuto{TotalOrderBinaryMinimumIdempotent}[thm];
\FormulaRefAuto{TotalOrderBinaryMaximumIdempotent}[thm];
\FormulaRefAuto{TotalOrderBinaryMinimumCommutative}[thm];
\FormulaRefAuto{TotalOrderBinaryMaximumCommutative}[thm];
\FormulaRefAuto{TotalOrderBinaryMinimumAssociative}[thm];
\FormulaRefAuto{TotalOrderBinaryMaximumAssociative}[thm]}\\
\addlinespace[.8ex]
\textbf{Natürliche Ordnung als Beispiel}
Die in Band 10 eingeführte natürliche Ordnung wird mit der
allgemeinen Definition der totalen Ordnung verbunden.
\UebQuellen{\FormulaRefAuto{NaturalOrderOperator}[def]}
&
\textit{Die natürlichen Zahlen sind total geordnet}
\UebFormel{\TotOrd{\mathbb N,\leq_{\mathbb N}}.}
Die einzelnen Ordnungsgesetze aus Band 10 erfüllen gemeinsam
die verlangten Strukturaxiome.
\UebQuellen{\FormulaRefAuto{NaturalNumbersTotalOrder}[thm]}\\
\end{BandUebersicht}

\chapter[Band 16: Wohlordnungen und Auswahlaxiom]{Band 16 auf einen Blick}
\noindent\textbf{Wohlordnungen und Auswahlaxiom}\par\medskip
\noindent\emph{Lesart.} Eine Wohlordnung ist eine totale Ordnung,
in der jede nichtleere Teilmenge ein kleinstes Element hat.
\emph{Wichtig für den Stand dieses Bandes:} Zermelos Konstruktionsprinzip
ist im Original als Axiom angesetzt. Die Richtung vom Auswahlaxiom zum
Wohlordnungssatz wird unter Verwendung dieses zusätzlichen Prinzips
hergeleitet; die Übersicht gibt es nicht als dort bewiesenen Rekursionssatz aus.
\begin{BandUebersicht}
\textbf{Wohlordnung}
Zusätzlich zu \(\TotOrd{A,\leq}\) wird gefordert:
\UebFormel{\begin{aligned}
&T\subseteq A\dsep T\neq\varnothing\\
&\hspace{10mm}\vdash\exists m\in T\,\forall x\in T\,(m\leq x).
\end{aligned}}
\UebQuellen{\FormulaRefAuto{Wohlordnung}[def];
\FormulaRefAuto{WellOrderPrinciple}[ax]}
&
\textit{Natürliches Beispiel}
\UebFormel{\WellOrd{\mathbb N,\leq_{\mathbb N}}.}
Der Satz verbindet die totale Ordnung aus Band 15 mit dem
Wohlordnungsprinzip aus Band 10.
\UebQuellen{\FormulaRefAuto{NaturalNumbersWellOrder}[thm]}\\
\addlinespace[.8ex]
\textbf{Wohlordenbarkeit und Wohlordnungssatz}
\UebFormel{\begin{aligned}
&\WellOrdable{A}\coloneqq\\
&\qquad\exists\preccurlyeq\;\WellOrd{A,\preccurlyeq},\\
&\WOP\coloneqq\forall A\,\WellOrdable{A}.
\end{aligned}}
\UebQuellen{\FormulaRefAuto{WellOrdable}[def];
\FormulaRefAuto{WOP}[def]}
&
\textit{Wohlordnungssatz liefert Auswahl}
\UebFormel{\WOP\vdash\ACsur.}
Die Richtung wird über die kleinsten Elemente der Fasern einer
Surjektion bewiesen.
\UebQuellen{\FormulaRefAuto{WOPToACsur}[thm]}\\
\addlinespace[.8ex]
\textbf{Auswahlrelation auf nichtleeren Teilmengen}
\UebFormel{\begin{aligned}
&\ChoiceRel{A}\coloneqq\\
&\qquad\{(X,a)\in\PowNE{A}\times A\\
&\hspace{20mm}\mid a\in X\}.
\end{aligned}}
\UebQuellen{\FormulaRefAuto{ChoiceRelation}[def]}
&
\textit{Die Auswahlrelation ist total}
\UebFormel{\begin{aligned}
&\TotRel{\ChoiceRel{A},\PowNE{A},A},\\
&(X,a)\in\ChoiceRel{A}\vdash a\in X.
\end{aligned}}
\UebQuellen{\FormulaRefAuto{ChoiceRelTotal}[thm];
\FormulaRefAuto{ChoiceRelElim}[thm]}\\
\addlinespace[.8ex]
\textbf{Auswahlfunktion auf der nichtleeren Potenzmenge}
\UebFormel{\begin{aligned}
&\ChoicePow{C}{A}\coloneqq\\
&\quad C\colon\PowNE{A}\to A\\
&\quad\land\forall X\in\PowNE{A}\,(C(X)\in X).
\end{aligned}}
\UebQuellen{\FormulaRefAuto{ChoicePow}[def]}
&
\textit{AC liefert eine solche Auswahlfunktion}
\UebFormel{\ACsur\vdash\exists C\,\ChoicePow{C}{A}.}
\UebQuellen{\FormulaRefAuto{ACsurChoicePow}[thm]}\\
\addlinespace[.8ex]
\textbf{Zermelos Konstruktionsprinzip -- Axiom im Original}
\UebFormel{\begin{aligned}
&\ChoicePow{C}{A}\\
&\hspace{10mm}\vdash\exists\preccurlyeq\;\WellOrd{A,\preccurlyeq}.
\end{aligned}}
Dieses Prinzip steht auf der Axiomseite des gegenwärtigen Aufbaus.
\UebQuellen{\FormulaRefAuto{ZermeloConstruction}[ax]}
&
\textit{Unter diesem Prinzip: Auswahl impliziert Wohlordnung}
\UebFormel{\begin{aligned}
&\exists C\,\ChoicePow{C}{A}\vdash\WellOrdable{A},\\
&\ACsur\vdash\WOP,\\
&\ACsur\eqvdash\WOP.
\end{aligned}}
Die letzte Zeile kombiniert diese Richtung mit dem unabhängig
von der Konstruktion bewiesenen Rückweg.
\UebQuellen{\FormulaRefAuto{ZermeloWO}[thm];
\FormulaRefAuto{ACsurToWOP}[thm];
\FormulaRefAuto{ACsurEquivWOP}[thm]}\\
\addlinespace[.8ex]
\textbf{Faserminimenge}
Für \(\WellOrd{B,\preccurlyeq}\), \(G\colon B\sur A\):
\UebFormel{\begin{aligned}
&L_{G,\preccurlyeq}\coloneqq\{b\in B\mid\forall c\in B\\
&\qquad(G(c)=G(b)\rightarrow b\preccurlyeq c)\}.
\end{aligned}}
\UebQuellen{\FormulaRefAuto{FiberMinSet}[def]}
&
\textit{Faserminima liefern eine Sektion}
\UebFormel{\begin{aligned}
&\forall X\in\Fib_G[A]\\
&\qquad\exists!b\,(b\in X\land b\in L_{G,\preccurlyeq}),\\
&\exists F\colon A\to B\;G\circ F=\Id_A.
\end{aligned}}
\UebQuellen{\FormulaRefAuto{WellOrdFiberChoiceSet}[thm];
\FormulaRefAuto{WellOrdFibersSection}[thm]}\\
\end{BandUebersicht}

\chapter[Band 17: Ganze Zahlen]{Band 17 auf einen Blick}
\noindent\textbf{Ganze Zahlen}\par\medskip
\noindent\emph{Lesart.}
In Klassenformeln sind \(a,b,c,d\in\mathbb N\), in der Ganzzahlarithmetik
\(z,w,u\in\mathbb Z\). Die Zahlzeichen \(0_{\mathbb Z}=\IntZero\) und
\(1_{\mathbb Z}=\IntOne\) bezeichnen die eingebetteten natürlichen Zahlen.
Die später ausdrücklich registrierten Axiome werden zuvor im Quotientenmodell
bewiesen. Die Übersicht unterscheidet diese beiden Rollen.

\begin{BandUebersicht}
\textbf{Differenzenpaare und ihre Äquivalenz}
\UebFormel{\begin{aligned}
D_{\mathbb Z}&\coloneqq\mathbb N\times\mathbb N,\\
(a,b)\IntPairRel(c,d)&\coloneqq a+d=b+c.
\end{aligned}}
\UebQuellen{\FormulaRefAuto{IntegerPairCarrier}[def];
\FormulaRefAuto{IntPairRelDef}[def]}
&
\textit{Äquivalenzrelation}
\UebFormel{\EqRel{D_{\mathbb Z},\IntPairRel}.}
Reflexivität, Symmetrie und Transitivität folgen aus der natürlichen
Addition und ihrer Kürzbarkeit.
\UebQuellen{\FormulaRefAuto{IntegerPairEquivalence}[thm]}
\\ \addlinespace[.8ex]

\textbf{Quotient und Ganzzahlklassen}
\UebFormel{\begin{aligned}
\mathbb Z&\coloneqq\QuotientSet{D_{\mathbb Z}}{\IntPairRel},\\
\IntClass a b&\coloneqq
\EqClass{(a,b)}{\IntPairRel}{D_{\mathbb Z}}.
\end{aligned}}
\UebQuellen{\FormulaRefAuto{IntegersAsQuotient}[def];
\FormulaRefAuto{IntegerClassDef}[def]}
&
\textit{Repräsentanten und Gleichheitskriterium}
\UebFormel{\begin{aligned}
&\IntClass a b=\IntClass c d\\
&\qquad\leftrightarrow a+d=b+c,\\
&z\in\mathbb Z\vdash
\exists a,b\in\mathbb N\;(z=\IntClass a b).
\end{aligned}}
\UebQuellen{\FormulaRefAuto{IntegerClassEquality}[thm];
\FormulaRefAuto{IntegerPairRepresentative}[thm]}
\\ \addlinespace[.8ex]

\textbf{Natürliche Einbettung und Teilmengenkopie}
\UebFormel{\begin{aligned}
&\IntEmbed\colon\mathbb N\to\mathbb Z,\qquad
\IntEmbed(n)\coloneqq\IntClass n0,\\
&\NatCopyInInt\coloneqq\IntEmbed[\mathbb N].
\end{aligned}}
\UebQuellen{\FormulaRefAuto{NaturalIntegerEmbedding}[def];
\FormulaRefAuto{NaturalCopyInIntegers}[def]}
&
\textit{Injektivität und Verträglichkeit}
Für \(m,n\in\mathbb N\):
\UebFormel{\begin{aligned}
&\IntEmbed\colon\mathbb N\inj\mathbb Z,\\
&\IntEmbed(m+n)=\IntEmbed(m)+\IntEmbed(n),\\
&\IntEmbed(mn)=\IntEmbed(m)\IntEmbed(n).
\end{aligned}}
\UebQuellen{\FormulaRefAuto{NaturalIntegerEmbeddingInjective}[thm];
\FormulaRefAuto{NaturalIntegerEmbeddingAdditive}[thm];
\FormulaRefAuto{NaturalIntegerEmbeddingMultiplicative}[thm]}
\\ \addlinespace[.8ex]

\textbf{Negation, Addition und Subtraktion}
\UebFormel{\begin{aligned}
-\IntClass a b&\coloneqq\IntClass b a,\\
\IntClass a b+\IntClass c d
&\coloneqq\IntClass{a+c}{b+d},\\
z-w&\coloneqq z+(-w).
\end{aligned}}
\UebQuellen{\FormulaRefAuto{IntegerNegationDef}[def];
\FormulaRefAuto{IntegerAdditionDef}[def];
\FormulaRefAuto{IntegerSubtractionDef}[def]}
&
\textit{Repräsentantenunabhängigkeit und additive Gesetze}
\UebFormel{\begin{aligned}
&(z+w)+u=z+(w+u),\\
&z+w=w+z,\qquad z+\IntZero=z,\\
&z+(-z)=\IntZero,\qquad-(-z)=z.
\end{aligned}}
\UebQuellen{\FormulaRefAuto{IntegerAdditionCompatible}[thm];
\FormulaRefAuto{IntAddAssociative}[thm];
\FormulaRefAuto{IntAddCommutative}[thm];
\FormulaRefAuto{IntAddZero}[thm];
\FormulaRefAuto{IntAddInverse}[thm];
\FormulaRefAuto{IntegerDoubleNegation}[thm]}
\\ \addlinespace[.8ex]

\textbf{Multiplikation}
\UebFormel{\begin{aligned}
&\IntClass a b\cdot\IntClass c d\\
&\qquad\coloneqq\IntClass{ac+bd}{ad+bc}.
\end{aligned}}
\UebQuellen{\FormulaRefAuto{IntegerMultiplicationDef}[def]}
&
\textit{Multiplikative Gesetze}
\UebFormel{\begin{aligned}
&(zw)u=z(wu),\qquad zw=wz,\\
&z\IntOne=z,\\
&z(w+u)=zw+zu.
\end{aligned}}
\UebQuellen{\FormulaRefAuto{IntegerMultiplicationCompatible}[thm];
\FormulaRefAuto{IntMulAssociative}[thm];
\FormulaRefAuto{IntMulCommutative}[thm];
\FormulaRefAuto{IntMulOne}[thm];
\FormulaRefAuto{IntMulLeftDistributive}[thm]}
\\ \addlinespace[.8ex]

\textbf{Arithmetische Axiome}
Die axiomatische Zusammenfassung fordert die eben modellierten
Abschluss-, Gruppen-, Einheits- und Distributivgesetze, insbesondere
\UebFormel{\begin{aligned}
&z+\IntZero=z,\qquad z+(-z)=\IntZero,\\
&z(w+u)=zw+zu.
\end{aligned}}
\UebQuellen{\FormulaRefAuto{IntegerArithmeticAxioms}[def];
\FormulaRefAuto{IntegerAddZeroAxiom}[ax];
\FormulaRefAuto{IntegerAddInverseAxiom}[ax];
\FormulaRefAuto{IntegerMulLeftDistributiveAxiom}[ax]}
&
\textit{Folgerungen aus dieser Axiomatik}
\UebFormel{\begin{aligned}
&z\IntZero=\IntZero,\qquad-\IntZero=\IntZero,\\
&z(-w)=-(zw),\\
&(-z)(-w)=zw.
\end{aligned}}
\UebQuellen{\FormulaRefAuto{IntegerMulZeroFromAxioms}[thm];
\FormulaRefAuto{IntegerNegativeZero}[thm];
\FormulaRefAuto{IntegerRightNegativeProduct}[thm];
\FormulaRefAuto{IntegerNegativeTimesNegative}[thm]}
\\ \addlinespace[.8ex]

\textbf{Schritte um eins}
Als Abbildungen \(\mathbb Z\to\mathbb Z\) werden eingeführt:
\UebFormel{\begin{aligned}
&(x\mapsto x+\IntOne)(z)\coloneqq z+\IntOne,\\
&(x\mapsto x-\IntOne)(z)\coloneqq z-\IntOne.
\end{aligned}}
\UebQuellen{\FormulaRefAuto{IntegerSuccessorDef}[def];
\FormulaRefAuto{IntegerPredecessorDef}[def]}
&
\textit{Inverse Schritte und Normalform}
\UebFormel{\begin{aligned}
&(z+\IntOne)-\IntOne=z,\\
&(z-\IntOne)+\IntOne=z,\\
&\exists n\in\mathbb N\;
 \bigl(z=\IntEmbed(n)\lor z=-\IntEmbed(n)\bigr).
\end{aligned}}
Insbesondere ist der Schritt nach rechts bijektiv.
\UebQuellen{\FormulaRefAuto{IntegerPredAfterSucc}[thm];
\FormulaRefAuto{IntegerSuccAfterPred}[thm];
\FormulaRefAuto{IntegerSuccessorBijective}[thm];
\FormulaRefAuto{IntegerSignedNormalForm}[thm]}
\\ \addlinespace[.8ex]

\textbf{Quotientenordnung}
\UebFormel{\begin{aligned}
&\IntClass a b\leq_{\mathbb Z}\IntClass c d
\coloneqq a+d\leq c+b,\\
&z<w\coloneqq z\leq w\land z\neq w.
\end{aligned}}
\UebQuellen{\FormulaRefAuto{IntegerOrderDef}[def];
\FormulaRefAuto{IntegerOrderNotation}[def];
\FormulaRefAuto{IntegerStrictOrderDef}[def]}
&
\textit{Totale und translationsverträgliche Ordnung}
\UebFormel{\begin{aligned}
&\TotOrd{\mathbb Z,\leq_{\mathbb Z}},\\
&z\leq w\leftrightarrow z+u\leq w+u,\\
&z<z+\IntOne.
\end{aligned}}
\UebQuellen{\FormulaRefAuto{IntegerOrderTotal}[thm];
\FormulaRefAuto{IntegerOrderTranslationModel}[thm];
\FormulaRefAuto{IntegerSuccessorPositive}[thm]}
\\ \addlinespace[.8ex]

\textbf{Positive ganze Zahlen}
\UebFormel{\PositiveIntegers\coloneqq
\{b\in\mathbb Z\mid\IntZero<b\}.}
\UebQuellen{\FormulaRefAuto{PositiveIntegersDef}[def]}
&
\textit{Kleinste positive Zahl und Faktorkürzung}
Für \(b\in\PositiveIntegers\) gilt
\UebFormel{\begin{aligned}
&\IntOne\leq b,\\
&z\leq w\leftrightarrow zb\leq wb,\\
&zb=wb\vdash z=w.
\end{aligned}}
\UebQuellen{\FormulaRefAuto{PositiveIntegerOneLowerBound}[thm];
\FormulaRefAuto{IntegerPositiveFactorOrder}[thm];
\FormulaRefAuto{IntegerPositiveFactorCancellation}[thm]}
\\ \addlinespace[.8ex]

\textbf{Zweiseitige Peano-Struktur}
Zur geordneten Arithmetik tritt das Induktionsaxiom mit den Daten
\UebFormel{\begin{aligned}
&P(\IntZero),\\
&\forall z\in\mathbb Z\;(P(z)\to P(z+\IntOne)),\\
&\forall z\in\mathbb Z\;(P(z)\to P(z-\IntOne)).
\end{aligned}}
\UebQuellen{\FormulaRefAuto{IntegerPeanoAxioms}[def];
\FormulaRefAuto{IntegerTwoSidedInductionAxiom}[ax]}
&
\textit{Modellnachweis und Induktion im Quotienten}
Aus den links genannten Daten folgt im konstruierten Modell
\UebFormel{\forall z\in\mathbb Z\;P(z).}
Der Quotient erfüllt die vollständige geordnete Ganzzahlarithmetik
einschließlich der inversen Schritte und dieser Induktion.
\UebQuellen{\FormulaRefAuto{IntegerTwoSidedInductionModel}[thm];
\FormulaRefAuto{IntQuotientModelsIntegerPeano}[thm]}
\\
\end{BandUebersicht}

\chapter[Band 18: Rationale Zahlen]{Band 18 auf einen Blick}
\noindent\textbf{Rationale Zahlen}\par\medskip
\noindent\emph{Lesart.}
Bruchpaare besitzen Zähler \(a,c\in\mathbb Z\) und positive Nenner
\(b,d\in\PositiveIntegers\). In rationalen Formeln sind \(p,q,r\in\mathbb Q\).
Die Zeichen \(\RatZero,\RatOne\) unterscheiden rationale Konstanten von
den ganzzahligen. Die Axiomatik des erzeugten geordneten Körpers folgt
im Band auf ihren Modellnachweis.

\begin{BandUebersicht}
\textbf{Bruchpaare und Äquivalenz}
\UebFormel{\begin{aligned}
\RatPairCarrier&\coloneqq
\mathbb Z\times\PositiveIntegers,\\
(a,b)\RatPairRel(c,d)&\coloneqq ad=cb.
\end{aligned}}
\UebQuellen{\FormulaRefAuto{RationalPairCarrier}[def];
\FormulaRefAuto{RationalPairRelDef}[def]}
&
\textit{Äquivalenzrelation auf dem Bruchträger}
\UebFormel{\EqRel{\RatPairCarrier,\RatPairRel}.}
Die Positivität der Nenner erlaubt bei der Transitivität
die Kürzung eines gemeinsamen Faktors.
\UebQuellen{\FormulaRefAuto{RationalPairEquivalence}[thm];
\FormulaRefAuto{RationalPairRelTransitive}[thm]}
\\ \addlinespace[.8ex]

\textbf{Rationale Zahlen und Bruchklassen}
\UebFormel{\begin{aligned}
\mathbb Q&\coloneqq
\QuotientSet{\RatPairCarrier}{\RatPairRel},\\
\RatClass a b&\coloneqq
\EqClass{(a,b)}{\RatPairRel}{\RatPairCarrier}.
\end{aligned}}
\UebQuellen{\FormulaRefAuto{RationalsAsQuotient}[def];
\FormulaRefAuto{RationalClassDef}[def]}
&
\textit{Gleichheit und Darstellung}
\UebFormel{\begin{aligned}
&\RatClass a b=\RatClass c d
\leftrightarrow ad=cb,\\
&q\in\mathbb Q\vdash
\exists a\in\mathbb Z\;\exists b\in\PositiveIntegers\\
&\hspace{5em}q=\RatClass a b.
\end{aligned}}
\UebQuellen{\FormulaRefAuto{RationalClassEquality}[thm];
\FormulaRefAuto{RationalPairRepresentative}[thm]}
\\ \addlinespace[.8ex]

\textbf{Ganzzahlige Einbettung und Zahlkopien}
\UebFormel{\begin{aligned}
&\RatEmbed\colon\mathbb Z\to\mathbb Q,\qquad
\RatEmbed(z)\coloneqq\RatClass z{\IntOne},\\
&\IntCopyInRat\coloneqq\RatEmbed[\mathbb Z],\\
&\NatCopyInRat\coloneqq\RatEmbed[\NatCopyInInt].
\end{aligned}}
\UebQuellen{\FormulaRefAuto{IntegerRationalEmbedding}[def];
\FormulaRefAuto{IntegerCopyInRationals}[def];
\FormulaRefAuto{NaturalCopyInRationals}[def]}
&
\textit{Einbettung des bisherigen Zahlbereichs}
Für \(z,w\in\mathbb Z\):
\UebFormel{\begin{aligned}
&\RatEmbed\colon\mathbb Z\inj\mathbb Q,\\
&\RatEmbed(z+w)=\RatEmbed(z)+\RatEmbed(w),\\
&\RatEmbed(zw)=\RatEmbed(z)\RatEmbed(w),\\
&\NatCopyInRat\subseteq\IntCopyInRat\subseteq\mathbb Q.
\end{aligned}}
\UebQuellen{\FormulaRefAuto{IntegerRationalEmbeddingInjective}[thm];
\FormulaRefAuto{IntegerRationalEmbeddingAdditive}[thm];
\FormulaRefAuto{IntegerRationalEmbeddingMultiplicative}[thm];
\FormulaRefAuto{EmbeddedNaturalIntegerRationalTower}[thm]}
\\ \addlinespace[.8ex]

\textbf{Negation und Addition durch Quotientenabstieg}
Die Definitionen wählen jeweils den eindeutig bestimmten rationalen
Wert aus den Repräsentanten:
\UebFormel{\begin{aligned}
&-q\coloneqq\iota u\in\mathbb Q:\\
&\quad\exists a\in\mathbb Z\;\exists b\in\PositiveIntegers\\
&\quad(q=\RatClass a b\land u=\RatClass{-a}b).
\end{aligned}}
Die Addition verwendet entsprechend den Wert
\(\RatClass{ad+cb}{bd}\) zu \(q=\RatClass a b,\ r=\RatClass c d\).
\UebQuellen{\FormulaRefAuto{RationalNegationDef}[def];
\FormulaRefAuto{RationalAdditionDef}[def]}
&
\textit{Klassenrechnung und additive Gesetze}
\UebFormel{\begin{aligned}
&-\RatClass a b=\RatClass{-a}b,\\
&\RatClass a b+\RatClass c d=\RatClass{ad+cb}{bd},\\
&(p+q)+r=p+(q+r),\\
&q+r=r+q,\qquad q+(-q)=\RatZero.
\end{aligned}}
\UebQuellen{\FormulaRefAuto{RationalNegationClassValue}[thm];
\FormulaRefAuto{RationalAdditionClassValue}[thm];
\FormulaRefAuto{RationalAddAssociative}[thm];
\FormulaRefAuto{RationalAddCommutative}[thm];
\FormulaRefAuto{RationalAddInverse}[thm]}
\\ \addlinespace[.8ex]

\textbf{Multiplikation durch Quotientenabstieg}
\UebFormel{\begin{aligned}
q r\coloneqq\iota u\in\mathbb Q\;&
\exists a,c\in\mathbb Z\\
&\exists b,d\in\PositiveIntegers:\\
&q=\RatClass a b\land r=\RatClass c d\\
&{}\land u=\RatClass{ac}{bd}.
\end{aligned}}
\UebQuellen{\FormulaRefAuto{RationalMultiplicationDef}[def]}
&
\textit{Multiplikative Gesetze}
\UebFormel{\begin{aligned}
&\RatClass a b\cdot\RatClass c d=\RatClass{ac}{bd},\\
&(pq)r=p(qr),\qquad qr=rq,\\
&q\RatOne=q,\qquad p(q+r)=pq+pr.
\end{aligned}}
\UebQuellen{\FormulaRefAuto{RationalMultiplicationClassValue}[thm];
\FormulaRefAuto{RationalMulAssociative}[thm];
\FormulaRefAuto{RationalMulCommutative}[thm];
\FormulaRefAuto{RationalMulOne}[thm];
\FormulaRefAuto{RationalDistributive}[thm]}
\\ \addlinespace[.8ex]

\textbf{Kehrwert und Division}
Für \(a\neq\IntZero\):
\UebFormel{\begin{aligned}
\RatRecipRep a b&\coloneqq
\begin{cases}
\RatClass b a,&\IntZero<a,\\
\RatClass{-b}{-a},&a<\IntZero,
\end{cases}\\
q/r&\coloneqq q\,r^{-1}\quad(r\neq\RatZero).
\end{aligned}}
Der Kehrwert \(q^{-1}\) ist der eindeutige Wert dieser Vorschrift
für einen Repräsentanten von \(q\neq\RatZero\).
\UebQuellen{\FormulaRefAuto{RationalReciprocalRepresentativeDef}[def];
\FormulaRefAuto{RationalReciprocalDef}[def];
\FormulaRefAuto{RationalDivisionDef}[def]}
&
\textit{Wohldefiniertheit und Inversengesetz}
\UebFormel{\begin{aligned}
&\RatClass a b=\RatZero\leftrightarrow a=\IntZero,\\
&q\neq\RatZero\vdash q^{-1}\in\mathbb Q,\\
&q\neq\RatZero\vdash qq^{-1}=\RatOne,\\
&r\neq\RatZero\vdash q/r\in\mathbb Q.
\end{aligned}}
\UebQuellen{\FormulaRefAuto{RationalClassZeroIffNumeratorZero}[thm];
\FormulaRefAuto{RationalReciprocalCompatible}[thm];
\FormulaRefAuto{RationalReciprocalClosure}[thm];
\FormulaRefAuto{RationalMulInverse}[thm];
\FormulaRefAuto{RationalDivisionClosure}[thm]}
\\ \addlinespace[.8ex]

\textbf{Quotientenordnung und strikter Vergleich}
\UebFormel{\begin{aligned}
&\RatOrderGraph\coloneqq\\
&\quad\{(\RatClass a b,\RatClass c d)\mid ad\leq cb\},\\
&q\leq_{\mathbb Q}r\coloneqq(q,r)\in\RatOrderGraph,\\
&q<r\coloneqq q\leq r\land q\neq r.
\end{aligned}}
Die Variablen der Mengenbeschreibung haben die eingangs angegebenen Typen.
\UebQuellen{\FormulaRefAuto{RationalOrderRelationDef}[def];
\FormulaRefAuto{RationalOrderDef}[def];
\FormulaRefAuto{RationalStrictOrderDef}[def]}
&
\textit{Ordnung und Arithmetik sind verträglich}
\UebFormel{\begin{aligned}
&\TotOrd{\mathbb Q,\leq_{\mathbb Q}},\\
&q\leq r\leftrightarrow p+q\leq p+r,\\
&\RatZero<p\land q<r\vdash pq<pr,\\
&z\leq w\leftrightarrow\RatEmbed(z)\leq\RatEmbed(w).
\end{aligned}}
In der letzten Zeile sind \(z,w\in\mathbb Z\).
\UebQuellen{\FormulaRefAuto{RationalOrderTotal}[thm];
\FormulaRefAuto{RationalOrderTranslation}[thm];
\FormulaRefAuto{RationalPositiveMulStrictOrder}[thm];
\FormulaRefAuto{IntegerRationalEmbeddingOrder}[thm]}
\\ \addlinespace[.8ex]

\textbf{Axiomatik des erzeugten geordneten Körpers}
Neben den Körper- und Ordnungsaxiomen wird ausdrücklich
die Erzeugung aus ganzen Zahlen verlangt:
\UebFormel{\begin{aligned}
q\in\mathbb Q\vdash
\exists a\in\mathbb Z\;\exists b\in\PositiveIntegers\\
q=\RatEmbed(a)\bigl(\RatEmbed(b)\bigr)^{-1}.
\end{aligned}}
\UebQuellen{\FormulaRefAuto{RationalOrderedFieldAxioms}[def];
\FormulaRefAuto{RationalFractionGenerationAxiom}[ax]}
&
\textit{Das Quotientenmodell erfüllt die Axiomatik}
\UebFormel{\begin{aligned}
&\mathbb Q=\QuotientSet{\RatPairCarrier}{\RatPairRel},\\
&\RationalOrderedField{\mathbb Q}.
\end{aligned}}
Auch die links als Axiom verwendete Brucherzeugung ist für das
konstruierte Modell ein bewiesenes Resultat.
\UebQuellen{\FormulaRefAuto{RationalFractionGenerationModel}[thm];
\FormulaRefAuto{RationalQuotientModelSatisfiesOrderedField}[thm]}
\\ \addlinespace[.8ex]

\textbf{Positive Nenner und Ganzzahlschranken}
Die axiomatische Weiterarbeit importiert
\UebFormel{\begin{aligned}
&b\in\PositiveIntegers\vdash b\in\mathbb Z
\land\IntZero<b,\\
&a\in\mathbb Z,\ b\in\PositiveIntegers\vdash\\
&\qquad\exists n\in\mathbb N\;a<\IntEmbed(n)b.
\end{aligned}}
\UebQuellen{\FormulaRefAuto{RationalPositiveDenominatorAxiom}[ax];
\FormulaRefAuto{RationalIntegerFractionBoundAxiom}[ax]}
&
\textit{Dichte und Archimedizität}
\UebFormel{\begin{aligned}
&q<r\vdash\exists s\in\mathbb Q\;(q<s<r),\\
&q\in\mathbb Q\vdash\exists n\in\mathbb N\\
&\qquad q<\RatEmbed(\IntEmbed(n)).
\end{aligned}}
Zwischen rationalen Zahlen liegt also wieder eine rationale Zahl;
die natürlichen Bilder sind dennoch nach oben unbeschränkt.
\UebQuellen{\FormulaRefAuto{RationalOrderDense}[thm];
\FormulaRefAuto{RationalArchimedean}[thm]}
\\
\end{BandUebersicht}

\chapter[Band 19: Reelle Zahlen]{Band 19 auf einen Blick}
\noindent\textbf{Reelle Zahlen}\par\medskip
\noindent\emph{Lesart.}
\(A,B,C\) sind reelle Zahlen, im Modell also Dedekindsche Schnitte;
\(p,q,r\) sind rationale Zahlen. In dieser Übersicht stehen
\(0,1,+,-,\cdot,\leq,<\) bei reellen Argumenten für die im Band konstruierten
Schnittoperationen und die Schnittordnung. Rationale Operationen innerhalb
einer Schnittbeschreibung bleiben rationale Operationen.

\begin{BandUebersicht}
\textbf{Dedekindsche Schnitte und reeller Träger}
\(\RealCut A\) wird durch die Schnittaxiome festgelegt:
\UebFormel{\begin{aligned}
&\varnothing\neq A\subsetneq\mathbb Q,\\
&p\in A,\ q\in\mathbb Q,\ q<p\vdash q\in A,\\
&p\in A\vdash\exists r\in A\;(p<r),\\
&\mathbb R\coloneqq
\{A\in\powerset(\mathbb Q)\mid\RealCut A\}.
\end{aligned}}
\UebQuellen{\FormulaRefAuto{DedekindCutPredicate}[def];
\FormulaRefAuto{DedekindCutNonemptyAxiom}[ax];
\FormulaRefAuto{DedekindCutDownwardAxiom}[ax];
\FormulaRefAuto{DedekindCutNoMaximumAxiom}[ax];
\FormulaRefAuto{RealCutCarrier}[def]}
&
\textit{Schnittkriterium und Elementcharakterisierung}
Genau die links genannten drei Bedingungen an \(A\) sind
äquivalent zur Schnitteigenschaft.
\UebFormel{A\in\mathbb R\eqvdash\RealCut A.}
Damit werden die axiomatisch beschriebenen Schnitte als Elemente
einer Menge zusammengefasst.
\UebQuellen{\FormulaRefAuto{DedekindCutCharacterization}[thm];
\FormulaRefAuto{RealCutMembership}[thm]}
\\ \addlinespace[.8ex]

\textbf{Hauptschnitt und rationale Einbettung}
\UebFormel{\begin{aligned}
\RealHat q&\coloneqq\{p\in\mathbb Q\mid p<q\},\\
\RealEmb&\colon\mathbb Q\to\mathbb R,\\
\RealEmb(q)&\coloneqq\RealHat q.
\end{aligned}}
\UebQuellen{\FormulaRefAuto{RealPrincipalCut}[def];
\FormulaRefAuto{RationalToRealEmbedding}[def]}
&
\textit{Injektive und ordnungstreue Einbettung}
\UebFormel{\begin{aligned}
&\RealCut{\RealHat q},\qquad
\RealEmb\colon\mathbb Q\inj\mathbb R,\\
&p\leq q\leftrightarrow\RealEmb(p)\leq\RealEmb(q),\\
&p<q\leftrightarrow\RealEmb(p)<\RealEmb(q).
\end{aligned}}
\UebQuellen{\FormulaRefAuto{RealPrincipalCutIsCut}[thm];
\FormulaRefAuto{RationalToRealEmbeddingInjective}[thm];
\FormulaRefAuto{RationalToRealEmbeddingOrder}[thm]}
\\ \addlinespace[.8ex]

\textbf{Schnittordnung}
\UebFormel{\begin{aligned}
A\RealLe B&\coloneqq A\subseteq B,\\
A\RealLt B&\coloneqq A\subsetneq B.
\end{aligned}}
Für die Vollständigkeit betrachtet man
\(\varnothing\neq\mathcal A\subseteq\mathbb R\) mit einer oberen
Schranke \(U\in\mathbb R\).
\UebQuellen{\FormulaRefAuto{RealCutOrder}[def];
\FormulaRefAuto{RealCutStrictOrder}[def]}
&
\textit{Totale und vollständige Ordnung}
\UebFormel{\begin{aligned}
&\TotOrd{\mathbb R,\RealLe},\\
&(\forall A\in\mathcal A\;A\leq U)\vdash\\
&\quad\bigcup\mathcal A\in\mathbb R,\qquad
\sup_{\mathbb R,\RealLe}\mathcal A=\bigcup\mathcal A.
\end{aligned}}
Eine nichtleere nach unten beschränkte reelle Menge besitzt
entsprechend ein reelles Infimum.
\UebQuellen{\FormulaRefAuto{RealCutTotalOrder}[thm];
\FormulaRefAuto{RealLeastUpperBoundProperty}[thm];
\FormulaRefAuto{RealCutSupremumUnion}[thm];
\FormulaRefAuto{RealGreatestLowerBoundProperty}[thm]}
\\ \addlinespace[.8ex]

\textbf{Kanonische Zahlkopien}
\UebFormel{\begin{aligned}
\RatCopyInReal&\coloneqq\RealEmb[\mathbb Q],\\
\IntCopyInReal&\coloneqq\RealEmb[\IntCopyInRat],\\
\NatCopyInReal&\coloneqq\RealEmb[\NatCopyInRat].
\end{aligned}}
\UebQuellen{\FormulaRefAuto{RationalCopyInReals}[def];
\FormulaRefAuto{IntegerCopyInReals}[def];
\FormulaRefAuto{NaturalCopyInReals}[def]}
&
\textit{Geschachtelter Zahlbereichsturm}
\UebFormel{\begin{aligned}
&\NatCopyInReal\subseteq\IntCopyInReal,\\
&\IntCopyInReal\subseteq\RatCopyInReal
\subseteq\mathbb R,\\
&\RealEmb\colon\mathbb Q\bij\RatCopyInReal.
\end{aligned}}
Die Identifikation erfolgt über die Einbettungen und ihre
Einschränkungen.
\UebQuellen{\FormulaRefAuto{EmbeddedNumberTowerInReals}[thm];
\FormulaRefAuto{RationalCopyRealBijection}[thm];
\FormulaRefAuto{IntegerCopyRealBijection}[thm];
\FormulaRefAuto{NaturalCopyRealBijection}[thm]}
\\ \addlinespace[.8ex]

\textbf{Addition und additive Null}
\UebFormel{\begin{aligned}
A+B\coloneqq\{q\in\mathbb Q\mid{}&
\exists a\in A\;\exists b\in B\\
&q<a+b\},\\
0&\coloneqq\RealHat{\RatZero}.
\end{aligned}}
\UebQuellen{\FormulaRefAuto{RealCutAddition}[def];
\FormulaRefAuto{RealCutZero}[def]}
&
\textit{Abschluss und additive Gesetze}
\UebFormel{\begin{aligned}
&A+B\in\mathbb R,\qquad0\in\mathbb R,\\
&(A+B)+C=A+(B+C),\\
&A+B=B+A,\qquad A+0=A.
\end{aligned}}
\UebQuellen{\FormulaRefAuto{RealCutAdditiveClosure}[thm];
\FormulaRefAuto{RealCutAdditiveLaws}[thm]}
\\ \addlinespace[.8ex]

\textbf{Negation eines Schnittes}
\UebFormel{\begin{aligned}
-A\coloneqq\{q\in\mathbb Q\mid{}&
\exists r\in\mathbb Q\\
&(r\notin A\land q<-r)\}.
\end{aligned}}
\UebQuellen{\FormulaRefAuto{RealCutNegation}[def]}
&
\textit{Additives Inverses und Ordnungsumkehr}
\UebFormel{\begin{aligned}
&-A\in\mathbb R,\qquad A+(-A)=0,\\
&-(-A)=A,\qquad -0=0,\\
&A\leq B\leftrightarrow-B\leq-A.
\end{aligned}}
\UebQuellen{\FormulaRefAuto{RealCutAdditiveClosure}[thm];
\FormulaRefAuto{RealCutAdditiveLaws}[thm];
\FormulaRefAuto{RealCutNegationOrder}[thm]}
\\ \addlinespace[.8ex]

\textbf{Positives Produkt und Vorzeichenfortsetzung}
Für \(0\leq A,B\):
\UebFormel{\begin{aligned}
A\RealPosMul B\coloneqq{}&0\cup
\{q\in\mathbb Q\mid\\
&\exists a\in A\;\exists b\in B\\
&(\RatZero<a\land\RatZero<b\land q<ab)\}.
\end{aligned}}
Das allgemeine Produkt entsteht durch die vier Vorzeichenfälle;
\(1\coloneqq\RealHat{\RatOne}\).
\UebQuellen{\FormulaRefAuto{RealCutPositiveProduct}[def];
\FormulaRefAuto{RealCutMultiplication}[def];
\FormulaRefAuto{RealCutOne}[def]}
&
\textit{Multiplikation und Ordnung}
\UebFormel{\begin{aligned}
&AB\in\mathbb R,\qquad(AB)C=A(BC),\\
&AB=BA,\qquad A1=A,\\
&A(B+C)=AB+AC,\\
&0\leq A,\ 0\leq B\vdash0\leq AB.
\end{aligned}}
\UebQuellen{\FormulaRefAuto{RealCutMultiplicativeClosure}[thm];
\FormulaRefAuto{RealCutMultiplicativeLaws}[thm];
\FormulaRefAuto{RealCutOrderArithmeticCompatibility}[thm]}
\\ \addlinespace[.8ex]

\textbf{Reziproker Schnitt}
Für \(0<A\):
\UebFormel{\begin{aligned}
&\RealPosInv A\coloneqq0\cup
\{q\in\mathbb Q\mid\\
&\quad\exists r\in\mathbb Q\;(\RatZero<r\land r\notin A\\
&\qquad{}\land q<r^{-1})\}.
\end{aligned}}
Für \(A<0\) ist \(A^{-1}\coloneqq-\RealPosInv{(-A)}\).
\UebQuellen{\FormulaRefAuto{RealCutPositiveReciprocal}[def];
\FormulaRefAuto{RealCutReciprocal}[def]}
&
\textit{Inversengesetz und Erhalt der rationalen Operationen}
\UebFormel{\begin{aligned}
&A\neq0\vdash A^{-1}\in\mathbb R,\quad AA^{-1}=1,\\
&\RealEmb(p+q)=\RealEmb(p)+\RealEmb(q),\\
&\RealEmb(pq)=\RealEmb(p)\RealEmb(q),\\
&p\neq\RatZero\vdash
\RealEmb(p^{-1})=(\RealEmb(p))^{-1}.
\end{aligned}}
\UebQuellen{\FormulaRefAuto{RealCutMultiplicativeClosure}[thm];
\FormulaRefAuto{RealCutMultiplicativeLaws}[thm];
\FormulaRefAuto{RationalToRealEmbeddingArithmetic}[thm]}
\\ \addlinespace[.8ex]

\textbf{Betrag und Abstand}
\UebFormel{\begin{aligned}
|A|&\coloneqq
\begin{cases}A,&0\leq A,\\-A,&A<0,\end{cases}\\
\RealDist A B&\coloneqq|A-B|.
\end{aligned}}
\UebQuellen{\FormulaRefAuto{RealAbsoluteValue}[def];
\FormulaRefAuto{RealDistance}[def]}
&
\textit{Betrags- und Abstandsgesetze}
\UebFormel{\begin{aligned}
&|AB|=|A||B|,\qquad|A+B|\leq|A|+|B|,\\
&0\leq\RealDist A B,\\
&\RealDist A B=0\leftrightarrow A=B,\\
&\RealDist A B=\RealDist B A,\\
&\RealDist A C\leq\RealDist A B+\RealDist B C.
\end{aligned}}
\UebQuellen{\FormulaRefAuto{RealAbsoluteValueLaws}[thm];
\FormulaRefAuto{RealDistanceLaws}[thm]}
\\ \addlinespace[.8ex]

\textbf{Intervalle und Epsilon-Umgebungen}
\UebFormel{\begin{aligned}
[A,B]_{\mathbb R}&\coloneqq\{x\in\mathbb R\mid A\leq x\leq B\},\\
(A,B)_{\mathbb R}&\coloneqq\{x\in\mathbb R\mid A<x<B\},\\
B_\varepsilon(x)&\coloneqq
\{y\in\mathbb R\mid|y-x|<\varepsilon\},
\end{aligned}}
wobei \(\varepsilon>0\). Halboffene Intervalle und Strahlen werden
durch die entsprechenden nichtstrikten oder einseitigen Bedingungen definiert.
\UebQuellen{\FormulaRefAuto{RealBoundedIntervals}[def];
\FormulaRefAuto{RealUnboundedIntervals}[def];
\FormulaRefAuto{RealEpsilonNeighborhood}[def]}
&
\textit{Umgebungen sind offene Intervalle}
\UebFormel{B_\varepsilon(x)=(x-\varepsilon,x+\varepsilon)_{\mathbb R}.}
Die Ordnungs- und die Abstandsbeschreibung stimmen damit überein.
\UebQuellen{\FormulaRefAuto{RealEpsilonNeighborhoodInterval}[thm]}
\\ \addlinespace[.8ex]

\textbf{Archimedische Einbettung}
Für \(\eta\colon\mathbb N\to K\) in einer total geordneten Menge:
\UebFormel{\begin{aligned}
&\operatorname{Arch}(K,\preccurlyeq,\eta)\coloneqq\\
&\quad\forall x\in K\;\exists n\in\mathbb N\\
&\qquad(x\preccurlyeq\eta(n)\land x\neq\eta(n)).
\end{aligned}}
Im reellen Modell ist
\(\eta=\RealEmb\circ\RatEmbed\circ\IntEmbed\).
\UebQuellen{\FormulaRefAuto{ArchimedeanOrderedEmbedding}[def]}
&
\textit{Archimedizität, Ganzzahlabschnitt und Dichte}
\UebFormel{\begin{aligned}
&\forall x\in\mathbb R\;\exists n\in\mathbb N\;(x<\eta(n)),\\
&\forall x\in\mathbb R\;\exists z\in\mathbb Z:\\
&\quad\RealEmb(\RatEmbed(z))\leq x,\\
&\quad x<\RealEmb(\RatEmbed(z+\IntOne)),\\
&x<y\vdash\exists q\in\mathbb Q\;(x<\RealEmb(q)<y).
\end{aligned}}
\UebQuellen{\FormulaRefAuto{RealArchimedeanProperty}[thm];
\FormulaRefAuto{RealIntegerPartProperty}[thm];
\FormulaRefAuto{RationalDenseInReal}[thm]}
\\ \addlinespace[.8ex]

\textbf{Vollständig angeordneter Körper}
Die Strukturdefinition verbindet die Körpergesetze, eine kompatible
totale Ordnung mit \(0_K<1_K\) und die Supremumseigenschaft:
\UebFormel{\begin{aligned}
&\varnothing\neq S\subseteq K,\quad
\exists u\in K\;\forall s\in S\;(s\preccurlyeq u)\\
&\qquad\vdash\sup_K S\text{ existiert in }K.
\end{aligned}}
\UebQuellen{\FormulaRefAuto{CompleteOrderedFieldAxioms}[def]}
&
\textit{Modell und Eindeutigkeit}
Das Schnittmodell ist ein vollständig angeordneter Körper.
Für jeden solchen Körper \(K\) gibt es genau eine Bijektion
\(\Phi\colon\mathbb R\bij K\), die die folgenden Gesetze erfüllt:
\UebFormel{\begin{aligned}
&\Phi(x+y)=\Phi(x)\boxplus\Phi(y),\\
&\Phi(xy)=\Phi(x)\boxtimes\Phi(y),\\
&x\leq y\leftrightarrow\Phi(x)\preccurlyeq\Phi(y),\\
&\Phi(0)=0_K,\qquad\Phi(1)=1_K.
\end{aligned}}
\UebQuellen{\FormulaRefAuto{RealCutsCompleteOrderedField}[thm];
\FormulaRefAuto{CompleteOrderedFieldUniqueness}[thm]}
\\
\end{BandUebersicht}

\chapter[Band 20: Endliche Mengen]{Band 20 auf einen Blick}
\noindent\textbf{Endliche Mengen}\par\medskip
\noindent\emph{Lesart.}
\(\Finite M\) ist das primitive Endlichkeitsprädikat; seine
Bijektionscharakterisierung ist ein Satz.
\(\NatSeg n=\{k\in\mathbb N\mid k<n\}\) bezeichnet einen
Peano-Anfangsabschnitt und \(\EqCard\) die Gleichmächtigkeit durch Bijektion.
Die im Original ausdrücklich angegebenen Auswahlabhängigkeiten der
Dedekind-/Tarski- und Unendlichkeitscharakterisierungen gelten auch hier.

\begin{BandUebersicht}
\textbf{Axiome: leere Menge und Adjunktion}
\UebFormel{\begin{aligned}
&\Finite\varnothing,\\
&\Finite A,\ a\notin A\vdash\Finite{A\cup\{a\}}.
\end{aligned}}
\UebQuellen{\FormulaRefAuto{Endlichkeitsaxiome}[def];
\FormulaRefAuto{FiniteEmptyAxiom}[ax];
\FormulaRefAuto{FiniteAdjunctionAxiom}[ax]}
&
\textit{Abschlüsse der Endlichkeit}
\UebFormel{\begin{aligned}
&\Finite A\vdash\Finite{A\cup\{a\}},\\
&\Finite A,\ \Finite B\vdash\Finite{A\cup B},\\
&B\subseteq A,\ \Finite A\vdash\Finite B,\\
&F\colon A\to B,\ \Finite A\vdash\Finite{F[A]}.
\end{aligned}}
\UebQuellen{\FormulaRefAuto{FiniteInsert}[thm];
\FormulaRefAuto{FiniteUnion}[thm];
\FormulaRefAuto{FiniteSubset}[thm];
\FormulaRefAuto{FiniteImage}[thm]}
\\ \addlinespace[.8ex]

\textbf{Axiom: Induktion über endliche Mengen}
Aus dem Anfang \(P(\varnothing)\) und dem Schritt
\UebFormel{\begin{aligned}
&\Finite A,\ P(A),\ a\notin A\\
&\qquad\vdash P(A\cup\{a\})
\end{aligned}}
folgt unter der Voraussetzung \(\Finite M\) die Aussage \(P(M)\).
\UebQuellen{\FormulaRefAuto{FiniteInductionRule}[ax]}
&
\textit{Bijektionscharakterisierung}
\UebFormel{\begin{aligned}
&n\in\mathbb N\vdash\Finite{\NatSeg n},\\
&\Finite M\eqvdash
\exists n\in\mathbb N\;(\NatSeg n\EqCard M).
\end{aligned}}
Für die von-Neumann-Realisierung gilt außerdem
\UebFormel{\Finite M\eqvdash
\exists n\in\mathbb N\;(n\EqCard M).}
\UebQuellen{\FormulaRefAuto{PeanoNatSegFinite}[thm];
\FormulaRefAuto{FiniteEqCardEquiv}[thm];
\FormulaRefAuto{FiniteVonNeumannEqCardEquiv}[thm]}
\\ \addlinespace[.8ex]

\textbf{Aufzählungsmengen}
Für \(a\colon\mathbb N\to A\) und \(n\in\mathbb N\):
\UebFormel{\operatorname{Enum}_a(n)\coloneqq a[\NatSeg n].}
\UebQuellen{\FormulaRefAuto{EnumSetDef}[def]}
&
\textit{Rekursion und Elementkriterium}
\UebFormel{\begin{aligned}
&\operatorname{Enum}_a(0)=\varnothing,\\
&\operatorname{Enum}_a(n+1)
=\operatorname{Enum}_a(n)\cup\{a_n\},\\
&x\in\operatorname{Enum}_a(n)\\
&\qquad\leftrightarrow\exists i\in\NatSeg n\;(x=a_i).
\end{aligned}}
\UebQuellen{\FormulaRefAuto{EnumSetZero}[thm];
\FormulaRefAuto{EnumSetSucc}[thm];
\FormulaRefAuto{EnumSetMembership}[thm]}
\\ \addlinespace[.8ex]

\textbf{Tarski-Endlichkeit}
\(\TarskiFinite(M)\) fordert für jede
\(\varnothing\neq B\subseteq\powerset(M)\):
\UebFormel{\exists x\in B\;\forall y\in B\;
(x\subseteq y\to y=x).}
\(\TarskiMinFinite(M)\) ersetzt hier \(x\subseteq y\) durch
\(y\subseteq x\).
\UebQuellen{\FormulaRefAuto{Tarski-endliche Menge}[def];
\FormulaRefAuto{minimal Tarski-endliche Menge}[def]}
&
\textit{Äquivalente Endlichkeitsbegriffe}
\UebFormel{\begin{aligned}
&\Finite M\leftrightarrow
\neg\exists H\colon\mathbb N\inj M,\\
&\Finite M\leftrightarrow\TarskiFinite(M),\\
&\TarskiFinite(M)\leftrightarrow\TarskiMinFinite(M).
\end{aligned}}
Die nichttriviale Rückrichtung aus dem Injektionsausschluss
verwendet im Band das Auswahlaxiom.
\UebQuellen{\FormulaRefAuto{\Finite{M}\eqvdash
\neg\exists H\colon\mathbb{N}\inj M}[thm];
\FormulaRefAuto{\Finite{M}\eqvdash\TarskiFinite(M)}[thm];
\FormulaRefAuto{TarskiMaxMinFiniteEquiv}[thm]}
\\ \addlinespace[.8ex]

\textbf{Kardinalzahl einer endlichen Menge}
Unter \(\Finite M\):
\UebFormel{\begin{aligned}
\FiniteCard(M)\coloneqq\iota n\;&
(n\in\mathbb N\\
&{}\land\NatSeg n\EqCard M).
\end{aligned}}
\UebQuellen{\FormulaRefAuto{FiniteCardinality}[def]}
&
\textit{Eindeutigkeit und Zählung}
\UebFormel{\begin{aligned}
&\exists!n\in\mathbb N\;(\NatSeg n\EqCard M),\\
&\FiniteCard(M)\in\mathbb N,\\
&\NatSeg{\FiniteCard(M)}\EqCard M.
\end{aligned}}
Die Definition wählt somit einen nachweislich eindeutigen
natürlichen Zählwert.
\UebQuellen{\FormulaRefAuto{FiniteCardWitnessUnique}[thm];
\FormulaRefAuto{FiniteCardWellDefined}[thm];
\FormulaRefAuto{FiniteCardNatural}[thm];
\FormulaRefAuto{FiniteCardCounts}[thm]}
\\ \addlinespace[.8ex]

\textbf{Binäre Teilmengencodierung}
Für \(X\subseteq\NatSeg n\):
\UebFormel{\begin{aligned}
\varepsilon_X(i)&\coloneqq
\begin{cases}1,&i\in X,\\0,&i\notin X,\end{cases}\\
c_n(X)&\coloneqq\sum_{i<n}\varepsilon_X(i)2^i.
\end{aligned}}
\UebQuellen{\FormulaRefAuto{BinarySubsetCode}[def]}
&
\textit{Potenzmengen zählen}
\UebFormel{c_n\colon\powerset(\NatSeg n)\bij\NatSeg{2^n}.}
Für \(\Finite A\) und \(\FiniteCard(A)=n\):
\UebFormel{\begin{aligned}
&\FiniteCard(\powerset(A))=2^n,\\
&\FiniteCard(\PowNE A)=2^n-1.
\end{aligned}}
Dabei ist \(\PowNE A\) die Menge der nichtleeren Teilmengen.
\UebQuellen{\FormulaRefAuto{BinarySubsetCodeBijective}[thm];
\FormulaRefAuto{FinitePowerSetCardinality}[thm];
\FormulaRefAuto{FiniteNonemptyPowerSetCardinality}[thm]}
\\ \addlinespace[.8ex]

\textbf{Maximalelementeigenschaft}
\UebFormel{\begin{aligned}
&\operatorname{HatMax}_{A,\preceq}(C)\coloneqq\\
&\quad(C\subseteq A\land C\neq\varnothing\\
&\qquad\to\exists m\in C\;\forall b\in C\\
&\qquad\qquad(m\preceq b\to b=m)).
\end{aligned}}
\UebQuellen{\FormulaRefAuto{FinitePosetMaximalProperty}[def]}
&
\textit{Endliche partielle Ordnungen}
Unter \(\PartOrd{A,\preceq}\),
\(\varnothing\neq B\subseteq A\) und \(\Finite B\):
\UebFormel{\begin{aligned}
&\exists m\in B\;\forall b\in B\;(m\preceq b\to b=m),\\
&\exists \ell\in B\;\forall b\in B\;(b\preceq\ell\to b=\ell).
\end{aligned}}
Ein maximales bzw. minimales Element muss kein größtes bzw.
kleinstes Element sein.
\UebQuellen{\FormulaRefAuto{FinitePosetMaximalElement}[thm];
\FormulaRefAuto{FinitePosetMinimalElement}[thm]}
\\ \addlinespace[.8ex]

\textbf{Unendlich, abzählbar und überabzählbar}
\UebFormel{\begin{aligned}
\Infinite A&\coloneqq\neg\Finite A,\\
\AtMostCountable A&\coloneqq A\CardLeq\mathbb N,\\
\CountablyInfinite A&\coloneqq A\EqCard\mathbb N,\\
\Uncountable A&\coloneqq\neg\AtMostCountable A.
\end{aligned}}
\UebQuellen{\FormulaRefAuto{InfiniteDef}[def];
\FormulaRefAuto{AtMostCountableDef}[def];
\FormulaRefAuto{CountablyInfiniteDef}[def];
\FormulaRefAuto{UncountableDef}[def]}
&
\textit{Unendlichkeit und abzählbare Alternative}
\UebFormel{\begin{aligned}
&\Infinite{\mathbb N},\\
&\Infinite A\leftrightarrow\exists H\colon\mathbb N\inj A,\\
&\AtMostCountable A\leftrightarrow\\
&\qquad(\Finite A\lor\CountablyInfinite A).
\end{aligned}}
Für die angegebenen Rückführungen auf Injektionen wird die im Band
festgehaltene Auswahlvoraussetzung benutzt.
\UebQuellen{\FormulaRefAuto{NaturalNumbersInfinite}[thm];
\FormulaRefAuto{InfiniteInjectionCharacterization}[thm];
\FormulaRefAuto{AtMostCountableDichotomy}[thm]}
\\ \addlinespace[.8ex]

\textbf{Einermengenabbildung und Diagonalmenge}
\UebFormel{\begin{aligned}
&\SingletonEmbedding A(x)\coloneqq\{x\},\\
&F\colon A\to\powerset(A)\vdash\\
&\qquad D_F^A\coloneqq\{x\in A\mid x\notin F(x)\}.
\end{aligned}}
\UebQuellen{\FormulaRefAuto{SingletonEmbeddingDef}[def];
\FormulaRefAuto{CantorDiagonalSetDef}[def]}
&
\textit{Cantors strikter Potenzmengensatz}
\UebFormel{\begin{aligned}
&\SingletonEmbedding A\colon A\inj\powerset(A),\\
&x\in A\vdash F(x)\neq D_F^A,\\
&A\CardLt\powerset(A).
\end{aligned}}
Insbesondere gibt es keine größte Mächtigkeitsstufe.
\UebQuellen{\FormulaRefAuto{SingletonPowerSetInjection}[thm];
\FormulaRefAuto{CantorDiagonalMissed}[thm];
\FormulaRefAuto{CantorStrictPowerSet}[thm];
\FormulaRefAuto{NoGreatestCardinality}[thm]}
\\ \addlinespace[.8ex]

\textbf{Kontinuumshypothese als Aussage}
\UebFormel{\begin{aligned}
\mathsf{CH}\coloneqq\neg\exists A\;(&\mathbb N\CardLt A\\
&{}\land A\CardLt\powerset(\mathbb N)).
\end{aligned}}
Dies ist die Definition einer Hypothese; der Band beweist sie nicht.
\UebQuellen{\FormulaRefAuto{ContinuumHypothesisRelational}[def]}
&
\textit{Bewiesene Unendlichkeitsstufen}
\UebFormel{\begin{aligned}
&\mathbb N\CardLt\powerset(\mathbb N),\\
&\powerset(\mathbb N)\CardLt
\powerset(\powerset(\mathbb N)),\\
&\mathbb N\CardLt\mathbb R,\qquad\Uncountable{\mathbb R}.
\end{aligned}}
Keine Folge natürlicher Indizes zählt sämtliche reellen Zahlen auf.
\UebQuellen{\FormulaRefAuto{FirstInfinitePowerSetLevels}[thm];
\FormulaRefAuto{NaturalsStrictlyLessPowerfulThanReals}[thm];
\FormulaRefAuto{RealsUncountable}[thm];
\FormulaRefAuto{CantorNoRealSurjection}[thm]}
\\
\end{BandUebersicht}

\chapter[Band 21: Folgen]{Band 21 auf einen Blick}
\noindent\textbf{Folgen}\par\medskip
\noindent\emph{Lesart.}
Eine Familie ist eine Funktion mit ausdrücklich angegebener Indexmenge.
Bei reellen Folgen sind \(a,b\colon\mathbb N\to\mathbb R\) und
\(L,M,c\in\mathbb R\); \(\varepsilon>0\) bezeichnet eine positive reelle
Zahl. Eine Teilfolgenabbildung \(\varphi\colon\mathbb N\to\mathbb N\)
ist stets strikt wachsend. Die letzten vier Zeilen verwenden die
eigenständigen Bezeichnungen der Mogiljanskaja-Familien.

\begin{BandUebersicht}
\textbf{Indizierte Familie und Wertemenge}
Für \(a\colon I\to M\):
\UebFormel{\begin{aligned}
(a_i)_{i\in I}&\coloneqq a,\qquad a_i\coloneqq a(i),\\
\operatorname{Bild}(a)&\coloneqq a[I].
\end{aligned}}
\UebQuellen{\FormulaRefAuto{IndexedFamilyDef}[def];
\FormulaRefAuto{IndexedFamilyImageDef}[def]}
&
\textit{Elementkriterium und surjektive Familien}
\UebFormel{\begin{aligned}
&i\in I\vdash a_i\in M,\\
&x\in\operatorname{Bild}(a)
\leftrightarrow\exists i\in I\;(x=a_i),\\
&\operatorname{Bild}(a)\subseteq M,\\
&a\colon I\sur M\leftrightarrow\operatorname{Bild}(a)=M.
\end{aligned}}
\UebQuellen{\FormulaRefAuto{IndexedFamilyCoordinateTyping}[thm];
\FormulaRefAuto{IndexedFamilyImageMembership}[thm];
\FormulaRefAuto{IndexedFamilyImageSubsetCodomain}[thm];
\FormulaRefAuto{IndexedFamilySurjectiveIffImage}[thm]}
\\ \addlinespace[.8ex]

\textbf{Einschränkung und Umindizierung}
Für \(J\subseteq I\), bzw. \(\psi\colon J\to I\):
\UebFormel{\begin{aligned}
(a_i)_{i\in J}&\coloneqq a\restriction_J,\\
(a_{\psi(j)})_{j\in J}&\coloneqq a\circ\psi.
\end{aligned}}
\UebQuellen{\FormulaRefAuto{IndexedFamilyRestrictionDef}[def];
\FormulaRefAuto{IndexedFamilyReindexingDef}[def]}
&
\textit{Werte bleiben über die Indizes kontrollierbar}
\UebFormel{\begin{aligned}
&(a\restriction_J)\colon J\to M,\\
&j\in J\vdash(a\restriction_J)_j=a_j,\\
&\operatorname{Bild}(a\circ\psi)=a[\psi[J]].
\end{aligned}}
Familien mit derselben Indexmenge sind genau dann gleich, wenn
alle entsprechenden Glieder übereinstimmen.
\UebQuellen{\FormulaRefAuto{IndexedFamilyRestrictionFacts}[thm];
\FormulaRefAuto{IndexedFamilyReindexingImage}[thm];
\FormulaRefAuto{IndexedFamilyExtensionality}[thm]}
\\ \addlinespace[.8ex]

\textbf{Endliche Indexkonventionen}
Für \(n\in\mathbb N\):
\UebFormel{\begin{aligned}
&\NatSeg n=\{i\in\mathbb N\mid i<n\},\\
&\PosNatSeg n\coloneqq
\{i\in\mathbb N\mid1\leq i\leq n\},\\
&\lambda_n\colon\NatSeg n\to\PosNatSeg n,\\
&\lambda_n(i)\coloneqq i+1.
\end{aligned}}
Endliche Folgen haben eine dieser beiden Mengen als Definitionsbereich.
\UebQuellen{\FormulaRefAuto{ZeroBasedFiniteSequenceDef}[def];
\FormulaRefAuto{PositiveNaturalSegmentDef}[def];
\FormulaRefAuto{PositiveFiniteSequenceDef}[def];
\FormulaRefAuto{FiniteSequenceIndexShiftDef}[def]}
&
\textit{Bijektive Indexverschiebung}
\UebFormel{\lambda_n\colon\NatSeg n\bij\PosNatSeg n.}
Für \(a\colon\NatSeg n\to M\) ist
\(a\circ\lambda_n^{-1}\colon\PosNatSeg n\to M\) die entsprechende
positiv indizierte Folge; zurück führt Komposition mit \(\lambda_n\).
\UebQuellen{\FormulaRefAuto{FiniteSequenceIndexShiftBijective}[thm];
\FormulaRefAuto{FiniteSequenceConventionConversion}[thm]}
\\ \addlinespace[.8ex]

\textbf{Unendliche Folgen, Beispiele und Rekursionsdaten}
\UebFormel{\begin{aligned}
&(a_n)_{n\in\mathbb N}\coloneqq a
\quad(a\colon\mathbb N\to M),\\
&(\operatorname{const}_{I,c})_i\coloneqq c,\\
&(\iota_{\mathbb N})_n\coloneqq n.
\end{aligned}}
Für eine Schrittvorschrift sind \(x_0\in M\) und \(F\colon M\to M\)
gegeben.
\UebQuellen{\FormulaRefAuto{InfiniteSequenceDef}[def];
\FormulaRefAuto{ConstantIndexedFamilyDef}[def];
\FormulaRefAuto{NaturalEnumerationSequenceDef}[def]}
&
\textit{Eindeutige Rekursion und verschachtelte Mengenfolgen}
\UebFormel{\begin{aligned}
&\exists!a\colon\mathbb N\to M:\\
&\qquad a_0=x_0,\quad
\forall n\in\mathbb N\;a_{n+1}=F(a_n).
\end{aligned}}
Für eine Mengenfolge \(S\) gilt:
\UebFormel{\begin{aligned}
&\forall k\in\mathbb N\;(S_k\subseteq S_{k+1}),\quad m\leq n\\
&\qquad\vdash S_m\subseteq S_n.
\end{aligned}}
\UebQuellen{\FormulaRefAuto{RecursiveSequenceExistenceUnique}[thm];
\FormulaRefAuto{SubsetSequenceNesting}[thm]}
\\ \addlinespace[.8ex]

\textbf{Restfolge und Teilfolge}
\UebFormel{\begin{aligned}
a^{\langle r\rangle}_n&\coloneqq a_{n+r}
\quad(r\in\mathbb N),\\
(a_{\varphi(n)})_{n\in\mathbb N}&\coloneqq a\circ\varphi.
\end{aligned}}
Die strikte Monotonie von \(\varphi\) ist Teil der
Teilfolgendefinition.
\UebQuellen{\FormulaRefAuto{SequenceTailDef}[def];
\FormulaRefAuto{SubsequenceDef}[def]}
&
\textit{Grenzwert bleibt erhalten}
\UebFormel{\begin{aligned}
&a_n\to L\leftrightarrow a^{\langle r\rangle}_n\to L,\\
&a_n\to L\vdash a_{\varphi(n)}\to L.
\end{aligned}}
Auch endlich viele Änderungen der Glieder verändern den Grenzwert nicht.
\UebQuellen{\FormulaRefAuto{SequenceTailSameLimit}[thm];
\FormulaRefAuto{SubsequenceOfConvergentSequence}[thm];
\FormulaRefAuto{FiniteModificationPreservesLimit}[thm]}
\\ \addlinespace[.8ex]

\textbf{Konvergenz und Beschränktheit}
\UebFormel{\begin{aligned}
&a_n\to L\coloneqq\\
&\quad\forall\varepsilon>0\;\exists N\in\mathbb N\\
&\qquad\forall n\geq N\;|a_n-L|<\varepsilon,\\
&\operatorname{Beschr}(a)\coloneqq\\
&\quad\exists C\geq0\;\forall n\in\mathbb N\;|a_n|\leq C.
\end{aligned}}
\UebQuellen{\FormulaRefAuto{RealSequenceConvergenceDef}[def];
\FormulaRefAuto{BoundedRealSequenceDef}[def]}
&
\textit{Eindeutigkeit, Beschränktheit und Ordnung}
\UebFormel{\begin{aligned}
&a_n\to L,\ a_n\to M\vdash L=M,\\
&a_n\to L\vdash\operatorname{Beschr}(a),\\
&a_n\to L,\ b_n\to M,\ \forall n\;a_n\leq b_n\\
&\qquad\vdash L\leq M.
\end{aligned}}
\UebQuellen{\FormulaRefAuto{RealSequenceLimitUnique}[thm];
\FormulaRefAuto{ConvergentRealSequenceBounded}[thm];
\FormulaRefAuto{RealSequenceLimitPreservesOrder}[thm]}
\\ \addlinespace[.8ex]

\textbf{Nullfolge}
\UebFormel{a\text{ ist Nullfolge}\coloneqq a_n\to0.}
Mit punktweiser Subtraktion ist \((a-L)_n\coloneqq a_n-L\).
\UebQuellen{\FormulaRefAuto{NullSequenceDef}[def];
\FormulaRefAuto{PointwiseRealSequenceOperationsDef}[def]}
&
\textit{Abschätzung und Einschließung}
\UebFormel{\begin{aligned}
&a_n\to L\leftrightarrow a_n-L\to0,\\
&b_n\to0,\ \forall n\;|a_n|\leq b_n\vdash a_n\to0,\\
&u_n\to L,\ v_n\to L,\ \forall n\;u_n\leq a_n\leq v_n\\
&\qquad\vdash a_n\to L.
\end{aligned}}
Hier sind auch \(u,v\) reelle Folgen.
\UebQuellen{\FormulaRefAuto{ConvergenceIffDifferenceNullSequence}[thm];
\FormulaRefAuto{NullSequenceMajorantCriterion}[thm];
\FormulaRefAuto{RealSequenceSqueezeTheorem}[thm]}
\\ \addlinespace[.8ex]

\textbf{Monotone Folgen}
\UebFormel{\begin{aligned}
\operatorname{Mon}_{\uparrow}(a)&\coloneqq
\forall n\in\mathbb N\;(a_n\leq a_{n+1}),\\
\operatorname{Mon}_{\downarrow}(a)&\coloneqq
\forall n\in\mathbb N\;(a_{n+1}\leq a_n).
\end{aligned}}
\UebQuellen{\FormulaRefAuto{MonotoneRealSequenceDef}[def]}
&
\textit{Monotone Konvergenz}
\UebFormel{\begin{aligned}
&\operatorname{Mon}_{\uparrow}(a),\
\exists C\;\forall n\;(a_n\leq C)\\
&\qquad\vdash a_n\to\sup_{\mathbb R,\RealLe}a[\mathbb N],\\
&\operatorname{Mon}_{\downarrow}(a),\
\exists C\;\forall n\;(C\leq a_n)\\
&\qquad\vdash a_n\to\inf_{\mathbb R,\RealLe}a[\mathbb N].
\end{aligned}}
\UebQuellen{\FormulaRefAuto{MonotoneRealSequenceConvergence}[thm]}
\\ \addlinespace[.8ex]

\textbf{Punktweise Rechenoperationen}
\UebFormel{\begin{aligned}
(a+b)_n&\coloneqq a_n+b_n,\\
(ab)_n&\coloneqq a_nb_n,\qquad(ca)_n\coloneqq ca_n,\\
(a/b)_n&\coloneqq a_n/b_n
\quad(\forall n\;b_n\neq0).
\end{aligned}}
\UebQuellen{\FormulaRefAuto{PointwiseRealSequenceOperationsDef}[def];
\FormulaRefAuto{PointwiseRealSequenceReciprocalQuotientDef}[def]}
&
\textit{Grenzwertregeln}
Unter \(a_n\to L,\ b_n\to M\):
\UebFormel{\begin{aligned}
&a_n+b_n\to L+M,\qquad ca_n\to cL,\\
&a_nb_n\to LM,\qquad |a_n|\to|L|,\\
&M\neq0,\ \forall n\;b_n\neq0
\vdash a_n/b_n\to L/M.
\end{aligned}}
\UebQuellen{\FormulaRefAuto{RealSequenceLinearLimitRules}[thm];
\FormulaRefAuto{RealSequenceProductLimitRule}[thm];
\FormulaRefAuto{AbsoluteValuePreservesSequenceLimits}[thm];
\FormulaRefAuto{RealSequenceQuotientLimitRule}[thm]}
\\ \addlinespace[.8ex]

\textbf{Cauchy-Folge}
\UebFormel{\begin{aligned}
&\operatorname{Cauchy}(a)\coloneqq\\
&\quad\forall\varepsilon>0\;\exists N\in\mathbb N\\
&\qquad\forall m,n\geq N\;|a_m-a_n|<\varepsilon,\\
&\operatorname{Konv}(a)\coloneqq
\exists L\in\mathbb R\;(a_n\to L).
\end{aligned}}
\UebQuellen{\FormulaRefAuto{RealCauchySequenceDef}[def];
\FormulaRefAuto{RealSequenceConvergentDivergentDef}[def]}
&
\textit{Cauchy-Kriterium der reellen Zahlen}
\UebFormel{\begin{aligned}
&\operatorname{Konv}(a)\leftrightarrow\operatorname{Cauchy}(a),\\
&\operatorname{Cauchy}(a)\vdash\operatorname{Beschr}(a).
\end{aligned}}
Die Existenz eines Grenzwerts ist hier eine Folgerung aus der
Vollständigkeit der reellen Ordnung.
\UebQuellen{\FormulaRefAuto{RealSequenceCauchyCriterion}[thm];
\FormulaRefAuto{RealCauchySequenceBounded}[thm]}
\\ \addlinespace[.8ex]

\textbf{Mogiljanskaja: Grundschichten und Familien}
Von hier an bezeichnen \(a,b\) die folgenden Mengenfamilien:
\UebFormel{\begin{aligned}
L&\coloneqq\{0\}\times\mathbb N,\qquad a_i\coloneqq(0,i),\\
D&\coloneqq\{1\}\times(\mathbb N\times\mathbb N),\\
b_{ij}&\coloneqq(1,(i,j)).
\end{aligned}}
\UebQuellen{\FormulaRefAuto{MogiljanskajaLayerLDef}[def];
\FormulaRefAuto{MogiljanskajaLayerDDef}[def];
\FormulaRefAuto{MogiljanskajaAElementsDef}[def];
\FormulaRefAuto{MogiljanskajaBElementsDef}[def]}
&
\textit{Indizes, Bildmengen und Trennung}
\UebFormel{\begin{aligned}
&a_i=a_j\vdash i=j,\\
&b_{ij}=b_{mn}\vdash i=m,\ j=n,\\
&a[\mathbb N]=L,\qquad
b[\mathbb N\times\mathbb N]=D,\\
&L\cap D=\varnothing.
\end{aligned}}
\UebQuellen{\FormulaRefAuto{MogiljanskajaAIndexUnique}[thm];
\FormulaRefAuto{MogiljanskajaBIndexUnique}[thm];
\FormulaRefAuto{MogiljanskajaLayerLImage}[thm];
\FormulaRefAuto{MogiljanskajaLayerDImage}[thm];
\FormulaRefAuto{MogiljanskajaLayersDisjoint}[thm]}
\\ \addlinespace[.8ex]

\textbf{Marker und Reserve}
\UebFormel{\begin{aligned}
&c_0\coloneqq b_{00},\quad c_1\coloneqq b_{11},
\quad c_2\coloneqq b_{10},\\
&U\coloneqq\{b_{n+2,j}\mid n,j\in\mathbb N\},\\
&V\coloneqq U\cup\{c_2\},\quad P\coloneqq\{c_0,c_1\},\\
&D'\coloneqq P\cup V.
\end{aligned}}
\UebQuellen{\FormulaRefAuto{MogiljanskajaDPrimeDef}[def]}
&
\textit{Ausgesparte Punkte und disjunkter Kern}
\UebFormel{\begin{aligned}
&c_0,c_1,c_2\in D,\qquad c_0,c_1,c_2\notin U,\\
&c_0,c_1,c_2\text{ sind paarweise verschieden},\\
&D'\subseteq D,\qquad b_{01}\notin D',\\
&c_2\notin P,\qquad P\cap U=\varnothing.
\end{aligned}}
\UebQuellen{\FormulaRefAuto{MogiljanskajaReserveMarkerFacts}[thm];
\FormulaRefAuto{MogiljanskajaReserveCoreFacts}[thm]}
\\ \addlinespace[.8ex]

\textbf{Zeilenparametrisierung und punktierte Verschiebung}
\UebFormel{\begin{aligned}
&\delta(n,j)\coloneqq b_{nj},\quad
\upsilon(n,j)\coloneqq b_{n+2,j},\\
&\theta\coloneqq\delta\circ\upsilon^{-1},\\
&H_n\coloneqq b_{n+1,0},\\
&\sigma\coloneqq(\SuccShift H)^{-1}.
\end{aligned}}
Die Nachfolgerverschiebung \(\SuccShift H\) wirkt längs der
injektiven Folge \(H\).
\UebQuellen{\FormulaRefAuto{MogiljanskajaRowParametrizationsDef}[def];
\FormulaRefAuto{MogiljanskajaThetaDef}[def];
\FormulaRefAuto{MogiljanskajaShiftSequenceDef}[def];
\FormulaRefAuto{MogiljanskajaSigmaDef}[def]}
&
\textit{Bijektionen und tatsächliche Bildmenge}
\UebFormel{\begin{aligned}
&\delta\colon\mathbb N^2\bij D,\qquad
\upsilon\colon\mathbb N^2\bij U,\\
&\theta\colon U\bij D,\qquad\sigma\colon U\bij V,\\
&H\colon\mathbb N\inj V,\qquad H_0=c_2,\\
&H[\mathbb N]=\{b_{n+1,0}\mid n\in\mathbb N\}
\subsetneq V.
\end{aligned}}
Hier ist \(\mathbb N^2=\mathbb N\times\mathbb N\).
\UebQuellen{\FormulaRefAuto{MogiljanskajaRowParametrizationsBijective}[thm];
\FormulaRefAuto{MogiljanskajaThetaBijective}[thm];
\FormulaRefAuto{MogiljanskajaSigmaBijective}[thm];
\FormulaRefAuto{MogiljanskajaShiftSequenceFacts}[thm];
\FormulaRefAuto{MogiljanskajaShiftSequenceImage}[thm]}
\\ \addlinespace[.8ex]

\textbf{Familien mit vorgegebener Bildmenge}
\UebFormel{\begin{aligned}
&\nu\coloneqq\sigma\circ\upsilon,\\
&I_{D'}\coloneqq(\{0\}\times\NatSeg2)
\cup(\{1\}\times\mathbb N^2),\\
&\omega((0,0))\coloneqq c_0,\qquad
\omega((0,1))\coloneqq c_1,\\
&\omega((1,(n,j)))\coloneqq\nu(n,j).
\end{aligned}}
\UebQuellen{\FormulaRefAuto{MogiljanskajaVParametrizationDef}[def];
\FormulaRefAuto{MogiljanskajaDPrimeIndexSetDef}[def];
\FormulaRefAuto{MogiljanskajaDPrimeParametrizationDef}[def]}
&
\textit{Die gewünschten Bildmengen werden exakt erreicht}
\UebFormel{\begin{aligned}
&\nu\colon\mathbb N^2\bij V,\qquad\nu[\mathbb N^2]=V,\\
&\omega\colon I_{D'}\bij D',\qquad\omega[I_{D'}]=D'.
\end{aligned}}
Zusammen mit \(a,b,\upsilon\) liegen damit ausdrücklich parametrisierte
Familien für \(L,D,U,V,D'\) vor.
\UebQuellen{\FormulaRefAuto{MogiljanskajaVParametrizationBijective}[thm];
\FormulaRefAuto{MogiljanskajaDPrimeParametrizationBijective}[thm];
\FormulaRefAuto{MogiljanskajaFamilyImages}[thm]}
\\
\end{BandUebersicht}

\chapter[Band 22: Gerichtete Graphen]{Band 22 auf einen Blick}
\noindent\textbf{Gerichtete Graphen}\par\medskip
Vorausgesetzt sind Relationen und endliche Folgen. Es gilt
\(\NatSeg n=\{i\in\mathbb N\mid i<n\}\); \(n\) zählt Kanten,
eine Knotenfolge hat daher \(n+1\) Einträge. Schleifen sind zugelassen.

\begin{BandUebersicht}
\textbf{Graph und Kantenrelation}
\UebFormel{\begin{aligned}
&\GraphStruct{V,E}\ \coloneqq\ E\subseteq V\times V.
\end{aligned}}
\UebQuellen{\FormulaRefAuto{GraphDef}[def]}
&
\textbf{Endpunkte liegen im Träger}
\UebFormel{\begin{aligned}
&\GraphStruct{V,E},\ (x,y)\in E\ \vdash\\
&\qquad x\in V\land y\in V.
\end{aligned}}
\UebQuellen{\FormulaRefAuto{GraphEdgeEndpointTyping}[thm]} \\
\addlinespace[.8ex]

\textbf{Ein- und Ausnachbarschaften}
Für \(\GraphStruct{V,E}\) und \(x\in V\):\UebFormel{\begin{aligned}
&\OutNeighbors E x=\{y\in V\mid(x,y)\in E\},\\
&\InNeighbors E x=\{y\in V\mid(y,x)\in E\}.
\end{aligned}}
\UebQuellen{\FormulaRefAuto{DirectedOutNeighborhoodDef}[def];
\FormulaRefAuto{DirectedInNeighborhoodDef}[def]}
&
\textbf{Mitgliedschaftskriterien}
\UebFormel{\begin{aligned}
&\GraphStruct{V,E},\ x,y\in V\ \vdash\\
&y\in\OutNeighbors E x\ \leftrightarrow\ (x,y)\in E,\\
&y\in\InNeighbors E x\ \leftrightarrow\ (y,x)\in E.
\end{aligned}}
\UebQuellen{\FormulaRefAuto{DirectedNeighborhoodMembership}[thm]} \\
\addlinespace[.8ex]

\textbf{Umkehrung und Einschränkung}
Für \(\GraphStruct{V,E}\) und \(U\subseteq V\):\UebFormel{\begin{aligned}
&\ReverseEdges E=\{(y,x)\mid(x,y)\in E\},\\
&\DirectedInducedEdges E U=E\cap(U\times U).
\end{aligned}}
\UebQuellen{\FormulaRefAuto{DirectedReverseEdgeRelationDef}[def];
\FormulaRefAuto{DirectedInducedEdgeRelationDef}[def]}
&
\textbf{Beide Konstruktionen erhalten Graphen}
\UebFormel{\begin{aligned}
&\GraphStruct{V,E}\ \vdash\
  \GraphStruct{V,\ReverseEdges E},\\
&\GraphStruct{V,E},\ U\subseteq V\ \vdash\\
&\qquad\GraphStruct{U,\DirectedInducedEdges E U}.
\end{aligned}}
\UebQuellen{\FormulaRefAuto{DirectedReverseGraph}[thm];
\FormulaRefAuto{DirectedInducedSubgraph}[thm]} \\
\addlinespace[.8ex]

\textbf{Walk und injektiver Pfad}
Bei festem \(\GraphStruct{V,E}\) bestehen die Walkaxiome aus\UebFormel{\begin{aligned}
&n\in\mathbb N,\quad p:\NatSeg{n+1}\to V,\\
&\forall i\in\NatSeg n\quad(p(i),p(i+1))\in E.\\
&\DirectedPath V E p n\ \leftrightarrow\\
&\quad\DirectedWalk V E p n
 \land p:\NatSeg{n+1}\inj V.
\end{aligned}}
\UebQuellen{\FormulaRefAuto{DirectedWalkDef}[def];
\FormulaRefAuto{DirectedPathDef}[def]}
&
\textbf{Typisierung und einzelne Schritte}
\UebFormel{\begin{aligned}
&\GraphStruct{V,E},\
 \DirectedWalk V E p n\ \vdash\\
&\qquad p(0)\in V\land p(n)\in V,\\
&\GraphStruct{V,E},\
 \DirectedWalk V E p n,\ i\in\NatSeg n\ \vdash\\
&\qquad(p(i),p(i+1))\in E.
\end{aligned}}
\UebQuellen{\FormulaRefAuto{DirectedWalkEndpointTyping}[thm];
\FormulaRefAuto{DirectedWalkStepEdge}[thm]} \\
\addlinespace[.8ex]

\textbf{Kreis und Kreisfreiheit}
Ein Kreis ist ein Walk positiver Länge mit\UebFormel{\begin{aligned}
&p(0)=p(n),\\
&\forall i,j\in\NatSeg n\
   (p(i)=p(j)\to i=j).
\end{aligned}}Kreisfreiheit schließt sämtliche solchen Kreise aus.
\UebQuellen{\FormulaRefAuto{DirectedCycleDef}[def];
\FormulaRefAuto{DirectedAcyclicGraphDef}[def]}
&
\textbf{Charakterisierung der Kreisfreiheit}
\UebFormel{\begin{aligned}
&\AcyclicGraphStruct{V,E}\ \eqvdash\\
&\quad\GraphStruct{V,E}\ \land\\
&\quad\neg\exists n\in\mathbb N\,\exists p\
          \DirectedCycle V E p n.
\end{aligned}}
\UebQuellen{\FormulaRefAuto{DirectedAcyclicGraphCharacterization}[thm]} \\
\addlinespace[.8ex]

\textbf{Erreichbarkeit}
Für \(u,v\in V\) bedeutet \(u\GraphReach E v\):\UebFormel{\begin{aligned}
&u=v\ \lor\ \exists n\in\mathbb N\,\exists p\\
&\quad(\DirectedWalk V E p n\\
&\qquad{}\land p(0)=u\land p(n)=v).
\end{aligned}}
\UebQuellen{\FormulaRefAuto{DirectedReachabilityDef}[def]}
&
\textbf{Reflexivität und Endpunkttypisierung}
\UebFormel{\begin{aligned}
&\GraphStruct{V,E},\ u\in V\ \vdash\
       u\GraphReach E u,\\
&\GraphStruct{V,E},\ u\GraphReach E v\ \vdash\\
&\qquad u\in V\land v\in V.
\end{aligned}}
\UebQuellen{\FormulaRefAuto{DirectedReachabilityReflexive}[thm];
\FormulaRefAuto{DirectedReachabilityEndpointTyping}[thm]} \\
\addlinespace[.8ex]
\end{BandUebersicht}


\chapter[Band 23: Ungerichtete Graphen]{Band 23 auf einen Blick}
\noindent\textbf{Ungerichtete Graphen}\par\medskip
Die Kanten bleiben geordnete Paare. Ein eigenes Symmetrieaxiom
ergänzt den gerichteten Graphbegriff; Schleifenfreiheit und Zusammenhang
werden hier noch nicht gefordert.

\begin{BandUebersicht}
\textbf{Symmetrische Graphstruktur}
\UebFormel{\begin{aligned}
&\UndirectedGraphStruct{V,E}\ \leftrightarrow\\
&\quad\GraphStruct{V,E}\land\\
&\quad\forall x,y\in V\\
&\qquad((x,y)\in E\to(y,x)\in E).
\end{aligned}}
\UebQuellen{\FormulaRefAuto{UndirectedGraphDef}[def]}
&
\textbf{Adjazenz ist richtungsunabhängig}
\UebFormel{\begin{aligned}
&\UndirectedGraphStruct{V,E},\
 x,y\in V\ \vdash\\
&\qquad (x,y)\in E\ \leftrightarrow\ (y,x)\in E.
\end{aligned}}
\UebQuellen{\FormulaRefAuto{UndirectedAdjacencySymmetric}[thm]} \\
\addlinespace[.8ex]

\textbf{Ungerichtete Nachbarschaft}
\UebFormel{\begin{aligned}
&\UndirectedGraphStruct{V,E},\ x\in V\ \vdash\\
&\GraphNeighbors E x
  \coloneqq\{y\in V\mid(x,y)\in E\}.
\end{aligned}}
\UebQuellen{\FormulaRefAuto{UndirectedNeighborhoodDef}[def]}
&
\textbf{Gegenseitige Nachbarschaft}
\UebFormel{\begin{aligned}
&\UndirectedGraphStruct{V,E},\
 x,y\in V\ \vdash\\
&y\in\GraphNeighbors E x
  \ \leftrightarrow\ x\in\GraphNeighbors E y.
\end{aligned}}
\UebQuellen{\FormulaRefAuto{UndirectedNeighborhoodSymmetric}[thm]} \\
\addlinespace[.8ex]

\textbf{Induzierter Teilgraph}
\UebFormel{\begin{aligned}
&U\subseteq V,\qquad
 \DirectedInducedEdges E U=E\cap(U\times U).
\end{aligned}}
\UebQuellen{\FormulaRefAuto{DirectedInducedEdgeRelationDef}[def]}
&
\textbf{Symmetrie bleibt erhalten}
\UebFormel{\begin{aligned}
&\UndirectedGraphStruct{V,E},\ U\subseteq V\
 \vdash\\
&\qquad\UndirectedGraphStruct{U,\DirectedInducedEdges E U}.
\end{aligned}}
\UebQuellen{\FormulaRefAuto{UndirectedInducedSubgraph}[thm]} \\
\addlinespace[.8ex]
\end{BandUebersicht}

\chapter[Band 24: Schlichte ungerichtete Graphen]{Band 24 auf einen Blick}
\noindent\textbf{Schlichte ungerichtete Graphen}\par\medskip
Ausgangspunkt sind ungerichtete relationale Graphen.
Schlichtheit ergänzt ausschließlich die Irreflexivität; Zusammenhang
gehört zur nächsten eigenständigen Strukturstufe.

\begin{BandUebersicht}
\textbf{Schlichtheit als Zusatzaxiom}
\UebFormel{\begin{aligned}
&\SimpleUndirectedGraphStruct{V,E}\ \leftrightarrow\\
&\quad\UndirectedGraphStruct{V,E}\land\\
&\quad\forall x\in V\quad(x,x)\notin E.
\end{aligned}}
\UebQuellen{\FormulaRefAuto{SimpleUndirectedGraphDef}[def]}
&
\textbf{Keine Knotenschleifen}
\UebFormel{\begin{aligned}
&\SimpleUndirectedGraphStruct{V,E},\
 x\in V\ \vdash\\
&\qquad(x,x)\notin E.
\end{aligned}}
\UebQuellen{\FormulaRefAuto{SimpleUndirectedGraphLoopFree}[thm]} \\
\addlinespace[.8ex]

\textbf{Einschränkung auf eine Knotenmenge}
\UebFormel{\begin{aligned}
&U\subseteq V,\qquad E[U]=E\cap(U\times U).
\end{aligned}}
\UebQuellen{\FormulaRefAuto{DirectedInducedEdgeRelationDef}[def]}
&
\textbf{Induzierte Teilgraphen sind wieder schlicht}
\UebFormel{\begin{aligned}
&\SimpleUndirectedGraphStruct{V,E},\
 U\subseteq V\ \vdash\\
&\qquad\SimpleUndirectedGraphStruct{U,E[U]}.
\end{aligned}}
\UebQuellen{\FormulaRefAuto{SimpleUndirectedInducedSubgraph}[thm]} \\
\addlinespace[.8ex]
\end{BandUebersicht}


\chapter[Band 25: Zusammenhängende Graphen]{Band 25 auf einen Blick}
\noindent\textbf{Zusammenhängende Graphen}\par\medskip
Vorausgesetzt sind ungerichtete Graphen und die gerichtete Erreichbarkeit
mittels endlicher Walks. Weder Schleifenfreiheit noch ein nichtleerer Träger
werden zusätzlich in die Definition aufgenommen.

\begin{BandUebersicht}
\textbf{Zusammenhang}
\UebFormel{\begin{aligned}
&\ConnectedGraphStruct{V,E}\ \leftrightarrow\\
&\quad\UndirectedGraphStruct{V,E}\land\\
&\quad\forall u,v\in V\quad u\GraphReach E v.
\end{aligned}}
\UebQuellen{\FormulaRefAuto{ConnectedGraphDef}[def]}
&
\textbf{Zugriff auf die Grundstruktur}
\UebFormel{\begin{aligned}
&\ConnectedGraphStruct{V,E}\ \vdash\\
&\quad\UndirectedGraphStruct{V,E}\ \land\\
&\quad\forall u,v\in V\quad u\GraphReach E v.
\end{aligned}}
\UebQuellen{\FormulaRefAuto{ConnectedGraphAccess}[thm]} \\
\addlinespace[.8ex]

\textbf{Verbindung konkreter Endpunkte}
Erreichbarkeit erlaubt gleiche Endpunkte ausdrücklich:\UebFormel{\begin{aligned}
&u,v\in V,\quad u=v\ \lor\\
&\exists n\in\mathbb N\,\exists p\
 (\DirectedWalk V E p n\\
&\hspace{6em}\land p(0)=u\land p(n)=v).
\end{aligned}}
\UebQuellen{\FormulaRefAuto{DirectedReachabilityDef}[def]}
&
\textbf{Beliebige Knoten sind erreichbar}
\UebFormel{\begin{aligned}
&\ConnectedGraphStruct{V,E},\
 u,v\in V\ \vdash\\
&\qquad u\GraphReach E v.
\end{aligned}}
\UebQuellen{\FormulaRefAuto{ConnectedGraphVertexReachability}[thm]} \\
\addlinespace[.8ex]
\end{BandUebersicht}


\chapter[Band 26: Bäume]{Band 26 auf einen Blick}
\noindent\textbf{Bäume}\par\medskip
Vorausgesetzt sind die vier Graphbände und endliche Folgen.
\(\mathcal P_{V,E}(u,v)\) besteht aus Paaren aus Kantenlänge und
injektiver Knotenfolge mit Endpunkten \(u,v\). Bäume müssen nicht endlich sein.

\begin{BandUebersicht}
\textbf{Endliche einfache Verbindungswege}
\UebFormel{\begin{aligned}
&\mathcal P_{V,E}(u,v)=\\
&\{(n,p)\mid n\in\mathbb N,\\
&\quad p\in\powerset(\mathbb N\times V),\\
&\quad\DirectedPath V E p n,\\
&\quad p(0)=u,\ p(n)=v\}.
\end{aligned}}
\UebQuellen{\FormulaRefAuto{FiniteSimplePathFamilyDef}[def]}
&
\textbf{Folgen und Endpunkte sind typisiert}
\UebFormel{\begin{aligned}
&(n,p)\in\mathcal P_{V,E}(u,v),\
 \GraphStruct{V,E}\ \vdash\\
&\quad p:\NatSeg{n+1}\to V,\qquad u,v\in V.
\end{aligned}}
\UebQuellen{\FormulaRefAuto{FiniteSimplePathSequence}[thm];
\FormulaRefAuto{FiniteSimplePathEndpointTyping}[thm]} \\
\addlinespace[.8ex]

\textbf{Wald und Baum}
Ein Wald ist schlicht und besitzt höchstens einen einfachen
Weg zwischen je zwei Knoten. Ein Baum ist schlicht, nichtleer und erfüllt\UebFormel{\begin{aligned}
&\forall u,v\in V\
 \exists!z\ (z\in\mathcal P_{V,E}(u,v)).
\end{aligned}}
\UebQuellen{\FormulaRefAuto{ForestDef}[def];
\FormulaRefAuto{TreeDef}[def]}
&
\textbf{Zusammenhang und Wegeindeutigkeit}
\UebFormel{\begin{aligned}
&\TreeStruct{V,E}\ \vdash\
 \ConnectedGraphStruct{V,E},\\
&\TreeStruct{V,E},\ u,v\in V,\\
&a,b\in\mathcal P_{V,E}(u,v)\ \vdash\ a=b.
\end{aligned}}
\UebQuellen{\FormulaRefAuto{TreeConnected}[thm];
\FormulaRefAuto{TreePathUnique}[thm]} \\
\addlinespace[.8ex]

\textbf{Verwurzelter Baum}
\UebFormel{\begin{aligned}
&\RootedTreeStruct{V,E,r}\ \leftrightarrow\\
&\qquad\TreeStruct{V,E}\land r\in V.
\end{aligned}}
\UebQuellen{\FormulaRefAuto{RootedTreeDef}[def]}
&
\textbf{Eindeutige Wurzelwege genügen}
\UebFormel{\begin{aligned}
&\RootedTreeStruct{V,E,r}\ \eqvdash\\
&\quad\SimpleUndirectedGraphStruct{V,E}\land r\in V\\
&\quad{}\land\forall x\in V\
 \exists!z\ (z\in\mathcal P_{V,E}(r,x)).
\end{aligned}}
\UebQuellen{\FormulaRefAuto{RootedTreeIffUniqueRootPaths}[thm]} \\
\addlinespace[.8ex]

\textbf{Eltern und Kinder}
Im verwurzelten Baum ist \(a\) Elternknoten von \(x\ne r\),
wenn es auf dessen Wurzelweg unmittelbar vor \(x\) steht:\UebFormel{\begin{aligned}
&\exists k\in\mathbb N\,\exists p\\
&\quad((k+1,p)\in\mathcal P_{V,E}(r,x)\\
&\qquad\land p(k)=a).
\end{aligned}}Kinder sind durch die umgekehrte Elternrelation definiert.
\UebQuellen{\FormulaRefAuto{TreeParentDef}[def];
\FormulaRefAuto{TreeChildDef}[def];
\FormulaRefAuto{TreeChildrenDef}[def]}
&
\textbf{Genau ein Elternknoten außerhalb der Wurzel}
\UebFormel{\begin{aligned}
&\RootedTreeStruct{V,E,r},\
 x\in V\setminus\{r\}\ \vdash\\
&\qquad\exists!a\in V\
 \TreeParentRel V E r a x.
\end{aligned}}Die Wurzel hat keinen Elternknoten.
\UebQuellen{\FormulaRefAuto{TreeParentUnique}[thm];
\FormulaRefAuto{RootHasNoParent}[thm]} \\
\addlinespace[.8ex]

\textbf{Blätter, innere Knoten und Vorfahren}
Im verwurzelten Baum sei \(K(x)=\TreeChildren V E r x\).\UebFormel{\begin{aligned}
&\operatorname{Blatt}(x)\leftrightarrow
      (x\in V\land K(x)=\varnothing),\\
&\operatorname{Innen}(x)\leftrightarrow
      (x\in V\land K(x)\ne\varnothing).
\end{aligned}}Ein Vorfahr liegt auf dem Wurzelweg zum betrachteten Knoten.
\UebQuellen{\FormulaRefAuto{TreeLeafInnerDef}[def];
\FormulaRefAuto{TreeAncestorDef}[def]}
&
\textbf{Zerlegung und ausgezeichnete Vorfahren}
Für \(\RootedTreeStruct{V,E,r}\) und \(x\in V\):\UebFormel{\begin{aligned}
&\operatorname{Blatt}(x)\lor\operatorname{Innen}(x),\\
&\operatorname{Blatt}(x)\ \vdash\neg\operatorname{Innen}(x),\\
&\TreeAncestorRel V E r x x,\\
&\TreeAncestorRel V E r r x.
\end{aligned}}
\UebQuellen{\FormulaRefAuto{TreeLeafInnerPartition}[thm];
\FormulaRefAuto{TreeAncestorReflexive}[thm];
\FormulaRefAuto{RootAncestorEveryVertex}[thm]} \\
\addlinespace[.8ex]

\textbf{Binäre und geordnete volle Binärbäume}
Binär bedeutet: höchstens zwei Kinder je Knoten. Voll bedeutet:
null oder genau zwei Kinder. Eine Ordnung benennt auf der Menge \(I\)
der inneren Knoten linkes und rechtes Kind:\UebFormel{\begin{aligned}
&\lambda,\rho:I\to V,\quad
 \lambda(x)\ne\rho(x),\\
&K(x)=\{\lambda(x),\rho(x)\}\quad(x\in I).
\end{aligned}}
\UebQuellen{\FormulaRefAuto{BinaryTreeDef}[def];
\FormulaRefAuto{FullBinaryTreeDef}[def];
\FormulaRefAuto{OrderedFullBinaryTreeDef}[def]}
&
\textbf{Innere Knoten haben zwei verschiedene Kinder}
\UebFormel{\begin{aligned}
&\FullBinaryTreeStruct{V,E,r},\
 \TreeInnerPred V E r x\ \vdash\\
&\quad\exists a,b\in K(x)\
 (a\ne b\land K(x)=\{a,b\}).
\end{aligned}}Vergessen von \(\lambda,\rho\) liefert einen vollen Binärbaum,
weiter einen binären und einen verwurzelten Baum.
\UebQuellen{\FormulaRefAuto{FullBinaryInnerHasTwoChildren}[thm];
\FormulaRefAuto{OrderedFullBinaryTreeForgetfulChain}[thm]} \\
\addlinespace[.8ex]
\end{BandUebersicht}

\chapter[Band 27: Endliche Wörter und Klammerungsbäume]{Band 27 auf einen Blick}
\noindent\textbf{Endliche Wörter und Klammerungsbäume}\par\medskip

% Dieselbe lokale Notation wie im Fachband, für die übernommene Übersicht.
\begingroup
% Abstand zur Tabellenlinie auch unter mehrzeiligen Formeln sichern.
\setlength{\aboverulesep}{6pt}
\newcommand{\BinaryAddresses}{\mathcal W_2}
\newcommand{\AddressPrefixMap}[1]{\operatorname{vor}_{#1}}
\newcommand{\AddressPrefix}[2]{\AddressPrefixMap{#1}(#2)}
\newcommand{\AddressGraft}[2]{\Gamma(#1,#2)}
\newcommand{\GoodAddressSet}[1]{\mathsf{VollPraef}(#1)}
\newcommand{\TreePositions}[2]{\operatorname{Pos}_{#1}(#2)}
\newcommand{\AddressEdges}[1]{E_{#1}^{\mathrm{adr}}}
\newcommand{\AddressInner}[1]{I_{#1}^{\mathrm{adr}}}
\newcommand{\AddressLeft}[1]{\lambda_{#1}^{\mathrm{adr}}}
\newcommand{\AddressRight}[1]{\rho_{#1}^{\mathrm{adr}}}
\newcommand{\TreePositionGraph}[2]{\mathcal G_{#1}^{\mathrm{adr}}(#2)}
\begingroup
\small
\setlength{\LTleft}{0pt}
\setlength{\LTright}{0pt}
\setlength{\LTpre}{.5\baselineskip}
\setlength{\LTpost}{0pt}
\renewcommand{\arraystretch}{1.08}
\newcommand{\BXXVIIBlickFormel}[1]{%
  \par\smallskip\noindent\(\displaystyle #1\)\par}
\newcommand{\BXXVIIBlickQuellen}[1]{%
  \par\nobreak\smallskip{\footnotesize #1}\par}

\noindent
\emph{Lesart.}
\(\NatSeg n=\{k\in\mathbb N\mid k<n\}\) ist der \(n\)-te Anfangsabschnitt.
Links von \(\vdash\) stehen die Voraussetzungen, rechts die Folgerung;
\(\eqvdash\) bezeichnet die entsprechende Äquivalenz.
\(\mathsf{BinOp}(A,\star)\) bedeutet, dass \(\star\) eine binäre Operation
auf \(A\) ist.

\begin{longtable}{@{}>{\raggedright\arraybackslash}p{.455\linewidth}%
  @{\hspace{.012\linewidth}\vrule width .35pt\hspace{.012\linewidth}}%
  >{\raggedright\arraybackslash}p{.51\linewidth}@{}}
\toprule
\textbf{Definitionen und Konstruktionen}&\textbf{Zentrale Resultate}\\
\midrule
\endfirsthead
\toprule
\textbf{Definitionen und Konstruktionen}&
  \textbf{Zentrale Resultate (Fortsetzung)}\\
\midrule
\endhead
\midrule
\multicolumn{2}{@{}r@{}}{%
  \footnotesize\itshape Fortsetzung auf der nächsten Seite}\\
\endfoot
\bottomrule
\endlastfoot

\multicolumn{2}{@{}l@{}}{\large\bfseries I. Endliche Wörter}\\*[.8ex]

\textbf{Wörter als kleinste abgeschlossene Menge}
\BXXVIIBlickFormel{%
  \begin{aligned}
  \mathcal E_A&\coloneqq\FinSubsets{\mathbb N\times A},\\[-2pt]
  J(w,a)&\coloneqq w\cup\{(\FiniteCard(w),a)\},\\[-2pt]
  \mathfrak C_A&\coloneqq
    \bigl\{C\in\powerset(\mathcal E_A)\bigm|{}\\[-2pt]
    &\quad\varnothing\in C\ \land{}\\[-2pt]
    &\quad\forall w\in C\,\forall a\in A:\\[-2pt]
    &\qquad J(w,a)\in C\bigr\},\\[-2pt]
  \FiniteWords A&\coloneqq\bigcap\mathfrak C_A.
  \end{aligned}}
Dabei sind \(w\in \mathcal E_A\) und \(a\in A\) die Voraussetzungen für
\(J(w,a)\). Der Umgebungsbereich \(\mathcal E_A\) enthält alle endlichen
Teilmengen; erst Anfang und Anfügung bestimmen die Wörter.
\BXXVIIBlickQuellen{%
  \FormulaRefAuto{FiniteOccurrenceSetUniverseDef}[def];
  \FormulaRefAuto{OccurrenceCardinalAdjunctionDef}[def];
  \FormulaRefAuto{WordGenerationFamilyDef}[def];
  \FormulaRefAuto{FiniteWordSetDef}[def]}
&
\textit{Bewiesener Anschluss an endliche Folgen}
\BXXVIIBlickFormel{%
  \begin{aligned}
  w\in\FiniteWords A\ \eqvdash\ {}&
    w\in\powerset(\mathbb N\times A)\ \land{}\\[-2pt]
  &\exists n\in\mathbb N\;(w\colon\NatSeg n\to A),\\[2pt]
  w\in\FiniteWords A\ \vdash\ {}&
    \exists!n\in\mathbb N\;(w\colon\NatSeg n\to A).
  \end{aligned}}
\BXXVIIBlickQuellen{%
  \FormulaRefAuto{FiniteWordSetMembership}[thm];
  \FormulaRefAuto{FiniteWordLengthUnique}[thm]}
Anfang und Anfügungsschritt liefern die Wortinduktion. Die lokale Regel
fasst beide Teilbeweise mit der Wortzugehörigkeit zusammen.
\BXXVIIBlickQuellen{\FormulaRefAuto{GeneratedWordInductionRule}[thm]}\\
\addlinespace[.8ex]

\textbf{Wortlänge}
\BXXVIIBlickFormel{%
  \begin{aligned}
  w\in\FiniteWords A\ \vdash\quad&\\[-7pt]
  \WordLength w\coloneqq{}&\FiniteCard(w).
  \end{aligned}}
Die Länge zählt die aufeinanderfolgenden Positionen und damit die
Anfügungsschritte. Grundlage ist die endliche Kardinalität aus Band~20.
\BXXVIIBlickQuellen{\FormulaRefAuto{FiniteWordLengthDef}[def];
  \FormulaRefAuto{FiniteWordLengthCardinality}[thm]}
&
\textit{Typisierung und Identifikation}
\BXXVIIBlickFormel{%
  \begin{aligned}
  w\in\FiniteWords A\ \vdash\quad
    &\WordLength w\in\mathbb N,\\[-2pt]
    &w\colon\NatSeg{\WordLength w}\to A,\\[2pt]
  n\in\mathbb N\dsep w\colon\NatSeg n\to A\ \vdash\quad
    &\WordLength w=n.
  \end{aligned}}
\BXXVIIBlickQuellen{%
  \FormulaRefAuto{FiniteWordTyping}[thm];
  \FormulaRefAuto{FiniteWordLengthFromTyping}[thm]}\\
\addlinespace[.8ex]

\textbf{Leerwort und Nichtleerwörter}
\BXXVIIBlickFormel{%
  \EmptyWord A\coloneqq\varnothing,
  \qquad
  \NonemptyWords A\coloneqq
    \FiniteWords A\setminus\{\EmptyWord A\}.}
\BXXVIIBlickQuellen{%
  \FormulaRefAuto{EmptyWordDef}[def];
  \FormulaRefAuto{NonemptyWordSetDef}[def]}
&
\textit{Länge und Elementkriterium}
\BXXVIIBlickFormel{%
  \begin{aligned}
  &\EmptyWord A\in\FiniteWords A,
    \qquad \WordLength{\EmptyWord A}=0,\\[2pt]
  &w\in\NonemptyWords A\ \eqvdash\quad
    w\in\FiniteWords A\land\WordLength w\neq0.
  \end{aligned}}
\BXXVIIBlickQuellen{%
  \FormulaRefAuto{EmptyWordLength}[thm];
  \FormulaRefAuto{NonemptyWordMembership}[thm]}\\
\addlinespace[.8ex]

\textbf{Buchstabenwort}
\BXXVIIBlickFormel{%
  a\in A\ \vdash\quad
  \LetterWord a\coloneqq\{(0,a)\}.}
\BXXVIIBlickQuellen{\FormulaRefAuto{LetterWordDef}[def]}
&
\textit{Grunddaten}
\BXXVIIBlickFormel{%
  a\in A\ \vdash\quad
  \begin{aligned}
  &\LetterWord a\colon\NatSeg 1\to A,\\[-2pt]
  &\LetterWord a\in\NonemptyWords A,\\[-2pt]
  &\WordLength{\LetterWord a}=1,
    \qquad \LetterWord a(0)=a.
  \end{aligned}}
\BXXVIIBlickQuellen{\FormulaRefAuto{LetterWordFacts}[thm]}\\
\midrule

\textbf{Primitive Anfügung eines Buchstabens}
\BXXVIIBlickFormel{%
  \sigma_A(u,a)=u\cup\{(\WordLength u,a)\}.}
Die Anfügung wird vor der allgemeinen Konkatenation eingeführt.
\BXXVIIBlickQuellen{\FormulaRefAuto{WordModelSuccessorDef}[def]}
&
\textit{Das konkrete Modell erfüllt W0 bis W3}
Die Minimalität folgt bereits aus der Konstruktion als kleinste
abgeschlossene Menge. Der Anschluss an endliche Folgen liefert
Typisierung, Trennung vom Leerwort und gemeinsame Injektivität.
\BXXVIIBlickQuellen{%
  \FormulaRefAuto{EarlyWordModelMinimality}[thm];
  \FormulaRefAuto{WordStructureConcreteModel}[thm]}\\
\addlinespace[.8ex]

\textbf{Frühe Wortaxiome}
\BXXVIIBlickFormel{%
  e\in W,\qquad s\colon W\times A\to W.}
W1: Der Anfang ist kein Nachfolger. W2: Die Anfügung ist gemeinsam
injektiv. W3: Jede Teilmenge mit Anfang und Anfügungsabschluss ist
der ganze Träger.
\BXXVIIBlickQuellen{\FormulaRefAuto{WordStructureDef}[def]}
&
\textit{Zerlegung und Induktion aus den Axiomen}
Jedes Wort ist der Anfang oder eindeutig ein Nachfolger.
Eine Eigenschaft, die für den Anfang gilt und bei Anfügung erhalten
bleibt, gilt für sämtliche Wörter der Struktur.
\BXXVIIBlickQuellen{%
  \FormulaRefAuto{WordStructureDecomposition}[thm];
  \FormulaRefAuto{WordStructureInduction}[thm];
  \FormulaRefAuto{WordStructureInductionRule}[thm]}
Die lokale Regel gilt in jeder Wortstruktur und damit auch im konkreten Modell.\\
\addlinespace[.8ex]

\textbf{Rekursionsdaten einer Wortstruktur}
\BXXVIIBlickFormel{%
  x_0\in X,\qquad r\colon X\times A\to X.}
Die Rekursionsfunktion wird als kleinster abgeschlossener Graph
unmittelbar aus W0 bis W3 gewonnen.
&
\textit{Axiomatische Rekursion und konkreter Modellfall}
\BXXVIIBlickFormel{%
  \begin{aligned}
  &\exists!F\colon W\to X:\\[-2pt]
  &F(e)=x_0,\\[-2pt]
  &F(s(u,a))=r(F(u),a).
  \end{aligned}}
\BXXVIIBlickQuellen{%
  \FormulaRefAuto{WordStructureRecursion}[thm];
  \FormulaRefAuto{FiniteWordRecursion}[thm]}\\
\midrule

\textbf{Konkatenation}
\BXXVIIBlickFormel{%
  \begin{aligned}
  u,v\in\FiniteWords A\ \vdash\quad&\\[-7pt]
  \WordConcat uv&\coloneqq u\cup\\[-2pt]
  &\quad\{(\WordLength u+j,v(j))\mid{}\\[-2pt]
  &\qquad j\in\NatSeg{\WordLength v}\}.
  \end{aligned}}
\BXXVIIBlickQuellen{\FormulaRefAuto{WordConcatenationDef}[def]}
&
\textit{Typisierung und Koordinaten}
\BXXVIIBlickFormel{%
  \begin{aligned}
  &\WordConcat uv\colon
    \NatSeg{\WordLength u+\WordLength v}\to A,\\[-2pt]
  &i\in\NatSeg{\WordLength u}\ \vdash\quad
    \WordConcat uv(i)=u(i),\\[-2pt]
  &j\in\NatSeg{\WordLength v}\ \vdash\quad
    \WordConcat uv(\WordLength u+j)=v(j).
  \end{aligned}}
\BXXVIIBlickQuellen{\FormulaRefAuto{WordConcatenationTyping}[thm]}\\
\addlinespace[.8ex]

\textit{Die zweite Wortkoordinate wird also um
\(\WordLength u\) verschoben.}
&
\textit{Länge und Rechengesetze}
\BXXVIIBlickFormel{%
  \begin{aligned}
  &\WordLength{\WordConcat uv}=\WordLength u+\WordLength v,\\[-2pt]
  &\WordConcat{\EmptyWord A}w=w=
    \WordConcat w{\EmptyWord A},\\[-2pt]
  &\WordConcat{\WordConcat uv}w=
    \WordConcat u{\WordConcat vw}.
  \end{aligned}}
\BXXVIIBlickQuellen{%
  \FormulaRefAuto{WordConcatenationLength}[thm];
  \FormulaRefAuto{WordConcatenationIdentity}[thm];
  \FormulaRefAuto{WordConcatenationAssociative}[thm]}
Die Rechengesetze werden mit der aus W3 gewonnenen Wortinduktion
bewiesen.\\
\addlinespace[.8ex]

\textit{Fester Schnitt}
\BXXVIIBlickFormel{%
  \WordLength u=\WordLength{u'}\dsep
  \WordConcat uv=\WordConcat{u'}{v'}.}
&
\textit{Kürzung am festgelegten Schnitt}
\BXXVIIBlickFormel{%
  u=u',\qquad v=v'.}
\BXXVIIBlickQuellen{%
  \FormulaRefAuto{WordConcatenationFixedCutCancellation}[thm]}\\
\midrule

\textbf{Konstruktoren nichtleerer Wörter}
\BXXVIIBlickFormel{%
  a\longmapsto\LetterWord a,
  \qquad
  (u,a)\longmapsto\WordConcat u{\LetterWord a}.}
&
\textit{Eindeutige Rechtszerlegung}
\BXXVIIBlickFormel{%
  \begin{aligned}
  w\in\NonemptyWords A\ \vdash\quad
  \exists!u\in\FiniteWords A\;
  \exists!a\in A\;{}\\[-2pt]
  w=\WordConcat u{\LetterWord a}.
  \end{aligned}}
\BXXVIIBlickQuellen{%
  \FormulaRefAuto{NonemptyWordRightDecomposition}[thm]}
\textit{Dabei darf der Präfix \(u\) das Leerwort sein.}\\
\addlinespace[.8ex]

\textbf{Rumpf und Endbuchstabe}
\BXXVIIBlickFormel{%
  \begin{aligned}
  &\WordRumpfMap A\colon\NonemptyWords A\to\FiniteWords A,\\[-2pt]
  &\WordEndMap A\colon\NonemptyWords A\to A.
  \end{aligned}}
\BXXVIIBlickQuellen{%
  \FormulaRefAuto{WordRumpfDef}[def];
  \FormulaRefAuto{WordEndDef}[def]}
&
\textit{Zerlegen und wieder zusammensetzen}
\BXXVIIBlickFormel{%
  \begin{aligned}
  &w\in\NonemptyWords A\ \vdash\\[-2pt]
  &\quad w=\WordConcat{\WordRumpf A w}{\LetterWord{\WordEnd A w}},\\[2pt]
  &u\in\FiniteWords A\dsep a\in A\ \vdash\\[-2pt]
  &\quad\WordRumpf A{\WordConcat u{\LetterWord a}}=u,\\[-2pt]
  &\quad\WordEnd A{\WordConcat u{\LetterWord a}}=a.
  \end{aligned}}
\BXXVIIBlickQuellen{%
  \FormulaRefAuto{WordRightDecompositionEquation}[thm];
  \FormulaRefAuto{WordRumpfAppend}[thm];
  \FormulaRefAuto{WordEndAppend}[thm]}\\
\addlinespace[.8ex]

\textbf{Induktionsdaten}
\BXXVIIBlickFormel{%
  \begin{aligned}
  &\forall a\in A\;P(\LetterWord a),\\[-2pt]
  &\forall u\in\NonemptyWords A\;\forall a\in A:\\[-2pt]
  &\qquad P(u)\to P(\WordConcat u{\LetterWord a}).
  \end{aligned}}
&
\textit{Induktion über nichtleere Wörter}
\BXXVIIBlickFormel{%
  \forall w\in\NonemptyWords A\;P(w).}
\BXXVIIBlickQuellen{\FormulaRefAuto{NonemptyWordInduction}[thm]}\\
\addlinespace[.8ex]

\textbf{Rekursionsdaten}
\BXXVIIBlickFormel{%
  f\colon A\to X,
  \qquad s\colon X\times A\to X.}
&
\textit{Rekursion über nichtleere Wörter}
\BXXVIIBlickFormel{%
  \begin{aligned}
  &\exists!F\colon\NonemptyWords A\to X:\\[-2pt]
  &F(\LetterWord a)=f(a),\\[-2pt]
  &F(\WordConcat u{\LetterWord a})=s(F(u),a).
  \end{aligned}}
\BXXVIIBlickQuellen{\FormulaRefAuto{NonemptyWordRecursion}[thm]}
Sie folgt aus der zuvor bewiesenen Rekursion für alle Wörter.\\
\addlinespace[.8ex]

\textbf{Linksfaltung}
\BXXVIIBlickFormel{%
  \begin{aligned}
  &\mathsf{BinOp}(A,\star)\ \vdash\\[-2pt]
  &\Pi_\star\coloneqq
    \iota F\colon\NonemptyWords A\to A,\\[-2pt]
  &\forall a\in A\quad F(\LetterWord a)=a,\\[-2pt]
  &\forall u\in\NonemptyWords A\;\forall a\in A:\\[-2pt]
  &\qquad F(\WordConcat u{\LetterWord a})=F(u)\star a.
  \end{aligned}}
\BXXVIIBlickQuellen{\FormulaRefAuto{LeftWordFoldDef}[def]}
&
\textit{Rekursionsgleichungen}
\BXXVIIBlickFormel{%
  \begin{aligned}
  &\WordFold\star{\LetterWord a}=a,\\[-2pt]
  &\WordFold\star{\WordConcat u{\LetterWord a}}
    =\WordFold\star u\star a.
  \end{aligned}}
\BXXVIIBlickQuellen{%
  \FormulaRefAuto{LeftWordFoldFunction}[thm];
  \FormulaRefAuto{LeftWordFoldEquations}[thm]}
\textit{Ohne ausgezeichnetes neutrales Element wird das Leerwort nicht
gefaltet.}\\
\midrule

\multicolumn{2}{@{}l@{}}{%
  \large\bfseries II. Wortcodes und ihre vollen planaren Positionsbäume}\\*[.8ex]

\textbf{Baumalphabet und Blattcode}
\BXXVIIBlickFormel{%
  \begin{aligned}
  \TreeAlphabet A&\coloneqq
    (\{0\}\times A)\cup\\[-2pt]
    &\qquad(\{1\}\times\mathbb N),\\[-2pt]
  a\in A\ \vdash\quad
  \LeafTree A a&\coloneqq\LetterWord{(0,a)}.
  \end{aligned}}
\BXXVIIBlickQuellen{%
  \FormulaRefAuto{TreeTokenAlphabetDef}[def];
  \FormulaRefAuto{TreeLeafCodeDef}[def]}
&
\textit{Blattinjektivität}
\BXXVIIBlickFormel{%
  a,b\in A\ \vdash\quad
  \bigl(\LeafTree A a=\LeafTree A b\ \leftrightarrow\ a=b\bigr).}
\BXXVIIBlickQuellen{\FormulaRefAuto{TreeConstructorNoConfusion}[thm]}\\
\addlinespace[.8ex]

\textbf{Knotencode}
\BXXVIIBlickFormel{%
  \begin{aligned}
  &S,T\in\NonemptyWords{\TreeAlphabet A}\ \vdash\\[-2pt]
  &\NodeTree A S T\coloneqq
    \LetterWord{(1,\WordLength S)}\frown\\[-2pt]
  &\hspace{7em}(\WordConcat ST).
  \end{aligned}}
\BXXVIIBlickQuellen{\FormulaRefAuto{TreeNodeCodeDef}[def]}
&
\textit{Konstruktortrennung und Knoteninjektivität}
\BXXVIIBlickFormel{%
  \begin{aligned}
  &\LeafTree A a\neq\NodeTree A S T,\\[-2pt]
  &\NodeTree A S T=\NodeTree A{S'}{T'}\ \vdash\\[-2pt]
  &\qquad S=S',\quad T=T'.
  \end{aligned}}
\BXXVIIBlickQuellen{\FormulaRefAuto{TreeConstructorNoConfusion}[thm]}
\textit{Die gespeicherte Länge \(\WordLength S\) legt den Schnitt des
Knotencodes fest.}\\
\midrule

\textbf{Erzeugungsschritt}
\BXXVIIBlickFormel{%
  \begin{aligned}
  &\TreeClosure A X\coloneqq\\[-2pt]
  &\quad\{\LeafTree A a\mid a\in A\}\cup\\[-2pt]
  &\quad\{\NodeTree A S T\mid S,T\in X\}.
  \end{aligned}}
\BXXVIIBlickQuellen{\FormulaRefAuto{TreeGenerationStepDef}[def]}
&
\textit{Abbildung und Monotonie}
\BXXVIIBlickFormel{%
  \begin{aligned}
  &\mathcal C_A\colon
    \powerset(\NonemptyWords{\TreeAlphabet A})
    \to\powerset(\NonemptyWords{\TreeAlphabet A}),\\[-2pt]
  &X\subseteq Y\ \vdash\quad
    \TreeClosure A X\subseteq\TreeClosure A Y.
  \end{aligned}}
\BXXVIIBlickQuellen{%
  \FormulaRefAuto{TreeGenerationStepFunction}[thm];
  \FormulaRefAuto{TreeGenerationStepMonotonicity}[thm]}\\
\addlinespace[.8ex]

\textbf{Baumstufen}
\BXXVIIBlickFormel{%
  n\in\mathbb N\ \vdash\quad
  \TreeStage A n\coloneqq
  \RecFun{\varnothing}{\mathcal C_A}(n).}
\BXXVIIBlickQuellen{\FormulaRefAuto{TreeStageDef}[def]}
&
\textit{Stufengleichungen und Verschachtelung}
\BXXVIIBlickFormel{%
  \begin{aligned}
  &\TreeStage A0=\varnothing,\\[-2pt]
  &n\in\mathbb N\ \vdash\quad
    \TreeStage A{n+1}=\TreeClosure A{\TreeStage A n},\\[-2pt]
  &m,n\in\mathbb N\dsep m\leq n\ \vdash\quad
    \TreeStage A m\subseteq\TreeStage A n.
  \end{aligned}}
\BXXVIIBlickQuellen{%
  \FormulaRefAuto{TreeStageZeroEquation}[thm];
  \FormulaRefAuto{TreeStageSuccessorEquation}[thm];
  \FormulaRefAuto{TreeStageNested}[thm]}\\
\addlinespace[.8ex]

\textbf{Menge der Klammerungsbäume}
\BXXVIIBlickFormel{%
  \begin{aligned}
  \BracketTreeSet A\coloneqq
  \{R\in\NonemptyWords{\TreeAlphabet A}\mid{}&
    \exists n\in\mathbb N\;{}\\[-2pt]
    &R\in\TreeStage A n\}.
  \end{aligned}}
\BXXVIIBlickQuellen{\FormulaRefAuto{FullPlanarBinaryTreeSetDef}[def]}
&
\textit{Das konkrete Modell erfüllt B0 bis B4}
\BXXVIIBlickFormel{%
  \begin{aligned}
  &\lambda_A\colon A\to\BracketTreeSet A,\\[-2pt]
  &\kappa_A\colon\BracketTreeSet A\times\BracketTreeSet A
      \to\BracketTreeSet A.
  \end{aligned}}
Die Konstruktoren sind injektiv und haben getrennte Bilder.
Die Minimalität folgt durch natürliche Induktion über die Baumstufen.
\BXXVIIBlickQuellen{%
  \FormulaRefAuto{TreeCodeMinimality}[thm];
  \FormulaRefAuto{BinaryTreeStructureConcreteModel}[thm]}\\
\addlinespace[.8ex]

\textbf{Frühe Baumaxiome}
\BXXVIIBlickFormel{%
  \ell\colon A\to T,\qquad k\colon T\times T\to T.}
B1 und B2: Beide Konstruktoren sind injektiv.
B3: Ihre Bilder sind getrennt.
B4: Jede Teilmenge mit allen Blättern und Knotenabschluss ist
der ganze Träger.
\BXXVIIBlickQuellen{\FormulaRefAuto{BinaryTreeStructureDef}[def]}
&
\textit{Zerlegung und Induktion aus den Axiomen}
Jeder Baum ist eindeutig ein beschriftetes Blatt oder ein Knoten
mit zwei geordneten Teilbäumen. Die Minimalität liefert die
Strukturinduktion für jede solche Baumstruktur.
\BXXVIIBlickQuellen{%
  \FormulaRefAuto{BinaryTreeStructureInduction}[thm];
  \FormulaRefAuto{TreeConstructorDecomposition}[thm]}\\
\midrule

\textbf{Daten der Strukturinduktion}
\BXXVIIBlickFormel{%
  \begin{aligned}
  &\forall a\in A\;P(\LeafTree A a),\\[-2pt]
  &\forall S,T\in\BracketTreeSet A:\\[-2pt]
  &\qquad P(S)\land P(T)\to\\[-2pt]
  &\qquad P(\NodeTree A S T).
  \end{aligned}}
&
\textit{Strukturinduktion}
\BXXVIIBlickFormel{%
  \forall R\in\BracketTreeSet A\;P(R).}
\BXXVIIBlickQuellen{%
  \FormulaRefAuto{FullPlanarBinaryTreeStructuralInduction}[thm]}\\
\addlinespace[.8ex]

\textbf{Daten der Strukturrekursion}
\BXXVIIBlickFormel{%
  f\colon A\to X,
  \qquad g\colon X\times X\to X.}
&
\textit{Axiomatische Rekursion und konkreter Modellfall}
\BXXVIIBlickFormel{%
  \begin{aligned}
  &\exists!E\colon\BracketTreeSet A\to X:\\[-2pt]
  &E(\LeafTree A a)=f(a),\\[-2pt]
  &E(\NodeTree A S T)=g(E(S),E(T)).
  \end{aligned}}
\BXXVIIBlickQuellen{%
  \FormulaRefAuto{BinaryTreeStructureRecursion}[thm];
  \FormulaRefAuto{FullPlanarBinaryTreeRecursion}[thm]}
Der allgemeine Satz folgt direkt aus B0 bis B4 durch einen kleinsten
abgeschlossenen Graphen; der angezeigte Satz ist seine Anwendung
auf das konkrete Modell.\\
\addlinespace[.8ex]

\textbf{Der benannte Baumrekursor}
\BXXVIIBlickFormel{%
  \begin{aligned}
  &f\colon A\to X,\quad g\colon X\times X\to X,\\[-2pt]
  &R\coloneqq\TreeRecursor A X f g.
  \end{aligned}}
\textit{Blätter durch Werte ersetzen, Teilbaumwerte verknüpfen.}
\BXXVIIBlickQuellen{\FormulaRefAuto{TreeRecursorDef}[def]}
&
\textit{Einheitliche Rechen- und Vergleichsregeln}
\BXXVIIBlickFormel{%
  \begin{aligned}
  &R\colon\BracketTreeSet A\to X,\\[-2pt]
  &a\in A\ \vdash\ R(\LeafTree A a)=f(a),\\[-2pt]
  &S,T\in\BracketTreeSet A\ \vdash\\[-2pt]
  &\quad R(\NodeTree A S T)=g(R(S),R(T)).
  \end{aligned}}
\textit{Eine gleich typisierte Funktion mit denselben Blatt- und
Knotenregeln ist gleich \(R\).}
\BXXVIIBlickQuellen{%
  \FormulaRefAuto{TreeRecursorEquations}[thm];
  \FormulaRefAuto{TreeRecursorIdentification}[thm]}\\
\addlinespace[.8ex]

\textbf{Positionsgraph eines Baumcodes}
\BXXVIIBlickFormel{%
  \begin{aligned}
  \BinaryAddresses&\coloneqq\FiniteWords{\{0,1\}},\\[-2pt]
  P_R&\coloneqq\TreePositions A R,\\[-2pt]
  \TreePositionGraph A R&\coloneqq
    \bigl(P_R,\AddressEdges{P_R},
      \EmptyWord{\{0,1\}},\\[-2pt]
  &\hspace{6em}\AddressLeft{P_R},\AddressRight{P_R}\bigr).
  \end{aligned}}
\BXXVIIBlickQuellen{%
  \FormulaRefAuto{TreePositionGraphDef}[def]}
&
\textit{Der Wortcode trägt einen besonderen Baum}
\BXXVIIBlickFormel{%
  \begin{aligned}
  P_R&\coloneqq\TreePositions A R,\\[-2pt]
  R\in\BracketTreeSet A\ \vdash\quad
    &\Finite{P_R},\\[-2pt]
    &\TreeStruct{P_R,\AddressEdges{P_R}}.
  \end{aligned}}
\BXXVIIBlickQuellen{%
  \FormulaRefAuto{BracketCodePositionGraphIsFiniteOrderedFullBinaryTree}[thm]}
\textit{Der Positionsgraph ist sogar ein geordneter voller Binärbaum.
Seine Knoten sind Adressen; daher bleiben auch gleiche Teilbaumcodes an
verschiedenen Positionen verschieden.}\\
\midrule

\multicolumn{2}{@{}l@{}}{%
  \large\bfseries III. Blattwörter, Klammerungen und Auswertung}\\*[.8ex]

\textbf{Einbuchstabenabbildung und Blattwort}
\BXXVIIBlickFormel{%
  \begin{aligned}
  &\operatorname{ein}_A\colon A\to\NonemptyWords A,
    \qquad \operatorname{ein}_A(a)\coloneqq\LetterWord a,\\[2pt]
  &\operatorname{wort}_A\coloneqq\\[-2pt]
  &\quad\TreeRecursor A{\NonemptyWords A}{\operatorname{ein}_A}
    {\BinaryOpFunction{\NonemptyWords A}{\frown}}.
  \end{aligned}}
\BXXVIIBlickQuellen{%
  \FormulaRefAuto{LetterWordFunctionDef}[def];
  \FormulaRefAuto{TreeLeafWordDef}[def]}
&
\textit{Rekursionsgleichungen des Blattwortes}
\BXXVIIBlickFormel{%
  \begin{aligned}
  &\operatorname{wort}_A\colon
    \BracketTreeSet A\to\NonemptyWords A,\\[-2pt]
  &\TreeLeafWord A{\LeafTree A a}=\LetterWord a,\\[-2pt]
  &\TreeLeafWord A{\NodeTree A S T}=\\[-2pt]
  &\qquad\WordConcat{\TreeLeafWord A S}{\TreeLeafWord A T}.
  \end{aligned}}
\BXXVIIBlickQuellen{\FormulaRefAuto{TreeLeafWordEquations}[thm]}\\
\addlinespace[.8ex]

\textbf{Ein Baum klammert ein Wort}
\BXXVIIBlickFormel{%
  \TreeBrackets A T w\coloneqq
    (T,w)\in\operatorname{wort}_A.}
\BXXVIIBlickQuellen{\FormulaRefAuto{TreeBracketingRelationDef}[def]}
\textit{Die Relation enthält zugleich die Baum- und Worttypisierung.}
&
\textit{Blattregel und Knotenregel}
\BXXVIIBlickFormel{%
  \begin{aligned}
  &a\in A\ \vdash\\[-2pt]
  &\quad\TreeBrackets A{\LeafTree A a}{\LetterWord a},\\[2pt]
  &\TreeBrackets A S u\dsep\TreeBrackets A T v\ \vdash\\[-2pt]
  &\quad\TreeBrackets A{\NodeTree A S T}{\WordConcat u v}.
  \end{aligned}}
\BXXVIIBlickQuellen{%
  \FormulaRefAuto{TreeBracketingLeafIntroduction}[thm];
  \FormulaRefAuto{TreeBracketingNodeIntroduction}[thm]}\\
\addlinespace[.8ex]

\textbf{Ein Blatt rechts anfügen}
\BXXVIIBlickFormel{%
  \begin{aligned}
  &T\in\BracketTreeSet A\dsep a\in A\ \vdash\\[-2pt]
  &\TreeAppendLeaf A T a\coloneqq\\[-2pt]
  &\qquad\NodeTree A T{\LeafTree A a}.
  \end{aligned}}
\BXXVIIBlickQuellen{\FormulaRefAuto{TreeAppendLeafDef}[def]}
&
\textit{Baum und Wort gemeinsam erweitern}
\BXXVIIBlickFormel{%
  \begin{aligned}
  &\TreeBrackets A T u\dsep a\in A\ \vdash\\[-2pt]
  &\quad\TreeBrackets A{\TreeAppendLeaf A T a}
    {\WordConcat u{\LetterWord a}},\\[2pt]
  &w\in\NonemptyWords A\ \vdash
    \exists T\,(\TreeBrackets A T w).
  \end{aligned}}
\BXXVIIBlickQuellen{%
  \FormulaRefAuto{TreeBracketingAppendLeafIntroduction}[thm];
  \FormulaRefAuto{TreeBracketingExistence}[thm]}\\
\addlinespace[.8ex]

\textbf{Der Linksbaum eines Wortes}
\BXXVIIBlickFormel{%
  \begin{aligned}
  &\LeftBracketingMap A\colon\NonemptyWords A\to\BracketTreeSet A,\\[-2pt]
  &a\in A\ \vdash\\[-2pt]
  &\quad\LeftBracketing A{\LetterWord a}=\LeafTree A a,\\[-2pt]
  &u\in\NonemptyWords A\dsep a\in A\ \vdash\\[-2pt]
  &\quad\LeftBracketing A{\WordConcat u{\LetterWord a}}\\[-2pt]
  &\qquad=\TreeAppendLeaf A{\LeftBracketing A u}a.
  \end{aligned}}
\BXXVIIBlickQuellen{%
  \FormulaRefAuto{LeftBracketingDef}[def];
  \FormulaRefAuto{LeftBracketingLetterEquation}[thm];
  \FormulaRefAuto{LeftBracketingAppendEquation}[thm]}
&
\textit{Eine konkrete Klammerung und ihre Auswertung}
\BXXVIIBlickFormel{%
  \begin{aligned}
  &w\in\NonemptyWords A\ \vdash\\[-2pt]
  &\quad\TreeBrackets A{\LeftBracketing A w}w,\\[2pt]
  &\mathsf{BinOp}(A,\star)\dsep w\in\NonemptyWords A\ \vdash\\[-2pt]
  &\quad\TreeEvaluation\star{\LeftBracketing A w}=\WordFold\star w.
  \end{aligned}}
\textit{Den Linksbaum auswerten ergibt die Linksfaltung.}
\BXXVIIBlickQuellen{%
  \FormulaRefAuto{LeftBracketingProperty}[thm];
  \FormulaRefAuto{LeftBracketingEvaluation}[thm]}\\
\addlinespace[.8ex]

\textbf{Klammerungsmenge eines Nichtleerwortes}
\BXXVIIBlickFormel{%
  \begin{aligned}
  &w\in\NonemptyWords A\ \vdash\\[-2pt]
  &\WordBracketings A w\coloneqq\\[-2pt]
  &\quad\{T\in\BracketTreeSet A\mid
    \TreeLeafWord A T=w\}.
  \end{aligned}}
\BXXVIIBlickQuellen{\FormulaRefAuto{WordBracketingSetDef}[def]}
&
\textit{Elementkriterium und Existenz}
\BXXVIIBlickFormel{%
  \begin{aligned}
  &w\in\NonemptyWords A\ \vdash\\[-2pt]
  &\bigl(T\in\WordBracketings A w\ \leftrightarrow\\[-2pt]
  &\quad(T\in\BracketTreeSet A\ \land
    \TreeLeafWord A T=w)\bigr),\\[2pt]
  &w\in\NonemptyWords A\ \vdash
    \WordBracketings A w\neq\varnothing.
  \end{aligned}}
\BXXVIIBlickQuellen{%
  \FormulaRefAuto{WordBracketingMembership}[thm];
  \FormulaRefAuto{WordBracketingExistence}[thm]}\\
\addlinespace[.8ex]

\textbf{Auswertung eines Klammerungsbaums}
\BXXVIIBlickFormel{%
  \begin{aligned}
  &\mathsf{BinOp}(A,\star)\ \vdash\\[-2pt]
  &\operatorname{ev}_\star\coloneqq\\[-2pt]
  &\quad\TreeRecursor A A{\Id_A}{\BinaryOpFunction A\star}.
  \end{aligned}}
\BXXVIIBlickQuellen{\FormulaRefAuto{TreeEvaluationDef}[def]}
&
\textit{Typisierung und Rekursionsgleichungen}
\BXXVIIBlickFormel{%
  \begin{aligned}
  &\operatorname{ev}_\star\colon\BracketTreeSet A\to A,\\[-2pt]
  &\TreeEvaluation\star{\LeafTree A a}=a,\\[-2pt]
  &\TreeEvaluation\star{\NodeTree A S T}=
    \TreeEvaluation\star S\star\TreeEvaluation\star T.
  \end{aligned}}
\BXXVIIBlickQuellen{\FormulaRefAuto{TreeEvaluationEquations}[thm]}
\textit{Erst die Assoziativität in Band~28 macht die Auswertung von der
gewählten Klammerung unabhängig.}\\

\end{longtable}
\endgroup

\section*{Vom konkreten Modell zu frühen Strukturaxiomen}

Das Anfügungsmodell beginnt im Bereich aller endlichen Teilmengen von
\(\mathbb N\times A\). Die Wortmenge wird als Schnitt aller Teilmengen
dieses Bereichs definiert, die \(\varnothing\) enthalten und unter
\(u\mapsto u\cup\{(\FiniteCard(u),a)\}\) für \(a\in A\) abgeschlossen sind
(\FormulaRefAuto{FiniteWordSetDef}[def]). Wie beim Aufbau der natürlichen
Zahlen durch eine kleinste induktive Menge folgt die Minimalität
unmittelbar aus dieser Konstruktion. Die endlichen Teilmengen und ihre
Kardinalzahlen liefert Band~20; eine Wortmenge wird dabei noch nicht
vorausgesetzt.

Anschließend wird bewiesen, dass genau die nullbasierten endlichen Folgen
aus Band~21 entstehen. Dieser Anschluss begründet die Positionsdarstellung
und die eindeutige Zerlegung. Auf der Wortmenge wird der Erzeugungsschritt
als Anfügungsfunktion \(\sigma_A\) gefasst, noch vor der Konkatenation.
\FormulaRefAuto{WordStructureDef}[def] bündelt die so bewiesenen
Eigenschaften als W0 bis W3 für beliebige Wortstrukturen.

Beim Baumaufbau werden zunächst die Konstruktoren und die endlichen
Erzeugungsstufen untersucht. Natürliche Induktion über diese Stufen
beweist die Minimalität des konkreten Baumträgers
(\FormulaRefAuto{TreeCodeMinimality}[thm]).
Danach führt \FormulaRefAuto{BinaryTreeStructureDef}[def]
die Bedingungen B0 bis B4 ein: typisierte, injektive Blatt- und
Knotenkonstruktoren mit getrennten Bildern sowie Minimalität.
Die beiden Axiomsysteme sind Bedingungen an Strukturen; ihre bereits
konstruierten Modelle erfordern keine zusätzlichen Mengenexistenzaxiome.

\widowpenalty=10000
Aus den frühen Axiomen folgen eindeutige Zerlegung und Induktion.
Die Rekursion wird unmittelbar durch den Schnitt aller unter den
vorgegebenen Rekursionsregeln abgeschlossenen Relationsgraphen bewiesen;
siehe \FormulaRefAuto{WordStructureRecursion}[thm] und
\FormulaRefAuto{BinaryTreeStructureRecursion}[thm].
Erst daraus folgen die konkreten Rekursionssätze und die Eindeutigkeit
der Strukturen bis auf einen eindeutig bestimmten strukturtreuen
Isomorphismus. Die späteren Wortzerlegungen, Induktionsbeweise,
Konkatenationsgesetze und Auswertungen greifen auf diese Ergebnisse
zurück. Die konkrete Darstellung kann deshalb gewechselt werden,
ohne die Aufbauregeln zu verändern.

Der Anschluss an allgemeine Bäume erfolgt über
\FormulaRefAuto{B27RankedParentSystemTree}[thm]: Eine endliche
Knotenmenge mit eindeutigen Eltern und geeigneten Höhenmarken trägt
einen verwurzelten Baum. Die volle geordnete binäre Verzweigung ist
eine zusätzliche Bedingung des Klammerungsmodells.
Für die Syntax in Band~48 dienen Konstruktortrennung, Induktion und
Rekursion als wiederverwendbares Beweismuster. Die dortigen
Konstruktoren unterschiedlicher Stelligkeit benötigen jeweils ihre
eigenen Aufbauregeln; eine binäre Codierung ersetzt diese Regeln nicht.
\endgroup


\chapter[Band 28: Halbgruppen]{Band 28 auf einen Blick}
\noindent\textbf{Halbgruppen}\par\medskip
Grundlage sind binäre Operationen sowie die Wörter und Klammerungsbäume
aus Band 27. In dieser Übersicht bedeutet \(A\cong B\) stets
Isomorphie der jeweils angegebenen Halbgruppen.
\(\mathcal P_+(A)=\powerset(A)\setminus\{\varnothing\}\) trägt das Mengenprodukt.
Ein neutrales Element wird nicht vorausgesetzt; Potenzen haben positive Exponenten.

\begin{BandUebersicht}
\textbf{Assoziative binäre Operation}
\UebFormel{\begin{aligned}
&\star:A\times A\to A,\\
&a,b,c\in A\ \vdash\
 (a\star b)\star c=a\star(b\star c).
\end{aligned}}Diese Axiome definieren \(\SemigroupAlg A\star\).
Die Linksfaltung eines Nichtleerwortes heißt \(\WordFold\star w\).
\UebQuellen{\FormulaRefAuto{SemigroupDef}[def];
\FormulaRefAuto{LeftWordFoldDef}[def]}
&
\textbf{Blöcke und Klammerungen}
Bei \(\SemigroupAlg A\star\) gilt:\UebFormel{\begin{aligned}
&u,v\in\NonemptyWords A\ \vdash\\
&\WordFold\star{\WordConcat u v}
 =\WordFold\star u\star\WordFold\star v,\\
&w\in\NonemptyWords A,\ S,T\in\WordBracketings A w\
 \vdash\\
&\qquad\TreeEvaluation\star S=\TreeEvaluation\star T.
\end{aligned}}Addition und Multiplikation auf \(\mathbb N\) sind Beispiele.
\UebQuellen{\FormulaRefAuto{SemigroupWordBlockLaw}[thm];
\FormulaRefAuto{SemigroupBracketingIndependence}[thm];
\FormulaRefAuto{NaturalAdditionSemigroup}[thm];
\FormulaRefAuto{NaturalMultiplicationSemigroup}[thm]} \\
\addlinespace[.8ex]

\textbf{Positive Potenzen und Idempotente}
Für \(\SemigroupAlg A\star\) und \(a\in A\):\UebFormel{\begin{aligned}
&a^1=a,\qquad a^{n+1}=a^n\star a\quad(n\ge1),\\
&E(A)=\{e\in A\mid e\star e=e\}.
\end{aligned}}
\UebQuellen{\FormulaRefAuto{SemigroupPositivePowerDef}[def];
\FormulaRefAuto{SemigroupPowerOne}[thm];
\FormulaRefAuto{SemigroupPowerSuccessor}[thm];
\FormulaRefAuto{SemigroupIdempotentSetDef}[def]}
&
\textbf{Potenzen rechnen assoziativ}
Für \(m,n\in\mathbb N\setminus\{0\}\):\UebFormel{\begin{aligned}
&\SemigroupAlg A\star,\ a\in A\ \vdash\\
&a^{m+n}=a^m\star a^n,\qquad
 a^m\star a^n=a^n\star a^m,\\
&\SemigroupAlg A\star,\ e\in E(A)\ \vdash\ e^n=e.
\end{aligned}}
\UebQuellen{\FormulaRefAuto{SemigroupPowerAdditionLaw}[thm];
\FormulaRefAuto{SemigroupPowersCommute}[thm];
\FormulaRefAuto{SemigroupIdempotentPower}[thm]} \\
\addlinespace[.8ex]

\textbf{Unterhalbgruppen und Erzeugung}
Bei \(\SemigroupAlg A\star\) heißt \(U\) Unterhalbgruppe, wenn\UebFormel{\begin{aligned}
&\varnothing\ne U\subseteq A,\quad
 \forall x,y\in U\ (x\star y\in U).\\
&\varnothing\ne X\subseteq A\ \vdash\\
&\langle X\rangle=\bigcap\mathcal U_X,\\
&\mathcal U_X=\{U\mid\mathsf{UHGr}(U;A,\star),\\
&\hspace{29mm}X\subseteq U\}.
\end{aligned}}
\UebQuellen{\FormulaRefAuto{SemigroupSubsemigroupDef}[def];
\FormulaRefAuto{SemigroupGeneratedSubsemigroupDef}[def]}
&
\textbf{Kleinste produktabgeschlossene Erweiterung}
Für die links angegebenen Voraussetzungen:\UebFormel{\begin{aligned}
&\mathsf{UHGr}(\langle X\rangle;A,\star),\
 X\subseteq\langle X\rangle,\\
&\mathsf{UHGr}(U;A,\star),\ X\subseteq U
 \ \vdash\ \langle X\rangle\subseteq U,\\
&a\in A\ \vdash\
 \langle a\rangle=\{a^n\mid n\in\mathbb N,\ n\ge1\}.
\end{aligned}}
\UebQuellen{\FormulaRefAuto{GeneratedSubsemigroupIsSubsemigroup}[thm];
\FormulaRefAuto{GeneratedSubsemigroupContainsGenerators}[thm];
\FormulaRefAuto{GeneratedSubsemigroupMinimal}[thm];
\FormulaRefAuto{SemigroupMonogenicPowerDescription}[thm]} \\
\addlinespace[.8ex]

\textbf{Ideale und zweiseitige Teilbarkeit}
In \(\SemigroupAlg A\star\) ist ein Ideal eine nichtleere
Teilmenge \(I\subseteq A\) mit \(AI\cup IA\subseteq I\).
Für \(a\in A\) setzt man\UebFormel{\begin{aligned}
&J(a)=\{a\}\cup Aa\cup aA\cup AaA,\\
&a\mid_A b\ \leftrightarrow\\
&\quad b=a\lor\exists x\in A\ (b=x\star a)\\
&\quad{}\lor\exists y\in A\ (b=a\star y)\\
&\quad{}\lor\exists x,y\in A\ (b=(x\star a)\star y).
\end{aligned}}
\UebQuellen{\FormulaRefAuto{SemigroupIdealDef}[def];
\FormulaRefAuto{SemigroupPrincipalTwoSidedIdealDef}[def];
\FormulaRefAuto{SemigroupTwoSidedDividesDef}[def]}
&
\textbf{Hauptideal und Teilbarkeit}
Für \(\SemigroupAlg A\star\) und \(a,b,c\in A\):\UebFormel{\begin{aligned}
&\mathsf{Id}(J(a);A,\star),\\
&a\mid_A b\ \leftrightarrow\ b\in J(a),\\
&a\mid_A a,\qquad
 (a\mid_A b\land b\mid_A c)\to a\mid_A c.
\end{aligned}}
\UebQuellen{\FormulaRefAuto{SemigroupPrincipalTwoSidedIdealIsIdeal}[thm];
\FormulaRefAuto{SemigroupTwoSidedDividesPrincipalIdeal}[thm];
\FormulaRefAuto{SemigroupDivisibilityReflexive}[thm];
\FormulaRefAuto{SemigroupTwoSidedDividesTransitive}[thm]} \\
\addlinespace[.8ex]

\textbf{Green-Relationen}
In einer festen Halbgruppe und für \(a,b\in A\) vergleichen
\(L(a),R(a),J(a)\) die linken, rechten und zweiseitigen Hauptideale:\UebFormel{\begin{aligned}
&a\mathcal L b\leftrightarrow L(a)=L(b),\\
&a\mathcal R b\leftrightarrow R(a)=R(b),\\
&a\mathcal J b\leftrightarrow J(a)=J(b),\\
&\mathcal H=\mathcal L\cap\mathcal R.
\end{aligned}}
\UebQuellen{\FormulaRefAuto{SemigroupGreenLDef}[def];
\FormulaRefAuto{SemigroupGreenRDef}[def];
\FormulaRefAuto{SemigroupGreenJDef}[def];
\FormulaRefAuto{SemigroupGreenHDef}[def]}
&
\textbf{Vier Äquivalenzrelationen}
Bei \(\SemigroupAlg A\star\) sind
\(\mathcal L,\mathcal R,\mathcal J,\mathcal H\) reflexiv, symmetrisch und
transitiv auf \(A\). Insbesondere:\UebFormel{\begin{aligned}
&a,b\in A\ \vdash\\
&a\mathcal J b\ \leftrightarrow\
 (a\mid_A b\land b\mid_A a).
\end{aligned}}
\UebQuellen{\FormulaRefAuto{SemigroupGreenLIsEquivalence}[thm];
\FormulaRefAuto{SemigroupGreenRIsEquivalence}[thm];
\FormulaRefAuto{SemigroupGreenJIsEquivalence}[thm];
\FormulaRefAuto{SemigroupGreenHIsEquivalence}[thm];
\FormulaRefAuto{SemigroupGreenJMutualTwoSidedDivisibility}[thm]} \\
\addlinespace[.8ex]

\textbf{Homomorphismen}
Für Halbgruppen \(A,B\):\UebFormel{\begin{aligned}
&f:A\to B,\\
&x,y\in A\ \vdash\\
&\quad f(x\star y)=f(x)\diamond f(y).
\end{aligned}}Ein Isomorphismus ist zusätzlich bijektiv.
\UebQuellen{\FormulaRefAuto{SemigroupHomDef}[def];
\FormulaRefAuto{SemigroupIsoDef}[def]}
&
\textbf{Komposition und Mengenbild}
Homomorphismen \(f:A\to B\), \(g:B\to C\) sind komponierbar.
Ein Isomorphismus \(f\) induziert durch \(X\mapsto f[X]\):\UebFormel{\begin{aligned}
&A\cong B\ \vdash\
 \mathcal P_+(A)\cong\mathcal P_+(B).
\end{aligned}}
\UebQuellen{\FormulaRefAuto{SemigroupHomComposition}[thm];
\FormulaRefAuto{SemigroupIsoInducesPowerSemigroupIso}[thm]} \\
\addlinespace[.8ex]

\textbf{Potenzhalbgruppe und Einermengenkopie}
Für \(\SemigroupAlg A\star\) und \(X,Y\in\mathcal P_+(A)\):
\UebFormel{\begin{aligned}
&X\star_{\mathcal P}Y
 =\{x\star y\mid x\in X,\ y\in Y\},\\
&\mathcal S(A)=\{\{a\}\mid a\in A\}.
\end{aligned}}
\UebQuellen{\FormulaRefAuto{PowerSemigroupProductDef}[def];
\FormulaRefAuto{SemigroupSingletonCopyDef}[def]}
&
\textbf{Mengenprodukte bleiben assoziativ}
\UebFormel{\begin{aligned}
&\SemigroupAlg A\star\ \vdash\\
&\quad\SemigroupAlg{\mathcal P_+(A)}{\star_{\mathcal P}},\\
&A\cong\mathcal S(A),\qquad
 \{a\}\star_{\mathcal P}\{b\}=\{a\star b\}.
\end{aligned}}Die letzte Gleichung gilt für \(a,b\in A\).
\UebQuellen{\FormulaRefAuto{NonemptyPowerSemigroup}[thm];
\FormulaRefAuto{SemigroupSingletonCopyIsomorphism}[thm];
\FormulaRefAuto{PowerSemigroupProductMembership}[thm]} \\
\addlinespace[.8ex]

\textbf{Koordinaten an einem Basiselement}
Für \(\SemigroupAlg A\star\), \(e\in A\):\UebFormel{\begin{aligned}
&\mathcal K_e=(Ae)\times(eA),\\
&(p,q)\square_e(r,s)=(p,s),\\
&\chi_e(x)=(x\star e,e\star x).
\end{aligned}}
\UebQuellen{\FormulaRefAuto{SemigroupCoordinateCarrierDef}[def];
\FormulaRefAuto{SemigroupCoordinateProductDef}[def];
\FormulaRefAuto{SemigroupCoordinateMapDef}[def]}
&
\textbf{Kriterium für eine Koordinatenisomorphie}
Unter den linken Voraussetzungen genügen\UebFormel{\begin{aligned}
&\forall x\in A\ (x\star x=x),\\
&\forall x,y,z\in A\ ((x\star y)\star z=x\star z)\\
&\qquad\vdash\ \SemigroupCoordinateIso A\star e.
\end{aligned}}
\UebQuellen{\FormulaRefAuto{SemigroupCoordinateIsoCriterion}[thm]} \\
\addlinespace[.8ex]

\textbf{Inflation}
Eine Halbgruppe \(S\supseteq G\) ist eine Inflation von \(G\),
wenn ein Rückzug \(r:S\to G\) erfüllt:\UebFormel{\begin{aligned}
&r|_G=\operatorname{id}_G,\\
&x,y\in S\ \vdash\ x\star y=r(x)\star r(y).
\end{aligned}}
\UebQuellen{\FormulaRefAuto{SemigroupInflationDef}[def]}
&
\textbf{Rekonstruktion durch passende Fasern}
Für Inflationen \(r:S\to G\), \(s:T\to H\) und
einen Halbgruppenisomorphismus \(\varphi:G\to H\) gilt:\UebFormel{\begin{aligned}
&\forall g\in G\quad
 r^{-1}[\{g\}]\sim s^{-1}[\{\varphi(g)\}]\\
&\qquad\vdash\ S\cong T.
\end{aligned}}\(\sim\) bedeutet hier Existenz einer Bijektion.
\UebQuellen{\FormulaRefAuto{InflationFiberIsomorphism}[thm]} \\
\addlinespace[.8ex]

\textbf{Das Mogiljanskaja-Paar}
Für die disjunkten Schichten \(L=\{a_i\}\), \(D=\{b_{ij}\}\)
aus Band 21:\UebFormel{\begin{aligned}
&S=L\cup D\cup\{\mathbf0\},\quad T=S\cup\{k\},\\
&a_i a_j=b_{ij},\qquad
 xy=\mathbf0\ \text{sonst}.
\end{aligned}}Die markierten Elemente \(\mathbf0,k\) liegen außerhalb der Schichten.
Beide Operationen folgen derselben Regel.
\UebQuellen{\FormulaRefAuto{MogiljanskajaPairDef}[def];
\FormulaRefAuto{MogiljanskajaRuleDef}[def]}
&
\textbf{Unendliches Gegenbeispiel, einschließlich formalem Anhang}
\UebFormel{\begin{aligned}
&\neg\Finite S,\qquad\neg\Finite T,\\
&S\not\cong T,\qquad
 \mathcal P_+(S)\cong\mathcal P_+(T).
\end{aligned}}Die im Anhang ausgeführte Bijektion \(\Phi\) erhält sämtliche
Mengenprodukte. Dies widerlegt die uneingeschränkte Rekonstruktion.
\UebQuellen{\FormulaRefAuto{MogiljanskajaPairCounterexample}[thm];
\FormulaRefAuto{MogiljanskajaArgumentM16PowerIsomorphism}[thm];
\FormulaRefAuto{MogiljanskajaArgumentM17CounterexampleWitness}[thm]} \\
\addlinespace[.8ex]
\end{BandUebersicht}

\chapter[Band 29: Kommutative Halbgruppen]{Band 29 auf einen Blick}
\noindent\textbf{Kommutative Halbgruppen}\par\medskip
Vorausgesetzt sind Halbgruppen, zweiseitige Teilbarkeit und
die elementaren Mengen- und Additionsgesetze. \(p\mid_A a\)
bezeichnet die zweiseitige Halbgruppenteilbarkeit.

\begin{BandUebersicht}
\textbf{Kommutativität}
\UebFormel{\begin{aligned}
&\CommutativeSemigroupAlg A\star\ \leftrightarrow\\
&\quad\SemigroupAlg A\star\land\\
&\quad\forall a,b\in A\ (a\star b=b\star a).
\end{aligned}}
\UebQuellen{\FormulaRefAuto{CommutativeSemigroupDef}[def]}
&
\textbf{Vertauschungsgesetze}
Für \(\CommutativeSemigroupAlg A\star\) und \(a,b,c,d\in A\):\UebFormel{\begin{aligned}
&(a\star b)\star c=(a\star c)\star b,\\
&(a\star b)\star(c\star d)\\
&\qquad=(a\star c)\star(b\star d).
\end{aligned}}
\UebQuellen{\FormulaRefAuto{CommutativeSemigroupThreeTermExchange}[thm];
\FormulaRefAuto{CommutativeSemigroupFourTermExchange}[thm]} \\
\addlinespace[.8ex]

\textbf{Addition und Vereinigung}
\UebFormel{\begin{aligned}
&+:\mathbb N^2\to\mathbb N,\\
&\cup:\powerset(M)^2\to\powerset(M).
\end{aligned}}
\UebQuellen{\FormulaRefAuto{NaturalAdditionSemigroup}[thm];
\FormulaRefAuto{SemigroupDef}[def]}
&
\textbf{Kommutative Beispiele}
\UebFormel{\begin{aligned}
&\CommutativeSemigroupAlg{\mathbb N}{+},\\
&\CommutativeSemigroupAlg{\powerset(M)}{\cup}.
\end{aligned}}
\UebQuellen{\FormulaRefAuto{NaturalAdditionCommutativeSemigroup}[thm];
\FormulaRefAuto{PowerSetUnionCommutativeSemigroup}[thm]} \\
\addlinespace[.8ex]

\textbf{Primelement ohne vorausgesetzte Einheit}
In einer kommutativen Halbgruppe ist \(p\) prim, wenn\UebFormel{\begin{aligned}
&p\in A,\quad\exists q\in A\ (p\nmid_A q),\\
&a,b\in A,\ p\mid_A(a\star b)\ \vdash\\
&\qquad p\mid_A a\lor p\mid_A b.
\end{aligned}}
\UebQuellen{\FormulaRefAuto{CommutativeSemigroupPrimeDef}[def]}
&
\textbf{Echtheit in Idealnotation}
Mit dem zweiseitigen Hauptideal \(J(p)\) gilt:\UebFormel{\begin{aligned}
&p,q\in A\ \vdash\\
&\quad(p\mid_A q\leftrightarrow q\in J(p)),\\
&p\text{ prim}\ \vdash\ J(p)\subsetneq A.
\end{aligned}}Die letzte Folgerung entfaltet die Echtheitsforderung.
\UebQuellen{\FormulaRefAuto{SemigroupTwoSidedDividesPrincipalIdeal}[thm];
\FormulaRefAuto{CommutativeSemigroupPrimeDef}[def]} \\
\addlinespace[.8ex]

\textbf{Homomorphismen und Isomorphismen}
Für zwei Strukturen derselben Klasse bedeutet ein Homomorphismus
\(H(f)\):\UebFormel{\begin{aligned}
&f:A\to B,\\
&x,y\in A\ \vdash\\
&\quad f(x\star y)=f(x)\diamond f(y).
\end{aligned}}Ein Isomorphismus \(I(f)\) ist ein bijektiver Homomorphismus:\UebFormel{\begin{aligned}
&I(f)\ \leftrightarrow H(f)\land f:A\bij B.
\end{aligned}}
\UebQuellen{\FormulaRefAuto{CommutativeSemigroupHomDef}[def];
\FormulaRefAuto{CommutativeSemigroupIsoDef}[def]}
&
\textbf{Komposition erhält die Struktur}
Für Homomorphismen \(f:A\to B\), \(g:B\to C\) zwischen
Strukturen dieser Klasse ist \(g\circ f\) wieder ein Homomorphismus.
Mit \(\bullet\) als Operation auf \(C\) gilt für \(x,y\in A\):\UebFormel{\begin{aligned}
&(g\circ f)(x\star y)\\
&\quad=(g\circ f)(x)\mathbin{\bullet}(g\circ f)(y).
\end{aligned}}Sind \(f,g\) Isomorphismen, so ist auch \(g\circ f\) ein Isomorphismus.
\UebQuellen{\FormulaRefAuto{SemigroupHomComposition}[thm];
\FormulaRefAuto{SemigroupIsoComposition}[thm]} \\
\addlinespace[.8ex]
\end{BandUebersicht}


\chapter[Band 30: Idempotente Halbgruppen]{Band 30 auf einen Blick}
\noindent\textbf{Idempotente Halbgruppen}\par\medskip
Idempotenz wird zusätzlich zur Assoziativität gefordert,
Kommutativität dagegen nicht. Die Linksnulloperation \(x\star y=x\)
ist bei verschiedenen \(x,y\) ein nichtkommutatives Beispiel.

\begin{BandUebersicht}
\textbf{Idempotente Operation}
\UebFormel{\begin{aligned}
&\IdempotentSemigroupAlg A\star\ \leftrightarrow\\
&\quad\SemigroupAlg A\star\land\\
&\quad\forall x\in A\ (x\star x=x).
\end{aligned}}
\UebQuellen{\FormulaRefAuto{IdempotentSemigroupDef}[def]}
&
\textbf{Vereinigung als Beispiel}
\UebFormel{\begin{aligned}
&\IdempotentSemigroupAlg{\powerset(M)}{\cup},\\
&X\in\powerset(M)\ \vdash\ X\cup X=X.
\end{aligned}}
\UebQuellen{\FormulaRefAuto{PowerSetUnionIdempotentSemigroup}[thm];
\FormulaRefAuto{IdempotentSemigroupIdempotenceAxiom}[ax]} \\
\addlinespace[.8ex]

\textbf{Homomorphismen und Isomorphismen}
Für zwei Strukturen derselben Klasse bedeutet ein Homomorphismus
\(H(f)\):\UebFormel{\begin{aligned}
&f:A\to B,\\
&x,y\in A\ \vdash\\
&\quad f(x\star y)=f(x)\diamond f(y).
\end{aligned}}Ein Isomorphismus \(I(f)\) ist ein bijektiver Homomorphismus:\UebFormel{\begin{aligned}
&I(f)\ \leftrightarrow H(f)\land f:A\bij B.
\end{aligned}}
\UebQuellen{\FormulaRefAuto{IdempotentSemigroupHomDef}[def];
\FormulaRefAuto{IdempotentSemigroupIsoDef}[def]}
&
\textbf{Komposition erhält die Struktur}
Für Homomorphismen \(f:A\to B\), \(g:B\to C\) zwischen
Strukturen dieser Klasse ist \(g\circ f\) wieder ein Homomorphismus.
Mit \(\bullet\) als Operation auf \(C\) gilt für \(x,y\in A\):\UebFormel{\begin{aligned}
&(g\circ f)(x\star y)\\
&\quad=(g\circ f)(x)\mathbin{\bullet}(g\circ f)(y).
\end{aligned}}Sind \(f,g\) Isomorphismen, so ist auch \(g\circ f\) ein Isomorphismus.
\UebQuellen{\FormulaRefAuto{SemigroupHomComposition}[thm];
\FormulaRefAuto{SemigroupIsoComposition}[thm]} \\
\addlinespace[.8ex]
\end{BandUebersicht}


\chapter[Band 31: Rechtecksbänder]{Band 31 auf einen Blick}
\noindent\textbf{Rechtecksbänder}\par\medskip
Vorausgesetzt sind Halbgruppen, idempotente Operationen und
Potenzhalbgruppen. \(\mathcal P_+(A)\) bezeichnet die nichtleeren Teilmengen
mit Mengenprodukt; \(\cong\) stets Halbgruppenisomorphie.

\begin{BandUebersicht}
\textbf{Rechtecksidentität}
\UebFormel{\begin{aligned}
&\RectangularBandAlg R\star\ \leftrightarrow\\
&\quad\SemigroupAlg R\star\land\\
&\quad\forall x,y\in R\ ((x\star y)\star x=x).
\end{aligned}}
\UebQuellen{\FormulaRefAuto{RectangularBandDef}[def]}
&
\textbf{Idempotenz und Dreigliedrigkeit}
\UebFormel{\begin{aligned}
&\RectangularBandAlg R\star\ \vdash\
 \IdempotentSemigroupAlg R\star,\\
&\RectangularBandAlg R\star,\ x,y,z\in R\ \vdash\\
&\qquad(x\star y)\star z=x\star z.
\end{aligned}}
\UebQuellen{\FormulaRefAuto{RectangularBandIsIdempotentSemigroup}[thm];
\FormulaRefAuto{RectangularBandThreeTermLaw}[thm]} \\
\addlinespace[.8ex]

\textbf{Produktdarstellung}
Für \(\SemigroupAlg R\star\) und \(e\in R\) sind die Koordinaten\UebFormel{\begin{aligned}
&\mathcal K_e=(Re)\times(eR),\\
&(p,q)\square_e(r,s)=(p,s),\\
&\chi_e(x)=(x\star e,e\star x).
\end{aligned}}
\UebQuellen{\FormulaRefAuto{SemigroupCoordinateCarrierDef}[def];
\FormulaRefAuto{SemigroupCoordinateProductDef}[def];
\FormulaRefAuto{SemigroupCoordinateMapDef}[def]}
&
\textbf{Charakterisierung nichtleerer Rechtecksbänder}
\UebFormel{\begin{aligned}
&R\ne\varnothing\ \vdash\\
&\bigl(\RectangularBandAlg R\star\ \leftrightarrow\\
&\qquad\forall e\in R\
 \SemigroupCoordinateIso R\star e\bigr).
\end{aligned}}
\UebQuellen{\FormulaRefAuto{RectangularBandProductRepresentation}[thm]} \\
\addlinespace[.8ex]

\textbf{Rekonstruktionsproblem für endliche Bänder}
Betrachtet werden endliche idempotente Halbgruppen
\((A,\star)\), \((B,\diamond)\). Eine Potenzhalbgruppenisomorphie
ist eine Bijektion \(\Psi:\mathcal P_+(A)\to\mathcal P_+(B)\) mit\UebFormel{\begin{aligned}
&\Psi(X\star_{\mathcal P}Y)
   =\Psi(X)\diamond_{\mathcal P}\Psi(Y).
\end{aligned}}
\UebQuellen{\FormulaRefAuto{IdempotentSemigroupDef}[def];
\FormulaRefAuto{SemigroupIsoDef}[def]}
&
\textbf{Bestimmtheit in der bewiesenen endlichen Klasse}
\UebFormel{\begin{aligned}
&\IdempotentSemigroupAlg A\star,\
 \IdempotentSemigroupAlg B\diamond,\\
&\Finite A,\ \Finite B,\
 \mathcal P_+(A)\cong\mathcal P_+(B)\ \vdash\\
&\qquad A\cong B.
\end{aligned}}Der Band führt dies über Bandzerlegung, Komponententransport
und endliches Anheben auf die Grundelemente. Die Endlichkeits- und
Idempotenzvoraussetzungen bleiben bestehen.
\UebQuellen{\FormulaRefAuto{FiniteIdempotentSemigroupsGloballyDetermined}[thm]} \\
\addlinespace[.8ex]
\end{BandUebersicht}


\chapter[Band 32: Kommutative idempotente Halbgruppen]{Band 32 auf einen Blick}
\noindent\textbf{Kommutative idempotente Halbgruppen}\par\medskip
Die beiden vorherigen Strukturaxiomatiken werden gemeinsam gefordert.
Ihre ordnungstheoretische Lesart wird in Band 45 entwickelt.

\begin{BandUebersicht}
\textbf{Gemeinsame Struktur}
\UebFormel{\begin{aligned}
&\CommutativeIdempotentSemigroupAlg A\star\
 \leftrightarrow\\
&\quad\CommutativeSemigroupAlg A\star\land\\
&\quad\IdempotentSemigroupAlg A\star.
\end{aligned}}
\UebQuellen{\FormulaRefAuto{CommutativeIdempotentSemigroupDef}[def]}
&
\textbf{Vereinigung erfüllt beide Axiome}
\UebFormel{\begin{aligned}
&\CommutativeIdempotentSemigroupAlg{\powerset(M)}{\cup},\\
&X,Y\in\powerset(M)\ \vdash\\
&\quad X\cup Y=Y\cup X,\qquad X\cup X=X.
\end{aligned}}
\UebQuellen{\FormulaRefAuto{PowerSetUnionCommutativeIdempotentSemigroup}[thm];
\FormulaRefAuto{CommutativeIdempotentSemigroupCommutativityAxiom}[ax];
\FormulaRefAuto{CommutativeIdempotentSemigroupIdempotenceAxiom}[ax]} \\
\addlinespace[.8ex]

\textbf{Homomorphismen und Isomorphismen}
Für zwei Strukturen derselben Klasse bedeutet ein Homomorphismus
\(H(f)\):\UebFormel{\begin{aligned}
&f:A\to B,\\
&x,y\in A\ \vdash\\
&\quad f(x\star y)=f(x)\diamond f(y).
\end{aligned}}Ein Isomorphismus \(I(f)\) ist ein bijektiver Homomorphismus:\UebFormel{\begin{aligned}
&I(f)\ \leftrightarrow H(f)\land f:A\bij B.
\end{aligned}}
\UebQuellen{\FormulaRefAuto{CommutativeIdempotentSemigroupHomDef}[def];
\FormulaRefAuto{CommutativeIdempotentSemigroupIsoDef}[def]}
&
\textbf{Komposition erhält die Struktur}
Für Homomorphismen \(f:A\to B\), \(g:B\to C\) zwischen
Strukturen dieser Klasse ist \(g\circ f\) wieder ein Homomorphismus.
Mit \(\bullet\) als Operation auf \(C\) gilt für \(x,y\in A\):\UebFormel{\begin{aligned}
&(g\circ f)(x\star y)\\
&\quad=(g\circ f)(x)\mathbin{\bullet}(g\circ f)(y).
\end{aligned}}Sind \(f,g\) Isomorphismen, so ist auch \(g\circ f\) ein Isomorphismus.
\UebQuellen{\FormulaRefAuto{SemigroupHomComposition}[thm];
\FormulaRefAuto{SemigroupIsoComposition}[thm]} \\
\addlinespace[.8ex]
\end{BandUebersicht}


\chapter[Band 33: Halbgruppen mit Nullelement]{Band 33 auf einen Blick}
\noindent\textbf{Halbgruppen mit Nullelement}\par\medskip
Das Nullelement absorbiert auf beiden Seiten. Es ist von einem
neutralen Element zu unterscheiden. Die konkrete Operation bleibt stets Teil
der Struktur.

\begin{BandUebersicht}
\textbf{Beidseitige Nullabsorption}
\UebFormel{\begin{aligned}
&\SemigroupWithZeroAlg A\star0\ \leftrightarrow\\
&\quad\SemigroupAlg A\star\land0\in A\land\\
&\quad\forall a\in A\
 (0\star a=0\land a\star0=0).
\end{aligned}}
\UebQuellen{\FormulaRefAuto{SemigroupWithZeroDef}[def]}
&
\textbf{Eindeutigkeit der Null}
\UebFormel{\begin{aligned}
&\SemigroupWithZeroAlg A\star0,\ z\in A,\\
&\forall a\in A\
 (z\star a=z\land a\star z=z)\ \vdash\\
&\qquad z=0.
\end{aligned}}
\UebQuellen{\FormulaRefAuto{SemigroupZeroUnique}[thm]} \\
\addlinespace[.8ex]

\textbf{Konstante Operation}
Für \(0\in A\) werde die Operation festgelegt durch\UebFormel{\begin{aligned}
&a,b\in A\ \vdash\ a\star b=0.
\end{aligned}}
\UebQuellen{\FormulaRefAuto{SemigroupWithZeroDef}[def]}
&
\textbf{Konstante und natürliche Beispiele}
\UebFormel{\begin{aligned}
&0\in A,\ \forall a,b\in A\ (a\star b=0)\
 \vdash\\
&\qquad\SemigroupWithZeroAlg A\star0,\\
&\SemigroupWithZeroAlg{\mathbb N}{\cdot}{0}.
\end{aligned}}
\UebQuellen{\FormulaRefAuto{ConstantSemigroupWithZero}[thm];
\FormulaRefAuto{NaturalMultiplicationSemigroupWithZero}[thm]} \\
\addlinespace[.8ex]

\textbf{Annihilatoren}
Für \(\SemigroupWithZeroAlg A\star0\) und \(X\subseteq A\)
schreiben wir kurz\UebFormel{\begin{aligned}
&L(X)=\{a\in A\mid\\
&\hspace{4em}\forall x\in X,\ a\star x=0\},\\
&R(X)=\{a\in A\mid\\
&\hspace{4em}\forall x\in X,\ x\star a=0\}.
\end{aligned}}
\UebQuellen{\FormulaRefAuto{SemigroupLeftAnnihilatorDef}[def];
\FormulaRefAuto{SemigroupRightAnnihilatorDef}[def]}
&
\textbf{Zweiseitige Bedingung}
\UebFormel{\begin{aligned}
&\operatorname{Ann}(X)=L(X)\cap R(X),\\
&a\in\operatorname{Ann}(X)\ \leftrightarrow\\
&\quad a\in A\land\forall x\in X\\
&\qquad(a\star x=0\land x\star a=0).
\end{aligned}}Dies entfaltet die drei Annihilatordefinitionen.
\UebQuellen{\FormulaRefAuto{SemigroupAnnihilatorDef}[def];
\FormulaRefAuto{SemigroupLeftAnnihilatorDef}[def];
\FormulaRefAuto{SemigroupRightAnnihilatorDef}[def]} \\
\addlinespace[.8ex]

\textbf{Homomorphismen und Isomorphismen}
Für zwei Strukturen derselben Klasse bedeutet ein Homomorphismus
\(H(f)\):\UebFormel{\begin{aligned}
&f:A\to B,\\
&x,y\in A\ \vdash\\
&\quad f(x\star y)=f(x)\diamond f(y).
\end{aligned}}Zusätzlich gilt \(f(0_A)=0_B\).Ein Isomorphismus \(I(f)\) ist ein bijektiver Homomorphismus:\UebFormel{\begin{aligned}
&I(f)\ \leftrightarrow H(f)\land f:A\bij B.
\end{aligned}}
\UebQuellen{\FormulaRefAuto{SemigroupWithZeroHomDef}[def];
\FormulaRefAuto{SemigroupWithZeroIsoDef}[def]}
&
\textbf{Komposition erhält die Struktur}
Für Homomorphismen \(f:A\to B\), \(g:B\to C\) zwischen
Strukturen dieser Klasse ist \(g\circ f\) wieder ein Homomorphismus.
Mit \(\bullet\) als Operation auf \(C\) gilt für \(x,y\in A\):\UebFormel{\begin{aligned}
&(g\circ f)(x\star y)\\
&\quad=(g\circ f)(x)\mathbin{\bullet}(g\circ f)(y).
\end{aligned}}Sind \(f,g\) Isomorphismen, so ist auch \(g\circ f\) ein Isomorphismus.
\UebQuellen{\FormulaRefAuto{SemigroupHomComposition}[thm];
\FormulaRefAuto{SemigroupIsoComposition}[thm]} \\
\addlinespace[.8ex]
\end{BandUebersicht}


\chapter[Band 34: Linkskürzbare Halbgruppen]{Band 34 auf einen Blick}
\noindent\textbf{Linkskürzbare Halbgruppen}\par\medskip
Vorausgesetzt sind Halbgruppen und die natürliche Addition.
Links- und Rechtskürzbarkeit werden als getrennte Eigenschaften geführt;
in Band 36 werden beide gemeinsam verlangt.

\begin{BandUebersicht}
\textbf{Kürzungsaxiomatik}
Zur Halbgruppenstruktur kommt die angegebene Kürzungsrichtung:\UebFormel{\begin{aligned}
&a,b,c\in A,\ a\star b=a\star c\ \vdash b=c.
\end{aligned}}Das Strukturprädikat lautet \(\LeftCancellativeSemigroupAlg A\star\).
\UebQuellen{\FormulaRefAuto{LeftCancellativeSemigroupDef}[def]}
&
\textbf{Die natürliche Addition}
\UebFormel{\begin{aligned}
&\LeftCancellativeSemigroupAlg{\mathbb N}{+},\\
&a,b,c\in\mathbb N,\ a+b=a+c\ \vdash b=c.
\end{aligned}}
\UebQuellen{\FormulaRefAuto{NaturalAdditionLeftCancellativeSemigroup}[thm];
\FormulaRefAuto{NaturalAdditionLeftCancellative}[thm]} \\
\addlinespace[.8ex]

\textbf{Homomorphismen und Isomorphismen}
Für zwei Strukturen derselben Klasse bedeutet ein Homomorphismus
\(H(f)\):\UebFormel{\begin{aligned}
&f:A\to B,\\
&x,y\in A\ \vdash\\
&\quad f(x\star y)=f(x)\diamond f(y).
\end{aligned}}Ein Isomorphismus \(I(f)\) ist ein bijektiver Homomorphismus:\UebFormel{\begin{aligned}
&I(f)\ \leftrightarrow H(f)\land f:A\bij B.
\end{aligned}}
\UebQuellen{\FormulaRefAuto{LeftCancellativeSemigroupHomDef}[def];
\FormulaRefAuto{LeftCancellativeSemigroupIsoDef}[def]}
&
\textbf{Komposition erhält die Struktur}
Für Homomorphismen \(f:A\to B\), \(g:B\to C\) zwischen
Strukturen dieser Klasse ist \(g\circ f\) wieder ein Homomorphismus.
Mit \(\bullet\) als Operation auf \(C\) gilt für \(x,y\in A\):\UebFormel{\begin{aligned}
&(g\circ f)(x\star y)\\
&\quad=(g\circ f)(x)\mathbin{\bullet}(g\circ f)(y).
\end{aligned}}Sind \(f,g\) Isomorphismen, so ist auch \(g\circ f\) ein Isomorphismus.
\UebQuellen{\FormulaRefAuto{SemigroupHomComposition}[thm];
\FormulaRefAuto{SemigroupIsoComposition}[thm]} \\
\addlinespace[.8ex]
\end{BandUebersicht}


\chapter[Band 35: Rechtskürzbare Halbgruppen]{Band 35 auf einen Blick}
\noindent\textbf{Rechtskürzbare Halbgruppen}\par\medskip
Vorausgesetzt sind Halbgruppen und die natürliche Addition.
Links- und Rechtskürzbarkeit werden als getrennte Eigenschaften geführt;
in Band 36 werden beide gemeinsam verlangt.

\begin{BandUebersicht}
\textbf{Kürzungsaxiomatik}
Zur Halbgruppenstruktur kommt die angegebene Kürzungsrichtung:\UebFormel{\begin{aligned}
&a,b,c\in A,\ b\star a=c\star a\ \vdash b=c.
\end{aligned}}Das Strukturprädikat lautet \(\RightCancellativeSemigroupAlg A\star\).
\UebQuellen{\FormulaRefAuto{RightCancellativeSemigroupDef}[def]}
&
\textbf{Die natürliche Addition}
\UebFormel{\begin{aligned}
&\RightCancellativeSemigroupAlg{\mathbb N}{+},\\
&a,b,c\in\mathbb N,\ b+a=c+a\ \vdash b=c.
\end{aligned}}
\UebQuellen{\FormulaRefAuto{NaturalAdditionRightCancellativeSemigroup}[thm];
\FormulaRefAuto{NaturalAdditionRightCancellative}[thm]} \\
\addlinespace[.8ex]

\textbf{Kommutativität und Linkskürzung}
\UebFormel{\begin{aligned}
&\CommutativeSemigroupAlg A\star,\qquad\\
&\LeftCancellativeSemigroupAlg A\star.
\end{aligned}}
\UebQuellen{\FormulaRefAuto{CommutativeSemigroupDef}[def];
\FormulaRefAuto{LeftCancellativeSemigroupDef}[def]}
&
\textbf{Daraus folgt Rechtskürzung}
\UebFormel{\begin{aligned}
&\CommutativeSemigroupAlg A\star,\\
&\LeftCancellativeSemigroupAlg A\star\ \vdash\\
&\qquad\RightCancellativeSemigroupAlg A\star.
\end{aligned}}
\UebQuellen{\FormulaRefAuto{CommutativeLeftCancellativeSemigroupImpliesRightCancellativeSemigroup}[thm]} \\
\addlinespace[.8ex]

\textbf{Homomorphismen und Isomorphismen}
Für zwei Strukturen derselben Klasse bedeutet ein Homomorphismus
\(H(f)\):\UebFormel{\begin{aligned}
&f:A\to B,\\
&x,y\in A\ \vdash\\
&\quad f(x\star y)=f(x)\diamond f(y).
\end{aligned}}Ein Isomorphismus \(I(f)\) ist ein bijektiver Homomorphismus:\UebFormel{\begin{aligned}
&I(f)\ \leftrightarrow H(f)\land f:A\bij B.
\end{aligned}}
\UebQuellen{\FormulaRefAuto{RightCancellativeSemigroupHomDef}[def];
\FormulaRefAuto{RightCancellativeSemigroupIsoDef}[def]}
&
\textbf{Komposition erhält die Struktur}
Für Homomorphismen \(f:A\to B\), \(g:B\to C\) zwischen
Strukturen dieser Klasse ist \(g\circ f\) wieder ein Homomorphismus.
Mit \(\bullet\) als Operation auf \(C\) gilt für \(x,y\in A\):\UebFormel{\begin{aligned}
&(g\circ f)(x\star y)\\
&\quad=(g\circ f)(x)\mathbin{\bullet}(g\circ f)(y).
\end{aligned}}Sind \(f,g\) Isomorphismen, so ist auch \(g\circ f\) ein Isomorphismus.
\UebQuellen{\FormulaRefAuto{SemigroupHomComposition}[thm];
\FormulaRefAuto{SemigroupIsoComposition}[thm]} \\
\end{BandUebersicht}


\chapter[Band 36: Kürzbare Halbgruppen]{Band 36 auf einen Blick}
\noindent\textbf{Kürzbare Halbgruppen}\par\medskip
Vorausgesetzt sind Halbgruppen und die natürliche Addition.
Links- und Rechtskürzbarkeit werden als getrennte Eigenschaften geführt;
in Band 36 werden beide gemeinsam verlangt.

\begin{BandUebersicht}
\textbf{Kürzungsaxiomatik}
Zur Halbgruppenstruktur kommt die angegebene Kürzungsrichtung:\UebFormel{\begin{aligned}
&a,b,c\in A,\ a\star b=a\star c\ \vdash b=c,\\
&a,b,c\in A,\ b\star a=c\star a\ \vdash b=c.
\end{aligned}}Das Strukturprädikat lautet \(\CancellativeSemigroupAlg A\star\).
\UebQuellen{\FormulaRefAuto{CancellativeSemigroupDef}[def]}
&
\textbf{Die natürliche Addition}
\UebFormel{\begin{aligned}
&\CancellativeSemigroupAlg{\mathbb N}{+},\\
&a,b,c\in\mathbb N\ \vdash\\
&(a+b=a+c\to b=c),\\
&(b+a=c+a\to b=c).
\end{aligned}}
\UebQuellen{\FormulaRefAuto{NaturalAdditionCancellativeSemigroup}[thm];
\FormulaRefAuto{NaturalAdditionCancellative}[thm]} \\
\addlinespace[.8ex]

\textbf{Homomorphismen und Isomorphismen}
Für zwei Strukturen derselben Klasse bedeutet ein Homomorphismus
\(H(f)\):\UebFormel{\begin{aligned}
&f:A\to B,\\
&x,y\in A\ \vdash\\
&\quad f(x\star y)=f(x)\diamond f(y).
\end{aligned}}Ein Isomorphismus \(I(f)\) ist ein bijektiver Homomorphismus:\UebFormel{\begin{aligned}
&I(f)\ \leftrightarrow H(f)\land f:A\bij B.
\end{aligned}}
\UebQuellen{\FormulaRefAuto{CancellativeSemigroupHomDef}[def];
\FormulaRefAuto{CancellativeSemigroupIsoDef}[def]}
&
\textbf{Komposition erhält die Struktur}
Für Homomorphismen \(f:A\to B\), \(g:B\to C\) zwischen
Strukturen dieser Klasse ist \(g\circ f\) wieder ein Homomorphismus.
Mit \(\bullet\) als Operation auf \(C\) gilt für \(x,y\in A\):\UebFormel{\begin{aligned}
&(g\circ f)(x\star y)\\
&\quad=(g\circ f)(x)\mathbin{\bullet}(g\circ f)(y).
\end{aligned}}Sind \(f,g\) Isomorphismen, so ist auch \(g\circ f\) ein Isomorphismus.
\UebQuellen{\FormulaRefAuto{SemigroupHomComposition}[thm];
\FormulaRefAuto{SemigroupIsoComposition}[thm]} \\
\addlinespace[.8ex]
\end{BandUebersicht}


\chapter[Band 37: Endliche Halbgruppen]{Band 37 auf einen Blick}
\noindent\textbf{Endliche Halbgruppen}\par\medskip
Vorausgesetzt werden die Halbgruppen aus Band~28, Endlichkeit und
endliche Kardinalitäten. Die Rekonstruktion aus der großen Potenzhalbgruppe
wird nur in den ausdrücklich angegebenen Fällen behauptet.
\(A\cong B\) bezeichnet hier einen Halbgruppenisomorphismus.

\begin{BandUebersicht}
\textbf{Endlicher Grundträger}
\UebFormel{\begin{aligned}
&\FiniteSemigroupAlg A\star\\
&\quad\leftrightarrow\SemigroupAlg A\star
 \land\Finite A,\\
&X\star_{\mathcal P}Y
 =\{x\star y\mid x\in X,\ y\in Y\}.
\end{aligned}}
\UebQuellen{\FormulaRefAuto{FiniteSemigroupDef}[def];
\FormulaRefAuto{PowerSemigroupProductDef}[def]}
&
\textbf{Die große Potenzhalbgruppe}
\UebFormel{\begin{aligned}
&\FiniteSemigroupAlg A\star\ \vdash\\
&\FiniteSemigroupAlg{\PowNE A}{\star_{\mathcal P}},\\
&\FiniteCard(\PowNE A)=2^{\FiniteCard(A)}-1.
\end{aligned}}
\UebQuellen{\FormulaRefAuto{FinitePowerSemigroup}[thm];
\FormulaRefAuto{FiniteSemigroupNonemptyPowerSetCardinality}[thm]} \\
\addlinespace[.8ex]

\textbf{Index und Periode}
Für \(a\in A\) bedeutet
\(\operatorname{IP}_{A,\star}(a;\mu,\lambda)\), wobei \(r,s\in\mathbb N\):\UebFormel{\begin{aligned}
&\mu,\lambda\in\mathbb N\setminus\{0\},\\
&a^{\mu+\lambda}=a^\mu,\\
&1\le r<s<\mu+\lambda\ \vdash\ a^r\ne a^s.
\end{aligned}}
\UebQuellen{\FormulaRefAuto{FiniteSemigroupIndexPeriodPairDef}[def]}
&
\textbf{Eindeutigkeit und Potenzfolge}
In jeder endlichen Halbgruppe besitzt jedes \(a\in A\) genau ein
solches Paar. Für dieses gilt:\UebFormel{\begin{aligned}
&r\in\mathbb N,\ \mu\le r\ \vdash\\
&\qquad a^{r+\lambda}=a^r,\\
&\FiniteCard(\langle a\rangle)=\mu+\lambda-1.
\end{aligned}}
\UebQuellen{\FormulaRefAuto{FiniteSemigroupIndexPeriod}[thm];
\FormulaRefAuto{FiniteSemigroupPowerEventuallyPeriodic}[thm];
\FormulaRefAuto{FiniteSemigroupMonogenicSize}[thm]} \\
\addlinespace[.8ex]

\textbf{Monogene Halbgruppe und Idempotente}
Für \(\SemigroupAlg A\star\) und \(a\in A\):
\UebFormel{\begin{aligned}
&\langle a\rangle=\{a^n\mid
 n\in\mathbb N,\ n\ne0\},\\
&E_{A,\star}=\{e\in A\mid e\star e=e\}.
\end{aligned}}
\UebQuellen{\FormulaRefAuto{SemigroupMonogenicSubsemigroupDef}[def];
\FormulaRefAuto{SemigroupMonogenicPowerDescription}[thm];
\FormulaRefAuto{SemigroupIdempotentSetDef}[def]}
&
\textbf{Idempotente Potenzen}
\UebFormel{\begin{aligned}
&\FiniteSemigroupAlg A\star,\ a\in A\ \vdash\\
&\quad\exists n\in\mathbb N\setminus\{0\}\ 
 a^n\in E_{A,\star},\\
&\FiniteSemigroupAlg A\star,\ A\ne\varnothing\\
&\quad\vdash E_{A,\star}\ne\varnothing.
\end{aligned}}
\UebQuellen{\FormulaRefAuto{FiniteSemigroupIdempotentPower}[thm];
\FormulaRefAuto{FiniteSemigroupIdempotentExists}[thm]} \\
\addlinespace[.8ex]

\textbf{Kleinstes Ideal}
Für eine nichtleere endliche Halbgruppe \((A,\star)\) bezeichne
\(\operatorname{Id}(I)\) ein nichtleeres zweiseitiges Ideal:\UebFormel{\begin{aligned}
&\mathcal I=\{I\subseteq A\mid\operatorname{Id}(I)\},\\
&K=\iota I\bigl(\operatorname{Id}(I)\land\\
&\qquad\forall J\in\mathcal I\ (I\subseteq J)\bigr).
\end{aligned}}
\UebQuellen{\FormulaRefAuto{FiniteSemigroupIdealFamilyDef}[def];
\FormulaRefAuto{FiniteSemigroupLeastIdealDef}[def]}
&
\textbf{Existenz und Eindeutigkeit}
\UebFormel{\begin{aligned}
&\FiniteSemigroupAlg A\star,\ A\ne\varnothing\\
&\vdash\operatorname{Id}(K)\ \land\
 \forall I\in\mathcal I\ (K\subseteq I).
\end{aligned}}Damit ist \(K\) das eindeutig bestimmte kleinste nichtleere
zweiseitige Ideal.
\UebQuellen{\FormulaRefAuto{FiniteSemigroupLeastIdeal}[thm];
\FormulaRefAuto{FiniteSemigroupLeastIdealCharacterization}[thm]} \\
\addlinespace[.8ex]

\textbf{Beschriftete Multiplikationstafeln}
Für \(n=\FiniteCard(A)\) und \(e:\NatSeg n\bij A\) wird die Operation
auf \(\NatSeg n\) transportiert:\UebFormel{\begin{aligned}
&m_e(i,j)=e^{-1}(e(i)\star e(j)).
\end{aligned}}Der kanonische Code \(\operatorname{can}(A,\star)\) ist das Paar
aus \(n\) und der Menge aller zeilenweise gelesenen Tafeltupel
\(w_n(m_e)\) für sämtliche solchen Beschriftungen \(e\).
\UebQuellen{\FormulaRefAuto{FiniteSemigroupMultiplicationTableDef}[def];
\FormulaRefAuto{FiniteSemigroupTableTupleDef}[def];
\FormulaRefAuto{FiniteSemigroupCanonicalTableCodeDef}[def]}
&
\textbf{Klassifikation durch den Code}
Für endliche Halbgruppen \((A,\star)\), \((B,\diamond)\):\UebFormel{\begin{aligned}
&(A,\star)\cong(B,\diamond)\\
&\quad\leftrightarrow\
 \operatorname{can}(A,\star)
 =\operatorname{can}(B,\diamond).
\end{aligned}}Umbenennungsorbits der assoziativen Tafeln sind genau
die Isomorphieklassen.
\UebQuellen{\FormulaRefAuto{FiniteSemigroupCanonicalTableCodeComplete}[thm];
\FormulaRefAuto{FiniteSemigroupRelabelOrbitsAreIsomorphismClasses}[thm]} \\
\addlinespace[.8ex]

\textbf{Einermengenschicht}
\UebFormel{\begin{aligned}
&\operatorname{Sing}(A)=\{\{a\}\mid a\in A\}.
\end{aligned}}Sei \(F\) ein Isomorphismus der großen Potenzhalbgruppen von
\((A,\star)\) und \((B,\diamond)\).
\UebQuellen{\FormulaRefAuto{SemigroupSingletonCopyDef}[def];
\FormulaRefAuto{SemigroupIsoDef}[def]}
&
\textbf{Ein rekonstruierbarer Teilfall}
\UebFormel{\begin{aligned}
&F[\operatorname{Sing}(A)]
 =\operatorname{Sing}(B)\\
&\qquad\vdash (A,\star)\cong(B,\diamond).
\end{aligned}}Für endliche Halbgruppen mit
\(\FiniteCard(A)=\FiniteCard(B)=2\) genügt bereits die Existenz von \(F\).
\UebQuellen{\FormulaRefAuto{PowerSemigroupSingletonLayerReconstruction}[thm];
\FormulaRefAuto{FinitePowerSemigroupReconstructionOrderTwo}[thm]} \\
\addlinespace[.8ex]

\textbf{Endliche Suche}
Auf dem festen Träger \(\NatSeg n\) werden Operationen
\(m:(\NatSeg n)^2\to\NatSeg n\) mit\UebFormel{\begin{aligned}
&m(m(i,j),k)=m(i,m(j,k))\\
&\hfill(i,j,k\in\NatSeg n)
\end{aligned}}betrachtet und anschließend nach Umbenennung zusammengefasst.
\UebQuellen{\FormulaRefAuto{FiniteSemigroupAssociativeOperationsDef}[def];
\FormulaRefAuto{FiniteSemigroupRelabelledOperationDef}[def]}
&
\textbf{Reichweite des Verfahrens}
Die Tafelkodierung erlaubt die vollständige Prüfung einer fest
vorgegebenen endlichen Ordnung. Sie beweist nicht die allgemeine
Implikation\UebFormel{\begin{aligned}
&(\PowNE A,\star_{\mathcal P})
 \cong(\PowNE B,\diamond_{\mathcal P})\\
&\qquad\Longrightarrow (A,\star)\cong(B,\diamond).
\end{aligned}}Diese bleibt im allgemeinen endlichen Fall eine Vermutung.
\UebQuellen{\FormulaRefAuto{FiniteSemigroupCanonicalTableCodeComplete}[thm];
\FormulaRefAuto{FinitePowerSemigroupReconstructionOrderTwo}[thm]} \\
\addlinespace[.8ex]
\end{BandUebersicht}


\chapter[Band 38: Monoide]{Band 38 auf einen Blick}
\noindent\textbf{Monoide}\par\medskip
Dieser Band ergänzt Halbgruppen um ein neutrales Element.
Die Aussagen zu Teilbarkeit verwenden Kommutativität; bei
\glqq prim impliziert irreduzibel\grqq{} kommt Kürzbarkeit hinzu.
Potenzprodukte und Produktquadrate stammen aus Band~28.

\begin{BandUebersicht}
\textbf{Monoidaxiome}
\UebFormel{\begin{aligned}
&\MonoidAlg A\star e\ \leftrightarrow\\
&\SemigroupAlg A\star\land e\in A\\
&\quad\land\forall a\in A\ (e\star a=a=a\star e).
\end{aligned}}
\UebQuellen{\FormulaRefAuto{MonoidDef}[def]}
&
\textbf{Eindeutigkeit und Beispiele}
\UebFormel{\begin{aligned}
&\MonoidAlg A\star e,\ \MonoidAlg A\star u
 \ \vdash e=u,\\
&\CommutativeMonoidAlg{\mathbb N}{+}{0},\\
&\CommutativeMonoidAlg{\mathbb N}{\cdot}{1}.
\end{aligned}}
\UebQuellen{\FormulaRefAuto{MonoidIdentityUnique}[thm];
\FormulaRefAuto{NaturalAdditiveCommutativeMonoid}[thm];
\FormulaRefAuto{NaturalMultiplicativeCommutativeMonoid}[thm]} \\
\addlinespace[.8ex]

\textbf{Einheiten}
Für ein Monoid \((A,\star,e)\):\UebFormel{\begin{aligned}
&\MonoidUnit{A,\star,e}u\ \leftrightarrow\\
&u\in A\land\exists v\in A\
 (u\star v=e=v\star u).
\end{aligned}}
\UebQuellen{\FormulaRefAuto{MonoidUnitDef}[def]}
&
\textbf{Einheiten der Potenzhalbgruppe}
Die große Potenzhalbgruppe ist ein Monoid mit Einheit \(\{e\}\).
Ihre invertierbaren Elemente sind genau die Einermengen invertierbarer
Grundelemente:\UebFormel{\begin{aligned}
&X\in\PowNE A\ \vdash\\
&\MonoidUnit{\PowNE A,\star_{\mathcal P},\{e\}}X\\
&\quad\leftrightarrow\exists u\in A\
 (X=\{u\}\land\MonoidUnit{A,\star,e}u).
\end{aligned}}
\UebQuellen{\FormulaRefAuto{PowerSemigroupMonoid}[thm];
\FormulaRefAuto{PowerMonoidUnitCharacterization}[thm]} \\
\addlinespace[.8ex]

\textbf{Teilbarkeit im kommutativen Monoid}
Bei \(\CommutativeMonoidAlg A\star e\) und \(a,b\in A\):\UebFormel{\begin{aligned}
&a\mid b\ \leftrightarrow\exists c\in A\ (b=a\star c).
\end{aligned}}
\UebQuellen{\FormulaRefAuto{CommutativeMonoidDividesDef}[def]}
&
\textbf{Einheiten als Teiler der Eins}
\UebFormel{\begin{aligned}
&\CommutativeMonoidAlg A\star e,\ u\in A\\
&\vdash\MonoidUnit{A,\star,e}u
 \ \leftrightarrow u\mid e.
\end{aligned}}
\UebQuellen{\FormulaRefAuto{CommutativeMonoidUnitIffDividesIdentity}[thm]} \\
\addlinespace[.8ex]

\textbf{Prim und irreduzibel}
Im kommutativen Monoid ist \(p\) irreduzibel, wenn \(p\) keine
Einheit ist und jede Zerlegung \(p=a\star b\) einen Einheitsfaktor besitzt.
Es ist prim, wenn es keine Einheit ist und\UebFormel{\begin{aligned}
&a,b\in A,\ p\mid a\star b\\
&\qquad\vdash p\mid a\lor p\mid b.
\end{aligned}}
\UebQuellen{\FormulaRefAuto{CommutativeMonoidIrreducibleDef}[def];
\FormulaRefAuto{CommutativeMonoidPrimeDef}[def]}
&
\textbf{Prim impliziert irreduzibel}
\UebFormel{\begin{aligned}
&\CommutativeMonoidAlg A\star e,\\
&\CancellativeSemigroupAlg A\star,\\
&\MonoidPrime{A,\star,e}p\\
&\qquad\vdash\MonoidIrreducible{A,\star,e}p.
\end{aligned}}
\UebQuellen{\FormulaRefAuto{CommutativeCancellativeMonoidPrimeIsIrreducible}[thm]} \\
\addlinespace[.8ex]

\textbf{Monoidmorphismen}
Für Monoide \(\mathfrak A=(A,\star,e_A)\),
\(\mathfrak B=(B,\diamond,e_B)\) bedeutet
\(\MonoidHom f{\mathfrak A}{\mathfrak B}\):\UebFormel{\begin{aligned}
&f:A\to B,\quad f(e_A)=e_B,\\
&x,y\in A\ \vdash\
 f(x\star y)=f(x)\diamond f(y).
\end{aligned}}
\UebQuellen{\FormulaRefAuto{MonoidHomDef}[def];
\FormulaRefAuto{MonoidIsoDef}[def]}
&
\textbf{Komposition und Einheiten}
Die Komposition zweier Monoidmorphismen ist wieder ein
Monoidmorphismus. Ein Halbgruppenisomorphismus zwischen Monoiden
erhält automatisch die neutralen Elemente und die Einheiten:\UebFormel{\begin{aligned}
&f(e_A)=e_B,\\
&\MonoidUnit{A,\star,e_A}u\\
&\qquad\leftrightarrow\MonoidUnit{B,\diamond,e_B}{f(u)}
 \quad(u\in A).
\end{aligned}}
\UebQuellen{\FormulaRefAuto{MonoidHomComposition}[thm];
\FormulaRefAuto{SemigroupIsoBetweenMonoidsIdentityAndUnits}[thm]} \\
\addlinespace[.8ex]

\textbf{Monoidales Produktquadrat}
Für eine nichtleere Halbgruppe \((A,\star)\) sei
\(A^2=\{a\star b\mid a,b\in A\}\).
Wenn \(\MonoidAlg{A^2}\star e\), setzt man\UebFormel{\begin{aligned}
&r_e:A\to A^2,\qquad r_e(a)=e\star a.
\end{aligned}}
\UebQuellen{\FormulaRefAuto{SemigroupSquareDef}[def];
\FormulaRefAuto{MonoidSquareInflationDef}[def]}
&
\textbf{Rückführung auf das Quadrat}
Unter diesen Voraussetzungen gilt:\UebFormel{\begin{aligned}
&r_e|_{A^2}=\operatorname{id}_{A^2},\\
&a,b\in A\ \vdash\
 a\star b=r_e(a)\star r_e(b).
\end{aligned}}Ferner besitzt \((\PowNE A)^2\) genau dann ein neutrales Element,
wenn \(A^2\) eines besitzt.
\UebQuellen{\FormulaRefAuto{MonoidSquareInflationProperties}[thm];
\FormulaRefAuto{PowerSemigroupSquareMonoidCriterion}[thm]} \\
\addlinespace[.8ex]

\textbf{Fasern der Rückführung}
Für die vorige Rückführung und \(g\in A^2\):\UebFormel{\begin{aligned}
&C_g=\{a\in A\mid r_e(a)=g\}.
\end{aligned}}
\UebQuellen{\FormulaRefAuto{MonoidSquareInflationDef}[def]}
&
\textbf{Potenzfasern}
\UebFormel{\begin{aligned}
&\{X\in\PowNE A\mid
 \{e\}\star_{\mathcal P}X=\{g\}\}\\
&\qquad=\PowNE{C_g}.
\end{aligned}}Die Faser über einer Einermenge ist somit genau die
nichtleere Potenzmenge der entsprechenden Grundfaser.
\UebQuellen{\FormulaRefAuto{MonoidSquarePowerFiber}[thm]} \\
\addlinespace[.8ex]
\end{BandUebersicht}


\chapter[Band 39: Halbringe]{Band 39 auf einen Blick}
\noindent\textbf{Halbringe}\par\medskip
Die additiven und multiplikativen Monoidaxiome werden hier verbunden.
Ein Halbring muss nicht \(0\ne1\) erfüllen; die Multiplikation muss
nicht kommutativ sein. Die natürliche Abbildung wird mittels der bereits
entwickelten Iteration konstruiert.

\begin{BandUebersicht}
\textbf{Halbringaxiome}
Ein Halbring \((R,+,\cdot,0,1)\) erfüllt:\UebFormel{\begin{aligned}
&\CommutativeMonoidAlg R+0,\\
&\MonoidAlg R\cdot1,\\
&x(y+z)=xy+xz,\\
&(x+y)z=xz+yz,\\
&0x=0=x0\qquad(x,y,z\in R).
\end{aligned}}Kommutativ heißt zusätzlich \(xy=yx\) für \(x,y\in R\).
\UebQuellen{\FormulaRefAuto{SemiringDef}[def];
\FormulaRefAuto{CommutativeUnitalSemiringDef}[def]}
&
\textbf{Natürliche Zahlen als Halbring}
\UebFormel{\begin{aligned}
&\SemiringAlg{\mathbb N}{+}{\cdot}{0}{1},\\
&\CommutativeSemiringAlg{\mathbb N}{+}{\cdot}{0}{1}.
\end{aligned}}Die additive Inversenbedingung eines Rings gilt hier nicht:
insbesondere hat \(1\) kein additives Inverses in \(\mathbb N\).
\UebQuellen{\FormulaRefAuto{NaturalUnitalSemiring}[thm];
\FormulaRefAuto{NaturalCommutativeUnitalSemiring}[thm];
\FormulaRefAuto{NaturalNumbersNotRing}[thm]} \\
\addlinespace[.8ex]

\textbf{Morphismen}
Für Halbringe \(\mathfrak R,\mathfrak S\) ist ein
Halbringhomomorphismus \(f:R\to S\) durch\UebFormel{\begin{aligned}
&f(x+y)=f(x)+f(y),\\
&f(xy)=f(x)f(y),\\
&f(0_R)=0_S,\qquad f(1_R)=1_S
\end{aligned}}gegeben. Ein Isomorphismus ist ein bijektiver solcher Morphismus.
\UebQuellen{\FormulaRefAuto{SemiringHomDef}[def];
\FormulaRefAuto{SemiringIsoDef}[def]}
&
\textbf{Komposition}
Für Halbringe \(\mathfrak R,\mathfrak S,\mathfrak T\):\UebFormel{\begin{aligned}
&\SemiringHom f{\mathfrak R}{\mathfrak S},\\
&\SemiringHom g{\mathfrak S}{\mathfrak T}\\
&\qquad\vdash
 \SemiringHom{g\circ f}{\mathfrak R}{\mathfrak T}.
\end{aligned}}
\UebQuellen{\FormulaRefAuto{SemiringHomComposition}[thm]} \\
\addlinespace[.8ex]

\textbf{Die natürliche Abbildung}
Für einen Halbring
\(\mathfrak R=(R,\oplus,\odot,0_R,1_R)\):\UebFormel{\begin{aligned}
&\sigma_R(x)=x\oplus1_R,\\
&\eta_R(0)=0_R,\\
&\eta_R(n+1)=\sigma_R(\eta_R(n))
 \quad(n\in\mathbb N).
\end{aligned}}
\UebQuellen{\FormulaRefAuto{SemiringSuccessorDef}[def];
\FormulaRefAuto{NaturalSemiringMapDef}[def]}
&
\textbf{Universelle Eigenschaft}
Die Abbildung \(\eta_R:\mathbb N\to R\) erfüllt für \(m,n\in\mathbb N\):\UebFormel{\begin{aligned}
&\eta_R(1)=1_R,\\
&\eta_R(m+n)=\eta_R(m)\oplus\eta_R(n),\\
&\eta_R(mn)=\eta_R(m)\odot\eta_R(n).
\end{aligned}}Sie ist der eindeutig bestimmte Halbringhomomorphismus
von \((\mathbb N,+,\cdot,0,1)\) nach \(\mathfrak R\).
\UebQuellen{\FormulaRefAuto{NaturalSemiringMapOne}[thm];
\FormulaRefAuto{NaturalSemiringMapAddition}[thm];
\FormulaRefAuto{NaturalSemiringMapMultiplication}[thm];
\FormulaRefAuto{NaturalSemiringUniversalHomomorphism}[thm]} \\
\addlinespace[.8ex]

\textbf{Nachfolgerinduktion im Ziel}
Die Induktionsbedingung lautet:\UebFormel{\begin{aligned}
&A\subseteq R,\quad 0_R\in A,\\
&\forall x\in A\ (\sigma_R(x)\in A)\\
&\qquad\Longrightarrow A=R.
\end{aligned}}
\UebQuellen{\FormulaRefAuto{NaturalSemiringMapSurjectiveIffSuccessorInduction}[thm]}
&
\textbf{Surjektivität und Injektivität}
Für einen Halbring \(R\) ist diese Bedingung äquivalent zu
\(\eta_R:\mathbb N\sur R\). Außerdem gilt:\UebFormel{\begin{aligned}
&\sigma_R:R\inj R,\quad
 \forall x\in R\ (\sigma_R(x)\ne0_R)\\
&\qquad\vdash\eta_R:\mathbb N\inj R.
\end{aligned}}
\UebQuellen{\FormulaRefAuto{NaturalSemiringMapSurjectiveIffSuccessorInduction}[thm];
\FormulaRefAuto{NaturalSemiringMapInjectivePeanoCriterion}[thm]} \\
\addlinespace[.8ex]

\textbf{Endlichkeit und Isomorphie}
Ein endlicher Halbring erfüllt zusätzlich
\(\Finite R\). Die kanonische Abbildung bleibt stets \(\eta_R\).
\UebQuellen{\FormulaRefAuto{FiniteSemiringDef}[def];
\FormulaRefAuto{NaturalSemiringMapDef}[def]}
&
\textbf{Grenzen und exaktes Kriterium}
\UebFormel{\begin{aligned}
&\SemiringAlg R\oplus\odot{0_R}{1_R},
 \ \Finite R\\
&\qquad\vdash\neg(\eta_R:\mathbb N\inj R).
\end{aligned}}Für beliebige Halbringe ist \(\eta_R\) genau dann ein
Halbringisomorphismus, wenn es injektiv und surjektiv ist.
\UebQuellen{\FormulaRefAuto{NaturalSemiringMapNotInjectiveIntoFiniteSemiring}[thm];
\FormulaRefAuto{NaturalSemiringMapIsoIffInjectiveAndSurjective}[thm]} \\
\addlinespace[.8ex]
\end{BandUebersicht}

\chapter[Band 40: Gruppen]{Band 40 auf einen Blick}
\noindent\textbf{Gruppen}\par\medskip
Vorausgesetzt werden die Halbgruppen und die Monoide aus Band~38.
Operation und neutrales Element gehören stets zur Strukturbezeichnung.

\begin{BandUebersicht}
\textbf{Gruppenaxiome}

\UebFormel{\begin{aligned}&\GroupAlg G\star e\ \leftrightarrow\\&\MonoidAlg G\star e\ \land\\&\forall x\in G\,\exists y\in G\\&\quad(x\star y=e\land y\star x=e).\end{aligned}}
\UebQuellen{\FormulaRefAuto{GroupDef}[def]}
&
\textbf{Inverse und Kürzung}
Für \(x,y,z\in G\):
\UebFormel{\begin{aligned}&(x\star y)^{-1}=y^{-1}\star x^{-1},\\&x\star y=x\star z\ \vdash y=z,\\&y\star x=z\star x\ \vdash y=z.\end{aligned}}
\UebQuellen{\FormulaRefAuto{GroupInverseUnique}[thm];
\FormulaRefAuto{GroupProductInverse}[thm];
\FormulaRefAuto{GroupLeftCancellation}[thm];
\FormulaRefAuto{GroupRightCancellation}[thm]} \\
\addlinespace[.8ex]

\textbf{Gruppenhomomorphismen}
Eine Abbildung zwischen Gruppen erhält die Operation; Isomorphismen sind bijektiv.
\UebFormel{f(x\star y)=f(x)\diamond f(y).}
\UebQuellen{\FormulaRefAuto{GroupHomDef}[def];
\FormulaRefAuto{GroupIsoDef}[def]}
&
\textbf{Erhaltung}

\UebFormel{\begin{aligned}&f(e_G)=e_H,\\&f(x^{-1})=f(x)^{-1}.\end{aligned}}
\UebQuellen{\FormulaRefAuto{GroupHomIdentityPreservation}[thm];
\FormulaRefAuto{GroupHomInversePreservation}[thm];
\FormulaRefAuto{GroupHomComposition}[thm]} \\
\addlinespace[.8ex]

\textbf{Einheitengruppe}

\UebFormel{\begin{aligned}&A^\times=\{u\in A\mid\MonoidUnit{A,\star,e}{u}\},\\&\MonoidAlg A\star e\ \vdash\GroupAlg{A^\times}\star e.\end{aligned}}
\UebQuellen{\FormulaRefAuto{MonoidUnitSetDef}[def];
\FormulaRefAuto{MonoidUnitSetIsGroup}[thm]}
&
\textbf{Rekonstruktion der Einheiten}
Für einen Isomorphismus \(F\) der großen Potenzhalbgruppen zweier Monoide:
\UebFormel{\begin{aligned}&F(A^\times)=B^\times,\\&F[\PowNE{A^\times}]=\PowNE{B^\times}.\end{aligned}}
\UebQuellen{\FormulaRefAuto{LargePowerSemigroupUnitGroupReconstruction}[thm]} \\
\addlinespace[.8ex]

\textbf{Gruppenpotenzhalbgruppe}

\UebFormel{\begin{aligned}&X\star_{\mathcal P}Y=\{x\star y\mid x\in X,y\in Y\},\\&\mathrm{Sing}(G)=\{\{g\}\mid g\in G\}.\end{aligned}}
\UebQuellen{\FormulaRefAuto{PowerSemigroupProductDef}[def];
\FormulaRefAuto{PowerGroupUnitsAreSingletons}[thm]}
&
\textbf{Einseitige Gruppenstarrheit}
Ist \(G\) eine Gruppe und \(B\) eine Halbgruppe, so trägt auch \(B\) eine Gruppenstruktur.
\UebFormel{\begin{aligned}&(\PowNE G,\star_{\mathcal P})\cong(\PowNE B,\diamond_{\mathcal P})\\&\qquad\vdash (G,\star)\cong(B,\diamond).\end{aligned}}
\UebQuellen{\FormulaRefAuto{GroupPowerSemigroupRigidity}[thm];
\FormulaRefAuto{PowerGroupsSingletonTransport}[thm]} \\
\addlinespace[.8ex]

\textbf{Endliche Kürzbarkeit}

\UebFormel{\CancellativeSemigroupAlg A\star,\quad\Finite A,\quad A\ne\varnothing.}
\UebQuellen{\FormulaRefAuto{CancellativeSemigroupDef}[def]}
&
\textbf{Gewinnung einer Gruppe}

\UebFormel{\exists e\in A\,\GroupAlg A\star e.}
\UebQuellen{\FormulaRefAuto{FiniteCancellativeSemigroupIsGroup}[thm]} \\
\addlinespace[.8ex]
\end{BandUebersicht}

\chapter[Band 41: Abelsche Gruppen]{Band 41 auf einen Blick}
\noindent\textbf{Abelsche Gruppen}\par\medskip
Dieser Band ergänzt die Gruppenaxiome aus Band~40 um Kommutativität.
Die Differenzenkonstruktion setzt ein kommutatives kürzbares Monoid voraus.

\begin{BandUebersicht}
\textbf{Abelsche Gruppe}

\UebFormel{\begin{aligned}&\AbelianGroupAlg G\star e\ \leftrightarrow\\&\GroupAlg G\star e\land\\&\forall x,y\in G\,(x\star y=y\star x).\end{aligned}}
\UebQuellen{\FormulaRefAuto{AbelianGroupDef}[def]}
&
\textbf{Morphismen}
Quelle und Ziel sind abelsche Gruppen. Ein Homomorphismus erhält die Operation; ein Isomorphismus ist zusätzlich bijektiv.
\UebFormel{\AbelianGroupIso f{G,\star,e}{H,\diamond,u}.}
\UebQuellen{\FormulaRefAuto{AbelianGroupHomDef}[def];
\FormulaRefAuto{AbelianGroupIsoDef}[def]} \\
\addlinespace[.8ex]

\textbf{Formale Differenzen}

\UebFormel{\begin{aligned}&(a,b)\sim(c,d)\ \leftrightarrow a\star d=b\star c,\\&K(M)=(M\times M)/{\sim},\\&[a,b]\oplus[c,d]=[a\star c,b\star d].\end{aligned}}
\UebQuellen{\FormulaRefAuto{GroupCompletionPairRelationDef}[def];
\FormulaRefAuto{GroupCompletionCarrierDef}[def];
\FormulaRefAuto{GroupCompletionOperationDef}[def]}
&
\textbf{Gruppenvervollständigung}

\UebFormel{\begin{aligned}&[a,b]=[c,d]\ \leftrightarrow a\star d=b\star c,\\&\AbelianGroupAlg{K(M)}\oplus{[e,e]}.\end{aligned}}
\UebQuellen{\FormulaRefAuto{GroupCompletionClassEquality}[thm];
\FormulaRefAuto{GroupCompletionClassInverse}[thm];
\FormulaRefAuto{GroupCompletionAbelianGroup}[thm]} \\
\addlinespace[.8ex]

\textbf{Kanonische Einbettung}

\UebFormel{\iota(a)=[a,e].}
\UebQuellen{\FormulaRefAuto{GroupCompletionEmbeddingDef}[def]}
&
\textbf{Einbettungssatz}

\UebFormel{\begin{aligned}&\iota:M\inj K(M),\quad\iota(e)=[e,e],\\&\iota(a\star b)=\iota(a)\oplus\iota(b).\end{aligned}}
\UebQuellen{\FormulaRefAuto{GroupCompletionEmbeddingTheorem}[thm]} \\
\addlinespace[.8ex]

\textbf{Ganzzahlige Addition}
Die Rechengesetze stammen aus der Quotientenkonstruktion in Band~17.
\UebFormel{\IntClass a b+\IntClass c d=\IntClass{a+c}{b+d}.}
\UebQuellen{\FormulaRefAuto{IntegerAdditionDef}[def]}
&
\textbf{Additives Grundbeispiel}

\UebFormel{\AbelianGroupAlg{\mathbb Z}{+}{\IntZero}.}
\UebQuellen{\FormulaRefAuto{IntegerAdditiveAbelianGroup}[thm]} \\
\addlinespace[.8ex]
\end{BandUebersicht}

\chapter[Band 42: Endliche Gruppen]{Band 42 auf einen Blick}
\noindent\textbf{Endliche Gruppen}\par\medskip
Der Band besitzt einen eigenen Begriff der endlichen Gruppe.
Sein ausgearbeiteter Schwerpunkt ist die Rekonstruktion endlicher
Halbgruppen mit einem \emph{Gruppenquadrat}. Benötigt werden insbesondere
Potenzhalbgruppen, Inflationen, Faserbijektionen und endliche
Potenzmengenkardinalitäten.

\begin{BandUebersicht}
\textbf{Eigene Struktur: endliche Gruppe}
\UebFormel{\begin{aligned}
&\mathsf{FinGrp}(G,\star,e)\ \leftrightarrow\\
&\GroupAlg G\star e\land\Finite G.
\end{aligned}}Endlichkeit ist hier ausdrücklich Bestandteil des Strukturbegriffs.
\UebQuellen{\FormulaRefAuto{FiniteGroupDef}[def]}
&
\textbf{Einordnung}
\UebFormel{\begin{aligned}
&\mathsf{FinGrp}(G,\star,e)\ \vdash\\
&\qquad\FiniteSemigroupAlg G\star.
\end{aligned}}Insbesondere sind die Sätze über endliche Halbgruppen auf den
zugrunde liegenden Halbgruppenträger anwendbar.
\UebQuellen{\FormulaRefAuto{FiniteGroupIsFiniteSemigroup}[thm]} \\
\addlinespace[.8ex]

\textbf{Gruppenquadrat und Potenzisomorphismus}
Für nichtleere Halbgruppen \((S,\star)\), \((T,\diamond)\) setze\UebFormel{\begin{aligned}
&G=S^2,\quad H=T^2,\\
&Q_S=\PowNE S,\\
&Q_T=\PowNE T.
\end{aligned}}Es sei \(G\) eine Gruppe mit Eins \(e\) und \(F\) ein
Isomorphismus von \((Q_S,\star_{\mathcal P})\) nach
\((Q_T,\diamond_{\mathcal P})\).
\UebQuellen{\FormulaRefAuto{SemigroupSquareDef}[def];
\FormulaRefAuto{SemigroupIsoDef}[def];
\FormulaRefAuto{GroupDef}[def]}
&
\textbf{Der Gruppenkern wird rekonstruiert}
Dann existieren \(u\in H\) und ein Gruppenisomorphismus
\(\varphi:(G,\star,e)\cong(H,\diamond,u)\) mit\UebFormel{\begin{aligned}
&g\in G\ \vdash\
 F(\{g\})=\{\varphi(g)\}.
\end{aligned}}Für diesen Kernschritt ist noch keine Endlichkeit nötig.
\UebQuellen{\FormulaRefAuto{PowerGroupSquareCoreReconstruction}[thm]} \\
\addlinespace[.8ex]

\textbf{Grundfasern der Inflation}
Sind nun beide Halbgruppen endlich und sind \(u,\varphi\) wie
zuvor gewählt, so seien\UebFormel{\begin{aligned}
&r(x)=e\star x,\qquad s(y)=u\diamond y,\\
&C_g=\{x\in S\mid r(x)=g\},\\
&D_h=\{y\in T\mid s(y)=h\}.
\end{aligned}}
\UebQuellen{\FormulaRefAuto{MonoidSquareInflationDef}[def]}
&
\textbf{Endliche Fasern sind gleichmächtig}
Unter allen Voraussetzungen des vorigen Schritts und
\(\Finite S,\Finite T\) gilt:\UebFormel{\begin{aligned}
&\forall g\in G\ \exists f_g:
 C_g\bij D_{\varphi(g)}.
\end{aligned}}Der Potenzisomorphismus liefert zunächst Potenzfaserbijektionen;
die Endlichkeit erlaubt den Rückschluss auf die Grundfasern.
\UebQuellen{\FormulaRefAuto{FiniteGroupSquareFiberBijections}[thm]} \\
\addlinespace[.8ex]

\textbf{Zusammenfügen über dem Gruppenquadrat}
Das Quadrat ist eine Unterhalbgruppe; die Rückführungen
\(r:S\to G\), \(s:T\to H\) machen \(S,T\) zu Inflationen.
Die vorige Faserinformation wird mit dem Gruppenisomorphismus verklebt.
\UebQuellen{\FormulaRefAuto{MonoidSquareInflationProperties}[thm];
\FormulaRefAuto{FiniteGroupSquareFiberBijections}[thm]}
&
\textbf{Rekonstruktionssatz}
Für endliche Halbgruppen \(S,T\) gilt:\UebFormel{\begin{aligned}
&S\ne\varnothing,\quad
 \GroupAlg{S^2}\star e,\\
&(\PowNE S,\star_{\mathcal P})
 \cong(\PowNE T,\diamond_{\mathcal P})\\
&\qquad\vdash (S,\star)\cong(T,\diamond).
\end{aligned}}Die Nichtleerheit von \(T\) folgt bereits aus diesen Voraussetzungen.
\UebQuellen{\FormulaRefAuto{FiniteGroupSquarePowerSemigroupReconstructionOneSidedNonempty}[thm]} \\
\addlinespace[.8ex]
\end{BandUebersicht}


\chapter[Band 43: Ringe mit Eins]{Band 43 auf einen Blick}
\noindent\textbf{Ringe mit Eins}\par\medskip
Ein Ring verbindet die abelschen Gruppen aus Band~41 mit Monoiden.
Die Multiplikation muss nicht kommutativ sein; \(0\ne1\) wird nicht vorausgesetzt.

\begin{BandUebersicht}
\textbf{Ringaxiome}

\UebFormel{\begin{aligned}&\AbelianGroupAlg R+0,\quad\MonoidAlg R\cdot1,\\&x(y+z)=xy+xz,\\&(x+y)z=xz+yz.\end{aligned}}
\UebQuellen{\FormulaRefAuto{RingDef}[def]}
&
\textbf{Null und Vorzeichen}

\UebFormel{\begin{aligned}&a0=0=0a,\quad-(-a)=a,\\&(-a)b=-(ab)=a(-b),\\&(-a)(-b)=ab.\end{aligned}}
\UebQuellen{\FormulaRefAuto{RingLeftZeroAbsorption}[thm];
\FormulaRefAuto{RingRightZeroAbsorption}[thm];
\FormulaRefAuto{RingDoubleNegative}[thm];
\FormulaRefAuto{RingNegativeTimesNegative}[thm]} \\
\addlinespace[.8ex]

\textbf{Ringhomomorphismen}

\UebFormel{\begin{aligned}&f(a+b)=f(a)+f(b),\\&f(ab)=f(a)f(b),\quad f(1_R)=1_S.\end{aligned}}
\UebQuellen{\FormulaRefAuto{RingHomDef}[def];
\FormulaRefAuto{RingIsoDef}[def]}
&
\textbf{Folgerungen}
Kompositionen sind wieder Ringhomomorphismen.
\UebFormel{f(0_R)=0_S,\qquad f(-a)=-f(a).}
\UebQuellen{\FormulaRefAuto{RingHomZeroPreservation}[thm];
\FormulaRefAuto{RingHomNegativePreservation}[thm];
\FormulaRefAuto{RingHomComposition}[thm]} \\
\addlinespace[.8ex]

\textbf{Ganzzahliger Ring}

\UebFormel{\RingAlg{\mathbb Z}{+}{\cdot}{\IntZero}{\IntOne}.}
\UebQuellen{\FormulaRefAuto{IntegerUnitalRing}[thm]}
&
\textbf{Kanonische Abbildung}
\(\eta_R\) ist die natürliche Halbringabbildung aus Band~39.
\UebFormel{\Phi_R(\IntClass a b)=\eta_R(a)\oplus(-\eta_R(b)).}
\UebQuellen{\FormulaRefAuto{IntegerRingMapDef}[def];
\FormulaRefAuto{IntegerRingMapClassValue}[thm]} \\
\addlinespace[.8ex]

\textbf{Erhaltung der Operationen}

\UebFormel{\begin{aligned}&\Phi_R(z+w)=\Phi_R(z)\oplus\Phi_R(w),\\&\Phi_R(zw)=\Phi_R(z)\odot\Phi_R(w).\end{aligned}}
\UebQuellen{\FormulaRefAuto{IntegerRingMapAddition}[thm];
\FormulaRefAuto{IntegerRingMapMultiplication}[thm];
\FormulaRefAuto{IntegerRingMapConstants}[thm]}
&
\textbf{Universelle Eigenschaft}
Für jeden Ring mit Eins ist \(\Phi_R\) der eindeutig bestimmte Ringhomomorphismus aus \(\mathbb Z\).
\UebFormel{\exists!f\,\RingHom f{\mathbb Z,+,\cdot,\IntZero,\IntOne}{R,\oplus,\odot,0_R,1_R}.}
\UebQuellen{\FormulaRefAuto{IntegerRingUniversalHomomorphism}[thm]} \\
\addlinespace[.8ex]
\end{BandUebersicht}

\chapter[Band 44: Kommutative Ringe mit Eins]{Band 44 auf einen Blick}
\noindent\textbf{Kommutative Ringe mit Eins}\par\medskip
Die Ringe mit Eins aus Band~43 werden durch Kommutativität der
Multiplikation spezialisiert.

\begin{BandUebersicht}
\textbf{Kommutativer Ring}

\UebFormel{\begin{aligned}&\CommutativeRingAlg R+\cdot01\ \leftrightarrow\\&\RingAlg R+\cdot01\land\\&\forall a,b\in R\,(a\cdot b=b\cdot a).\end{aligned}}
\UebQuellen{\FormulaRefAuto{CommutativeRingDef}[def]}
&
\textbf{Grundstruktur}
Alle allgemeinen Ringsätze aus Band~43 bleiben anwendbar.
\UebFormel{\CommutativeRingAlg R+\cdot01\vdash\RingAlg R+\cdot01.}
\UebQuellen{\FormulaRefAuto{CommutativeRingBaseRingAxiom}[ax]} \\
\addlinespace[.8ex]

\textbf{Morphismen}
Quelle und Ziel sind kommutative Ringe; \(f\) ist ein Ringhomomorphismus.
\UebFormel{\CommutativeRingHom f{R,+,\cdot,0,1}{S,\oplus,\odot,0_S,1_S}.}
\UebQuellen{\FormulaRefAuto{CommutativeRingHomDef}[def]}
&
\textbf{Isomorphismen}
Zur Homomorphie kommt die Bijektivität hinzu.
\UebFormel{\CommutativeRingIso f{R,+,\cdot,0,1}{S,\oplus,\odot,0_S,1_S}.}
\UebQuellen{\FormulaRefAuto{CommutativeRingIsoDef}[def];
\FormulaRefAuto{CommutativeRingIsoBijectiveAxiom}[ax]} \\
\addlinespace[.8ex]

\textbf{Multiplikation ganzer Zahlen}

\UebFormel{z,w\in\mathbb Z\vdash z\cdot w=w\cdot z.}
\UebQuellen{\FormulaRefAuto{IntegerMulCommutativeAxiom}[ax]}
&
\textbf{Grundbeispiel}

\UebFormel{\CommutativeRingAlg{\mathbb Z}{+}{\cdot}{\IntZero}{\IntOne}.}
\UebQuellen{\FormulaRefAuto{IntegerCommutativeRing}[thm]} \\
\addlinespace[.8ex]
\end{BandUebersicht}

\chapter[Band 45: Halbverbände und Verbände]{Band 45 auf einen Blick}
\noindent\textbf{Halbverbände und Verbände}\par\medskip
Dieser Band stellt Halbverbände eigenständig axiomatisch dar und
verbindet sie mit partiellen Ordnungen. Wir schreiben die induzierte
Supremumsordnung kurz als \(\le\). Endlichkeit und ein Nullelement
werden erst bei den entsprechend bezeichneten Resultaten vorausgesetzt.

\begin{BandUebersicht}
\textbf{Halbverbandsaxiome}
Eine Halbverbandsoperation \(\star:A\times A\to A\) erfüllt für
\(x,y,z\in A\):\UebFormel{\begin{aligned}
&(x\star y)\star z=x\star(y\star z),\\
&x\star y=y\star x,\qquad x\star x=x.\\
&x\le y\ \leftrightarrow\
 (x,y\in A\land x\star y=y).
\end{aligned}}
\UebQuellen{\FormulaRefAuto{Semilattice}[def];
\FormulaRefAuto{JoinOrder}[def]}
&
\textbf{Partielle Ordnung und Paarsupremum}
\UebFormel{\begin{aligned}
&\Semi{A,\star}\ \vdash\PartOrd{A,\le},\\
&\Semi{A,\star},\ x,y\in A\ \vdash\\
&\qquad x\star y=\sup_{A,\le}\{x,y\}.
\end{aligned}}
\UebQuellen{\FormulaRefAuto{JoinSemilatticeInducesOrder}[thm];
\FormulaRefAuto{JoinIsPairSupremum}[thm]} \\
\addlinespace[.8ex]

\textbf{Morphismen von Halbverbänden}
Für Halbverbände \((A,\star)\), \((B,\diamond)\) ist
\(f:A\to B\) ein Homomorphismus, wenn\UebFormel{\begin{aligned}
&x,y\in A\ \vdash\\
&\quad f(x\star y)=f(x)\diamond f(y).
\end{aligned}}Ein Isomorphismus ist ein bijektiver Homomorphismus.
\UebQuellen{\FormulaRefAuto{SemilatticeHomomorphism}[def];
\FormulaRefAuto{SemilatticeIsomorphism}[def]}
&
\textbf{Komposition}
Sind \(f:A\to B\) und \(g:B\to C\) Halbverbandshomomorphismen,
dann ist auch \(g\circ f:A\to C\) einer; insbesondere gilt\UebFormel{\begin{aligned}
&(g\circ f)(x\star y)\\
&\quad=(g\circ f)(x)\mathbin{\bullet}(g\circ f)(y).
\end{aligned}}Für Isomorphismen gilt dieselbe Abgeschlossenheit.
\UebQuellen{\FormulaRefAuto{SemilatticeHomComposition}[thm];
\FormulaRefAuto{SemilatticeIsoComposition}[thm]} \\
\addlinespace[.8ex]

\textbf{Vom Paarsupremum zur Algebra}
In einer partiell geordneten Menge \((A,\le)\) mit allen
Paarsuprema wird definiert:\UebFormel{\begin{aligned}
&x\star y=\sup_{A,\le}\{x,y\}
 \quad(x,y\in A).
\end{aligned}}
\UebQuellen{\FormulaRefAuto{PairJoinOperation}[def]}
&
\textbf{Rückgewinnung der Ordnung}
Diese Operation ist eine Halbverbandsoperation und gewinnt
die Ausgangsordnung zurück:\UebFormel{\begin{aligned}
&\Semi{A,\star},\\
&x,y\in A\ \vdash\
 (x\le y\leftrightarrow x\star y=y).
\end{aligned}}
\UebQuellen{\FormulaRefAuto{PairJoinCharacterization}[thm];
\FormulaRefAuto{PairJoinRecoversOrder}[thm]} \\
\addlinespace[.8ex]

\textbf{Infimale Lesart und Mengenbeispiele}
Infimal liest man eine Halbverbandsoperation durch\UebFormel{\begin{aligned}
&x\le_{\mathrm{inf}}y
 \ \leftrightarrow x\star y=x
 \quad(x,y\in A).
\end{aligned}}Auf \(\powerset(M)\) stehen dafür \(\cap\) und \(\cup\) zur Verfügung.
\UebQuellen{\FormulaRefAuto{MeetOrderEvaluation}[thm];
\FormulaRefAuto{PowerSetUnionSemilattice}[thm];
\FormulaRefAuto{PowerSetIntersectionSemilattice}[thm]}
&
\textbf{Dualität und Inklusion}
Bei der infimalen Lesart liefert \(x\star y\) das Paarinfimum.
Auf der Potenzmenge gilt\UebFormel{\begin{aligned}
&X\cup Y=Y\ \leftrightarrow X\subseteq Y,\\
&X\cap Y=X\ \leftrightarrow X\subseteq Y\\
&\hfill(X,Y\subseteq M).
\end{aligned}}
\UebQuellen{\FormulaRefAuto{MeetIsPairInfimum}[thm];
\FormulaRefAuto{PowerSetUnionOrderIsInclusion}[thm];
\FormulaRefAuto{PowerSetIntersectionOrderIsInclusion}[thm]} \\
\addlinespace[.8ex]

\textbf{Verband und Verbandshomomorphismus}
Ein Verband \((A,\MeetOp,\JoinOp)\) besitzt zwei
Halbverbandsoperationen und erfüllt\UebFormel{\begin{aligned}
&x\MeetOp(x\JoinOp y)=x,\\
&x\JoinOp(x\MeetOp y)=x\quad(x,y\in A).
\end{aligned}}Ein Verbandshomomorphismus erhält beide Operationen.
\UebQuellen{\FormulaRefAuto{AlgebraicLattice}[def];
\FormulaRefAuto{LatticeHomomorphism}[def];
\FormulaRefAuto{LatticeIsomorphism}[def]}
&
\textbf{Eine gemeinsame Ordnung}
\UebFormel{\begin{aligned}
&\LatticeAlg{A,\MeetOp,\JoinOp},\ x,y\in A\\
&\vdash (x\MeetOp y=x\leftrightarrow x\JoinOp y=y),\\
&\LatticeAlg{\powerset(M),\cap,\cup}.
\end{aligned}}Verbandshomomorphismen sind unter Komposition abgeschlossen.
\UebQuellen{\FormulaRefAuto{LatticeOrdersCoincide}[thm];
\FormulaRefAuto{PowerSetLattice}[thm];
\FormulaRefAuto{LatticeHomComposition}[thm]} \\
\addlinespace[.8ex]

\textbf{Null und Eins}
Ein Supremums-Halbverband mit Null besitzt \(0\in A\) mit
\(0\JoinOp x=x\). Ein beschränkter besitzt zusätzlich \(1\in A\) mit
\(x\JoinOp1=1\), jeweils für \(x\in A\).
\UebQuellen{\FormulaRefAuto{ZeroJoinSemilattice}[def];
\FormulaRefAuto{BoundedJoinSemilattice}[def]}
&
\textbf{Schranken und endliche Vollständigkeit}
Null ist das kleinste und Eins das größte Element.
Jeder endliche Supremums-Halbverband mit Null besitzt ein größtes Element:\UebFormel{\begin{aligned}
&\Finite A,\ \ZeroJoinSemi{A,\JoinOp,0}\ \vdash\\
&\quad\exists t\in A\ \forall x\in A\ (x\le t).
\end{aligned}}
\UebQuellen{\FormulaRefAuto{ZeroJoinLeast}[thm];
\FormulaRefAuto{BoundedJoinTopGreatest}[thm];
\FormulaRefAuto{FiniteZeroJoinSemilatticeHasGreatest}[thm];
\FormulaRefAuto{FiniteZeroJoinSemilatticeIsBounded}[thm]} \\
\addlinespace[.8ex]

\textbf{Gemeinsame untere Schranken}
Für \(x,y\in A\) setze\UebFormel{\begin{aligned}
&D(x,y)=\{d\in A\mid d\le x,\ d\le y\}.
\end{aligned}}In einem Supremums-Halbverband mit Null ist \(0\in D(x,y)\).
\UebQuellen{\FormulaRefAuto{ZeroJoinLeast}[thm]}
&
\textbf{Endliche Halbverbände mit Null haben Paarinfima}
\UebFormel{\begin{aligned}
&\Finite A,\ \ZeroJoinSemi{A,\JoinOp,0},\\
&x,y\in A\ \vdash\\
&\exists m\in D(x,y)\ \forall d\in D(x,y)\
 (d\le m).
\end{aligned}}Dieses größte untere Element ist das Paarinfimum.
\UebQuellen{\FormulaRefAuto{FiniteZeroJoinSemilatticePairInfima}[thm]} \\
\addlinespace[.8ex]

\textbf{Irreduzible und Trenner}
In einem Halbverband mit Null heißt \(j\in A\setminus\{0\}\)
\(\JoinOp\)-irreduzibel, wenn\UebFormel{\begin{aligned}
&x,y\in A,\quad x\JoinOp y=j\\
&\qquad\vdash x=j\lor y=j.\\
&\operatorname{Sep}(x,y)
 =\{d\in A\mid d\le x,\ d\not\le y\}.
\end{aligned}}
\UebQuellen{\FormulaRefAuto{JoinIrreducible}[def];
\FormulaRefAuto{JoinSeparatorSet}[def]}
&
\textbf{Irreduzible trennen}
Minimale Trenner sind irreduzibel. Ist außerdem \(A\) endlich,
so gilt für \(x,y\in A\):\UebFormel{\begin{aligned}
&x\not\le y\ \vdash\\
&\exists j\in\JoinIrrSet{A,\JoinOp,0}\\
&\qquad(j\le x\land j\not\le y).
\end{aligned}}
\UebQuellen{\FormulaRefAuto{MinimalSeparatorJoinIrreducible}[thm];
\FormulaRefAuto{JoinIrreduciblesSeparate}[thm]} \\
\addlinespace[.8ex]

\textbf{Echte untere Elemente}
In der im Band verwendeten Notation für einen Halbverband \(L\)
mit Nullelement \(\varnothing\):\UebFormel{\begin{aligned}
&D_L(P)=\\
&\quad\{Q\in L\mid Q\le P,\ Q\ne P\}.
\end{aligned}}
\UebQuellen{\FormulaRefAuto{JoinIrreducibleProperLowerSet}[def]}
&
\textbf{Größtes echtes unteres Element}
\UebFormel{\begin{aligned}
&\Finite L,\ \ZeroJoinSemi{L,\JoinOp,\varnothing},\\
&\JoinIrr{L,\JoinOp,\varnothing}P\ \vdash\\
&\exists P_\ast\in D_L(P)\
 \forall Q\in D_L(P)\ (Q\le P_\ast).
\end{aligned}}Diese Eigenschaft trägt später den Beweis über private Punkte.
\UebQuellen{\FormulaRefAuto{JoinIrreducibleLargestProperLower}[thm]} \\
\addlinespace[.8ex]
\end{BandUebersicht}


\chapter[Band 46: Frankls Vermutung]{Band 46 auf einen Blick}
\noindent\textbf{Frankls Vermutung}\par\medskip
Die allgemeine Frankl-Vermutung wird in diesem Band
\emph{nicht bewiesen}. Bewiesen werden Spezialfälle, Reduktionen und die
Äquivalenz von Mengen- und Halbverbandsform. Vorausgesetzt werden
Mengenfamilien, Quotienten, endliche Kardinalitäten und Band~45.
Die Häufigkeitsaussagen werden wie im Original durch Injektionen formuliert.

\begin{BandUebersicht}
\textbf{Frankl-Familie}
\UebFormel{\begin{aligned}
&\FranklFam M\ \leftrightarrow\\
&M\ne\varnothing\land\bigcup M\ne\varnothing\\
&\quad\land\Finite M\\
&\quad\land\UnionClosedFam(M).\\
&\UnionClosedFam(M)\ \leftrightarrow\\
&\quad\forall X,Y\in M\\
&\qquad(X\cup Y\in M).\\
&M_a=\{X\in M\mid a\in X\},\\
&M_{\bar a}=\{X\in M\mid a\notin X\}.
\end{aligned}}
\UebQuellen{\FormulaRefAuto{Frankl-Familie}[def];
\FormulaRefAuto{Vereinigungsabgeschlossene Familie}[def];
\FormulaRefAuto{Enthaltende Teilfamilie}[def];
\FormulaRefAuto{Vermeidende Teilfamilie}[def]}
&
\textbf{Zerlegung und Endlichkeit der Teilfamilien}
\UebFormel{\begin{aligned}
&M\setminus M_{\bar a}=M_a,\\
&\Finite M\ \vdash\Finite{M_a}\land\Finite{M_{\bar a}}.
\end{aligned}}
\UebQuellen{\FormulaRefAuto{FamContainingAsAvoidingRelativeComplement}[thm];
\FormulaRefAuto{FiniteContainingSubfamily}[thm];
\FormulaRefAuto{FiniteAvoidingSubfamily}[thm]} \\
\addlinespace[.8ex]

\textbf{Halbhäufigkeit und Frankl-Eigenschaft}
Mit den eben definierten Teilfamilien:\UebFormel{\begin{aligned}
&\HalfOcc{a}{M}\ \leftrightarrow\\
&a\in\bigcup M\land
 \exists f:M_{\bar a}\inj M_a,\\
&\Frankl{M}\ \leftrightarrow\exists a\ \HalfOcc{a}{M}.
\end{aligned}}Bei endlichem \(M\) heißt dies: \(a\) liegt in mindestens der
Hälfte der Familienglieder.
\UebQuellen{\FormulaRefAuto{Halbhäufiges Element}[def];
\FormulaRefAuto{Frankl-Eigenschaft}[def]}
&
\textbf{Bewiesener Spezialfall: Zweiermenge}
\UebFormel{\begin{aligned}
&\FranklFam M,\quad a\ne b,\quad \{a,b\}\in M\\
&\qquad\vdash\Frankl{M}.
\end{aligned}}Die bloße Vereinigungsgeschlossenheit einer beliebigen
Frankl-Familie wird dadurch nicht erledigt.
\UebQuellen{\FormulaRefAuto{FranklFamPairMemberFrankl}[thm]} \\
\addlinespace[.8ex]

\textbf{Ununterscheidbarkeitsquotient}
Für \(a,b\in\bigcup M\) setzt man\UebFormel{\begin{aligned}
&a\sim_M b\ \leftrightarrow\\
&\qquad\forall X\in M\ (a\in X\leftrightarrow b\in X).
\end{aligned}}Ersetzt man jedes \(a\in X\) durch seine Klasse \([a]\), entsteht
die Quotientenfamilie \(\OccQuotFam M\).
\UebQuellen{\FormulaRefAuto{SameOccRelDef}[def];
\FormulaRefAuto{OccQuotMapChar}[thm]}
&
\textbf{Reduktion auf endliche Glieder}
\UebFormel{\begin{aligned}
&\FranklFam M\ \vdash\\
&\qquad\FranklFam{\OccQuotFam M},\\
&\Finite M,\ Y\in\OccQuotFam M\\
&\qquad\vdash\Finite Y,\\
&\Frankl{\OccQuotFam M}\ \vdash\Frankl{M}.
\end{aligned}}Es genügt daher, die Vermutung für Frankl-Familien mit
ausschließlich endlichen Gliedern zu beweisen.
\UebQuellen{\FormulaRefAuto{OccQuotFamFranklFam}[thm];
\FormulaRefAuto{OccQuotMembersFinite}[thm];
\FormulaRefAuto{OccQuotFranklTransport}[thm];
\FormulaRefAuto{FranklFiniteMembersReduction}[thm]} \\
\addlinespace[.8ex]

\textbf{Adjunktion der leeren Menge}
\UebFormel{\begin{aligned}
&M^\circ=M\cup\{\varnothing\}.
\end{aligned}}
\UebQuellen{\FormulaRefAuto{FranklEmptyAdjunction}[def]}
&
\textbf{Normalisierung und Rücktransport}
Für \(\FranklFam M\) gilt:
\UebFormel{\begin{aligned}
&\FranklFam M\ \vdash\FranklFam{M^\circ},\\
&\HalfOcc{a}{M^\circ}\ \vdash\HalfOcc{a}{M},\\
&\Frankl{M^\circ}\ \vdash\Frankl{M}.
\end{aligned}}Die globale Vermutung ist äquivalent zu ihrer Einschränkung
auf Familien, welche die leere Menge enthalten.
\UebQuellen{\FormulaRefAuto{FranklEmptyAdjunctionPreservesFamily}[thm];
\FormulaRefAuto{FranklEmptyAdjunctionHalfOccurrenceBack}[thm];
\FormulaRefAuto{FranklEmptyAdjunctionFranklBack}[thm];
\FormulaRefAuto{FranklEmptyFamilyGlobalEquivalence}[thm]} \\
\addlinespace[.8ex]

\textbf{Halbverbandsform}
Für \(\ZeroJoinSemi{A,\JoinOp,0}\) und dessen Ordnung \(\le\)
sei \(\uparrow j=\{x\in A\mid j\le x\}\). Dann bedeutet\UebFormel{\begin{aligned}
&\FranklSemi{A,\JoinOp,0}\ \leftrightarrow\\
&\exists j\in\JoinIrrSet{A,\JoinOp,0}\
 \exists f:\uparrow j\inj(A\setminus\uparrow j).
\end{aligned}}Im endlichen Fall soll also der Hauptfilter eines
Irreduziblen höchstens die halbe Größe von \(A\) haben.
\UebQuellen{\FormulaRefAuto{FranklSemilatticeProperty}[def]}
&
\textbf{Bewiesene Halbverbandsfälle}
Ein \(\JoinOp\)-irreduzibles größtes Element \(t\) genügt:\UebFormel{\begin{aligned}
&\JoinIrr{A,\JoinOp,0}t,\quad
 \forall x\in A\ (x\le t)\\
&\qquad\vdash\FranklSemi{A,\JoinOp,0}.
\end{aligned}}Auch jede endliche, total geordnete solche Struktur mit
größtem Element \(t\ne0\) erfüllt die Eigenschaft.
\UebQuellen{\FormulaRefAuto{JoinIrreducibleTopFrankl}[thm];
\FormulaRefAuto{FranklJoinChains}[thm]} \\
\addlinespace[.8ex]

\textbf{Kanonische Profilfamilie}
Für einen endlichen Halbverband \((A,\JoinOp,0)\) mit Null setze
\(J=\JoinIrrSet{A,\JoinOp,0}\) und\UebFormel{\begin{aligned}
&J(x)=\{j\in J\mid j\le x\},\\
&\mathcal U=\{J\setminus J(x)\mid x\in A\},\\
&\Phi(x)=J\setminus J(x).
\end{aligned}}
\UebQuellen{\FormulaRefAuto{IrreduciblesBelow}[def];
\FormulaRefAuto{CanonicalProfileFamilies}[def];
\FormulaRefAuto{CanonicalProfileMapDefinition}[def]}
&
\textbf{Bijektion und Häufigkeitstransport}
\UebFormel{\begin{aligned}
&\Phi:A\bij\mathcal U,\qquad
 \UnionClosedFam(\mathcal U).
\end{aligned}}Ist \(A\ne\{0\}\), dann ist \(\mathcal U\) eine Frankl-Familie,
und\UebFormel{\begin{aligned}
&\HalfOcc{j}{\mathcal U}\ \vdash\
 \FranklSemi{A,\JoinOp,0}.
\end{aligned}}
\UebQuellen{\FormulaRefAuto{CanonicalProfileBijection}[thm];
\FormulaRefAuto{CanonicalProfileFamilyUnionClosed}[thm];
\FormulaRefAuto{CanonicalUnionFamilyIsFranklFamily}[thm];
\FormulaRefAuto{CanonicalHalfOccurrenceTransport}[thm]} \\
\addlinespace[.8ex]

\textbf{Komplementduale Konstruktion}
Für eine Frankl-Familie \(M\):\UebFormel{\begin{aligned}
&U=\bigcup M,\\
&L=\{U\setminus X\mid X\in M^\circ\},\\
&X\JoinOp Y=\bigcap\mathcal Z_{X,Y},\\
&\mathcal Z_{X,Y}=\{Z\in L\mid X\cup Y\subseteq Z\}.
\end{aligned}}
\UebQuellen{\FormulaRefAuto{FranklComplementConstruction}[def]}
&
\textbf{Ein endlicher Halbverband mit Null}
\UebFormel{\begin{aligned}
&\Finite L,\qquad\InterClosedFam L,\\
&\ZeroJoinSemi{L,\JoinOp,\varnothing},\\
&X,Y\in L\ \vdash\
 (X\le Y\leftrightarrow X\subseteq Y).
\end{aligned}}
\UebQuellen{\FormulaRefAuto{FranklComplementFamilyFinite}[thm];
\FormulaRefAuto{FranklComplementFamilyIntersectionClosed}[thm];
\FormulaRefAuto{FranklComplementFiniteLattice}[thm];
\FormulaRefAuto{FranklComplementOrderRecovered}[thm]} \\
\addlinespace[.8ex]

\textbf{Private Punkte}
In der vorigen Komplementkonstruktion besitzt jedes
\(\JoinOp\)-irreduzible \(P\in L\) einen privaten Punkt \(a\in P\):\UebFormel{\begin{aligned}
&Q\in L\ \vdash\
 (a\in Q\leftrightarrow P\le Q).
\end{aligned}}Der Hauptfilter \(H=\uparrow P\) ist damit genau die
Teilfamilie der \(a\) enthaltenden Mengen.
\UebQuellen{\FormulaRefAuto{JoinIrreduciblePrivatePoint}[thm]}
&
\textbf{Rücktransport zur Mengenform}
Für einen solchen privaten Punkt gilt:\UebFormel{\begin{aligned}
&\exists f:H\inj(L\setminus H)\\
&\qquad\vdash\HalfOcc{a}{M}.
\end{aligned}}Ein Irreduzibles mit höchstens halbgroßem Hauptfilter liefert
somit ein halbhäufiges Element der ursprünglichen Familie.
\UebQuellen{\FormulaRefAuto{FranklComplementPrivatePointTransport}[thm]} \\
\addlinespace[.8ex]

\textbf{Die beiden globalen Aussagen}
Die Mengenform lautet\UebFormel{\begin{aligned}
&\FranklSetStatement\ \leftrightarrow\\
&\quad\forall M\\
&\qquad(\FranklFam M\to\Frankl{M}).
\end{aligned}}Die Halbverbandsform \(\FranklJoinStatement\) fordert
\(\FranklSemi{A,\JoinOp,0}\) für jeden endlichen Halbverband mit Null
und mindestens einem von Null verschiedenen Element.
\UebQuellen{\FormulaRefAuto{FranklGlobalSetStatement}[def];
\FormulaRefAuto{FranklGlobalJoinStatement}[def]}
&
\textbf{Bewiesene Äquivalenz, offene Vermutung}
\UebFormel{\begin{aligned}
&\FranklSetStatement\leftrightarrow\FranklJoinStatement.
\end{aligned}}Der Satz beweist die beiden Transporte zwischen den globalen
Aussagen. Er beweist keine dieser Aussagen für sich.
\UebQuellen{\FormulaRefAuto{FranklSetImpliesSemilattice}[thm];
\FormulaRefAuto{FranklSemilatticeImpliesSet}[thm];
\FormulaRefAuto{FranklSetSemilatticeEquivalence}[thm]} \\
\addlinespace[.8ex]
\end{BandUebersicht}


\chapter[Band 47: Metrische Räume und Vollständigkeit]{Band 47 auf einen Blick}
\noindent\textbf{Metrische Räume und Vollständigkeit}\par\medskip
\noindent\emph{Lesart.}
Soweit nicht anders angegeben, ist \((X,d)\) ein metrischer Raum,
sind \(A,K,U\subseteq X\), \(x,y,z\in X\), \(a\colon\mathbb N\to X\),
und Radien sowie \(\varepsilon,\delta\) sind positive reelle Zahlen.
Bei Abbildungen \(f\colon X\to Y\) sind \(d_X,d_Y\) die jeweiligen Metriken.
\(\varphi\colon\mathbb N\to\mathbb N\) bezeichnet stets eine strikt
wachsende Teilfolgenabbildung.

\begin{BandUebersicht}
\textbf{Metrik und ihre Axiome}
\(\operatorname{Metrik}_X(d)\) umfasst
\UebFormel{\begin{aligned}
&d\colon X\times X\to\mathbb R,\qquad0\leq d(x,y),\\
&d(x,y)=0\leftrightarrow x=y,\\
&d(x,y)=d(y,x),\\
&d(x,z)\leq d(x,y)+d(y,z).
\end{aligned}}
\UebQuellen{\FormulaRefAuto{MetricSpaceDef}[def];
\FormulaRefAuto{MetricFunctionAxiom}[ax];
\FormulaRefAuto{MetricPositivityAxiom}[ax];
\FormulaRefAuto{MetricDefinitenessAxiom}[ax];
\FormulaRefAuto{MetricSymmetryAxiom}[ax];
\FormulaRefAuto{MetricTriangleAxiom}[ax]}
&
\textit{Umgekehrte Dreiecksungleichung}
\UebFormel{|d(x,z)-d(y,z)|\leq d(x,y).}
Ein Abstand von einem festen Punkt kann sich höchstens um den
Abstand zwischen den beiden bewegten Punkten ändern.
\UebQuellen{\FormulaRefAuto{MetricReverseTriangleInequality}[thm]}
\\ \addlinespace[.8ex]

\textbf{Reelle und diskrete Metrik}
\UebFormel{\begin{aligned}
d_{\mathbb R}(x,y)&\coloneqq|x-y|,\\
\delta_X(x,y)&\coloneqq
\begin{cases}0,&x=y,\\1,&x\neq y.\end{cases}
\end{aligned}}
Die erste Formel gilt für \(x,y\in\mathbb R\), die zweite auf jeder Menge \(X\).
\UebQuellen{\FormulaRefAuto{RealMetricDistanceFunctionDef}[def];
\FormulaRefAuto{DiscreteMetricDef}[def]}
&
\textit{Modelle und diskrete Konvergenz}
\UebFormel{\begin{aligned}
&\operatorname{Metrik}_{\mathbb R}(d_{\mathbb R}),\\
&\operatorname{Metrik}_X(\delta_X),\\
&a_n\MetricTo{\delta_X}x
\leftrightarrow\exists N\;\forall n\geq N\;(a_n=x).
\end{aligned}}
\UebQuellen{\FormulaRefAuto{RealDistanceIsMetric}[thm];
\FormulaRefAuto{DiscreteMetricIsMetric}[thm];
\FormulaRefAuto{DiscreteMetricConvergenceEventuallyConstant}[thm]}
\\ \addlinespace[.8ex]

\textbf{Teilraummetrik}
Für \(Y\subseteq X\):
\UebFormel{d_Y\coloneqq d\restriction_{Y\times Y}.}
\UebQuellen{\FormulaRefAuto{MetricSubspaceRestrictionDef}[def]}
&
\textit{Einschränkung erhält die metrischen Begriffe}
\UebFormel{\operatorname{Metrik}_Y(d_Y).}
Für \(a\colon\mathbb N\to Y\) und \(y\in Y\):
\UebFormel{\begin{aligned}
&a_n\MetricTo{d_Y}y\leftrightarrow a_n\MetricTo d y,\\
&\operatorname{Cauchy}_{d_Y}(a)
\leftrightarrow\operatorname{Cauchy}_d(a).
\end{aligned}}
\UebQuellen{\FormulaRefAuto{MetricSubspaceRestrictionIsMetric}[thm];
\FormulaRefAuto{MetricSubspaceConvergenceAgree}[thm];
\FormulaRefAuto{MetricSubspaceCauchyAgree}[thm]}
\\ \addlinespace[.8ex]

\textbf{Offene Kugeln und offene Mengen}
\UebFormel{\begin{aligned}
&\MetricBall d x r\coloneqq\{y\in X\mid d(y,x)<r\},\\
&\operatorname{Offen}_d(U)\coloneqq\\
&\quad\forall x\in U\;\exists r>0\;
(\MetricBall d x r\subseteq U).
\end{aligned}}
\UebQuellen{\FormulaRefAuto{MetricOpenBallDef}[def];
\FormulaRefAuto{MetricOpenClosedSetDef}[def]}
&
\textit{Kugelverkleinerung und Offenheit}
\UebFormel{\begin{aligned}
&y\in\MetricBall d x r\vdash\exists s>0\\
&\qquad\MetricBall d y s\subseteq\MetricBall d x r,\\
&\operatorname{Offen}_d(\MetricBall d x r).
\end{aligned}}
Für \(d_{\mathbb R}\) sind diese Kugeln die reellen
Epsilon-Umgebungen.
\UebQuellen{\FormulaRefAuto{MetricBallShrinking}[thm];
\FormulaRefAuto{MetricOpenBallsAreOpen}[thm];
\FormulaRefAuto{RealMetricBallIsEpsilonNeighborhood}[thm]}
\\ \addlinespace[.8ex]

\textbf{Metrische Konvergenz}
\UebFormel{\begin{aligned}
a_n\MetricTo d x\coloneqq{}&
\forall\varepsilon>0\;\exists N\in\mathbb N\\
&\forall n\geq N\;d(a_n,x)<\varepsilon.
\end{aligned}}
\UebQuellen{\FormulaRefAuto{MetricSequenceConvergenceDef}[def]}
&
\textit{Eindeutigkeit und Beständigkeit des Grenzwerts}
\UebFormel{\begin{aligned}
&a_n\MetricTo d x,\ a_n\MetricTo d y\vdash x=y,\\
&a_n\MetricTo d x\vdash a_{\varphi(n)}\MetricTo d x.
\end{aligned}}
Endliche Änderungen und Restfolgen erhalten ebenfalls den Grenzwert;
für \(d_{\mathbb R}\) stimmt die Definition mit der reellen Konvergenz überein.
\UebQuellen{\FormulaRefAuto{MetricSequenceLimitUnique}[thm];
\FormulaRefAuto{MetricSubsequencePreservesLimit}[thm];
\FormulaRefAuto{MetricFiniteModificationPreservesLimit}[thm];
\FormulaRefAuto{MetricSequenceTailSameLimit}[thm];
\FormulaRefAuto{RealAndMetricSequenceConvergenceAgree}[thm]}
\\ \addlinespace[.8ex]

\textbf{Beschränkte Mengen und Folgen}
\UebFormel{\begin{aligned}
&\operatorname{Beschr}_d(A)\coloneqq\\
&\quad(A=\varnothing\lor
\exists c\in X\;\exists r>0\\
&\qquad A\subseteq\MetricBall d c r),\\
&\operatorname{Beschr}_d(a)\coloneqq
\operatorname{Beschr}_d(a[\mathbb N]).
\end{aligned}}
\UebQuellen{\FormulaRefAuto{MetricBoundedSetSequenceDef}[def]}
&
\textit{Notwendige, aber keine hinreichende Bedingung}
\UebFormel{a_n\MetricTo d x\vdash\operatorname{Beschr}_d(a).}
Die diskrete Menge \((\mathbb N,\delta_{\mathbb N})\) ist beschränkt,
aber nicht folgenkompakt.
\UebQuellen{\FormulaRefAuto{ConvergentMetricSequenceBounded}[thm];
\FormulaRefAuto{BoundedDiscreteNaturalNumbersNotSequentiallyCompact}[thm]}
\\ \addlinespace[.8ex]

\textbf{Metrische Cauchy-Folge}
\UebFormel{\begin{aligned}
&\operatorname{Cauchy}_d(a)\coloneqq\\
&\quad\forall\varepsilon>0\;\exists N\in\mathbb N\\
&\qquad\forall m,n\geq N\;d(a_m,a_n)<\varepsilon.
\end{aligned}}
\UebQuellen{\FormulaRefAuto{MetricCauchySequenceDef}[def]}
&
\textit{Cauchy-Folgen und konvergente Teilfolgen}
\UebFormel{\begin{aligned}
&a_n\MetricTo d x\vdash\operatorname{Cauchy}_d(a),\\
&\operatorname{Cauchy}_d(a)\vdash\operatorname{Beschr}_d(a),\\
&\operatorname{Cauchy}_d(a),\
a_{\varphi(n)}\MetricTo d x\\
&\qquad\vdash a_n\MetricTo d x.
\end{aligned}}
\UebQuellen{\FormulaRefAuto{MetricConvergentSequenceIsCauchy}[thm];
\FormulaRefAuto{MetricCauchySequenceBounded}[thm];
\FormulaRefAuto{MetricCauchyWithConvergentSubsequence}[thm]}
\\ \addlinespace[.8ex]

\textbf{Vollständigkeit}
\UebFormel{\begin{aligned}
&\operatorname{Vollst}(X,d)\coloneqq\\
&\quad\forall a\colon\mathbb N\to X\\
&\qquad(\operatorname{Cauchy}_d(a)\to\\
&\qquad\exists x\in X\;(a_n\MetricTo d x)).
\end{aligned}}
\UebQuellen{\FormulaRefAuto{CompleteMetricSpaceDef}[def]}
&
\textit{Vollständige und unvollständige Modelle}
\UebFormel{\begin{aligned}
&\operatorname{Vollst}(\mathbb R,d_{\mathbb R}),\\
&\operatorname{Vollst}(X,\delta_X),\\
&\neg\operatorname{Vollst}
(\mathbb R\setminus\{0\},(d_{\mathbb R})_{\mathbb R\setminus\{0\}}).
\end{aligned}}
Die letzte Aussage wird durch eine Cauchy-Folge mit fehlendem
Grenzwert \(0\) im Teilraum belegt.
\UebQuellen{\FormulaRefAuto{RealMetricSpaceComplete}[thm];
\FormulaRefAuto{DiscreteMetricSpaceComplete}[thm];
\FormulaRefAuto{PuncturedRealLineIncomplete}[thm]}
\\ \addlinespace[.8ex]

\textbf{Abgeschlossenheit}
\UebFormel{\begin{aligned}
\operatorname{Abgeschlossen}_d(A)
\coloneqq\operatorname{Offen}_d(X\setminus A).
\end{aligned}}
\UebQuellen{\FormulaRefAuto{MetricOpenClosedSetDef}[def]}
&
\textit{Grenzwerte und vollständige Teilräume}
Ist \(A\) abgeschlossen, so gilt für \(a\colon\mathbb N\to A\):
\UebFormel{a_n\MetricTo d x\vdash x\in A.}
Ist außerdem \(X\) vollständig, dann
\UebFormel{\operatorname{Vollst}(A,d_A).}
\UebQuellen{\FormulaRefAuto{MetricClosedSetContainsSequenceLimits}[thm];
\FormulaRefAuto{ClosedSubspaceOfCompleteMetricSpace}[thm]}
\\ \addlinespace[.8ex]

\textbf{Stetigkeit}
\UebFormel{\begin{aligned}
&\operatorname{Stetig}_{d_X,d_Y}(f,x)\coloneqq\\
&\quad\forall\varepsilon>0\;\exists\delta>0\;\forall u\in X:\\
&\qquad d_X(u,x)<\delta\\
&\qquad\to d_Y(f(u),f(x))<\varepsilon.
\end{aligned}}
Stetigkeit von \(f\) bedeutet Stetigkeit in jedem \(x\in X\).
\UebQuellen{\FormulaRefAuto{MetricContinuityDef}[def]}
&
\textit{Folgenkriterium und Komposition}
\UebFormel{\begin{aligned}
&\operatorname{Stetig}_{d_X,d_Y}(f,x)\\
&\ \leftrightarrow
\forall a\colon\mathbb N\to X:\\
&\qquad(a_n\MetricTo{d_X}x
\to f(a_n)\MetricTo{d_Y}f(x)).
\end{aligned}}
Für stetige \(f\colon X\to Y\) und \(g\colon Y\to Z\)
ist auch \(g\circ f\) stetig.
\UebQuellen{\FormulaRefAuto{MetricSequentialContinuityCriterion}[thm];
\FormulaRefAuto{MetricContinuousComposition}[thm]}
\\ \addlinespace[.8ex]

\textbf{Isometrie}
\UebFormel{\begin{aligned}
&\operatorname{Isometrie}_{d_X,d_Y}(f)\coloneqq\\
&\quad\forall x,u\in X\\
&\qquad d_Y(f(x),f(u))=d_X(x,u).
\end{aligned}}
\UebQuellen{\FormulaRefAuto{MetricIsometryDef}[def]}
&
\textit{Erhalt der metrischen Struktur}
Unter der Isometrieannahme:
\UebFormel{\begin{aligned}
&\operatorname{Stetig}_{d_X,d_Y}(f),\\
&\operatorname{Cauchy}_{d_X}(a)\vdash
\operatorname{Cauchy}_{d_Y}(f\circ a).
\end{aligned}}
\UebQuellen{\FormulaRefAuto{MetricIsometryPreservesStructure}[thm]}
\\ \addlinespace[.8ex]

\textbf{Folgenkompaktheit}
\UebFormel{\begin{aligned}
&\operatorname{Folgenkompakt}_d(K)\coloneqq\\
&\quad\forall a\colon\mathbb N\to K\;\exists x\in K\\
&\quad\exists\varphi\text{ strikt wachsend}:\\
&\qquad a_{\varphi(n)}\MetricTo d x.
\end{aligned}}
Der Grenzwert muss zur betrachteten Teilmenge \(K\) gehören.
\UebQuellen{\FormulaRefAuto{MetricSequentialCompactnessDef}[def]}
&
\textit{Strukturelle Folgerungen}
Jede folgenkompakte Teilmenge ist beschränkt und abgeschlossen;
als metrischer Teilraum ist sie vollständig.
\UebFormel{\operatorname{Folgenkompakt}_d(K)
\vdash\operatorname{Vollst}(K,d_K).}
Endliche Teilmengen sind folgenkompakt, und stetige Bilder
folgenkompakter Mengen sind folgenkompakt.
\UebQuellen{\FormulaRefAuto{MetricSequentiallyCompactSetBounded}[thm];
\FormulaRefAuto{MetricSequentiallyCompactSetClosed}[thm];
\FormulaRefAuto{MetricSequentiallyCompactSetComplete}[thm];
\FormulaRefAuto{FiniteMetricSetSequentiallyCompact}[thm];
\FormulaRefAuto{MetricContinuousImageSequentiallyCompact}[thm]}
\\ \addlinespace[.8ex]

\textbf{Reelle Spezialisierung}
Die metrischen Begriffe werden mit
\UebFormel{X=\mathbb R,\qquad d_{\mathbb R}(x,y)=|x-y|}
verwendet. Eine reelle Folge ist beschränkt, wenn ihre Wertemenge
in einer Kugel liegt.
\UebQuellen{\FormulaRefAuto{RealMetricDistanceFunctionDef}[def];
\FormulaRefAuto{MetricBoundedSetSequenceDef}[def]}
&
\textit{Bolzano--Weierstraß und kompakte Intervalle}
Jede beschränkte reelle Folge besitzt eine konvergente Teilfolge:
\UebFormel{\begin{aligned}
&a\colon\mathbb N\to\mathbb R,\ \operatorname{Beschr}(a)\\
&\qquad\vdash\exists L\in\mathbb R\;\exists\varphi
\text{ strikt wachsend}:\\
&\hspace{7em}a_{\varphi(n)}\to L.
\end{aligned}}
Für reelle \(u\leq v\) ist \([u,v]_{\mathbb R}\) folgenkompakt.
\UebQuellen{\FormulaRefAuto{RealSequenceMonotoneSubsequence}[thm];
\FormulaRefAuto{RealSequenceBolzanoWeierstrass}[thm];
\FormulaRefAuto{RealClosedBoundedIntervalSequentiallyCompact}[thm]}
\\
\end{BandUebersicht}

\noindent\emph{Anschluss.}
Aus einer Norm entsteht der Abstand \(d(x,y)=\lVert x-y\rVert\).
Der abschließende Ausblick ordnet dies in die spätere Theorie der Vektorräume ein.
Die hier entwickelten Begriffe benötigen auf \(X\) selbst keine
Addition oder Skalarmultiplikation.\par

\chapter[Band 48: Axiomatische Mengenlehre II]{Band 48 auf einen Blick}
\noindent\textbf{Von den Mengenaxiomen zu Unabhängigkeitsbeweisen}\par\medskip
\noindent\emph{Lesart.}
Dieser Band setzt die Mengenlehre aus Band~3 und die bereits
entwickelten Begriffe der Funktionen, Wohlordnungen und reellen
Zahlen voraus. Er wiederholt diese Grundlagen nicht, sondern
ergänzt die Metatheorie und die beiden Modellkonstruktionen für
die Unabhängigkeitsfragen der Kapitel~4 und~5.

\section*{Kerninhalt}
\begin{itemize}
 \item Syntax, Erfüllung, Korrektheit und ein Beweis des
       Vollständigkeitssatzes durch Henkin-Erweiterungen.
 \item Ordinale, transfinite und fundierte Rekursion,
       Mostowski-Kollaps, Kardinalarithmetik und Skolemhüllen.
 \item Konstruktible Hierarchie, Reflexion, Kondensation und
       relative Konsistenz von ZFC mit GCH.
 \item Mengenforcing über beliebigen abzählbaren Modellen,
       Wahrheitslemma, Erhaltung von ZFC und ccc-Kardinalerhaltung.
 \item Cohen-Forcing und Zählung schöner Namen bis zum Modell
       mit \(2^{\aleph_0}=2^{\aleph_1}=\aleph_2\).
\end{itemize}

\section*{Erreichte Ergebnisse}
Die Vorwärtsrichtung
\(A\approx B\Rightarrow\mathcal P(A)\approx\mathcal P(B)\),
Cantors strikte Ungleichung und die endliche Rückrichtung
werden in Beweistabellen hergeleitet.

Unter der ausdrücklich genannten Voraussetzung
\(\operatorname{Con}(\mathsf{ZFC})\) ist die universelle Aussage
\[
 \forall A,B\
 (\mathcal P(A)\approx\mathcal P(B)\to A\approx B)
\]
weder aus ZFC beweisbar noch in ZFC widerlegbar; siehe
\FormulaRefAuto{B48PCIndependence}[thm].
Dasselbe gilt für die Kontinuumshypothese
\(2^{\aleph_0}=\aleph_1\); siehe
\FormulaRefAuto{B48CHIndependence}[thm].
Die Gleichheit \(|\mathbb R|=2^{\aleph_0}\) und die
übliche Zwischenmächtigkeitsformulierung werden ebenfalls bewiesen.

\section*{Didaktische Stellung}
Der Aufbau folgt dem Schema Definition, Hilfssatz,
Beweistabelle und Anwendung. Der Unvollständigkeitssatz ist
keine Voraussetzung. Weder die Konsistenz von ZFC noch
die Existenz eines transitiven ZFC-Modells wird
voraussetzungslos behauptet. Die Tabellen sind mathematische
Beweise, keine maschinell geprüfte Formalisierung.


\end{document}
